% DER_XVIII-d: Grothendieck's text with the commentary set in, then the overview, inside a group that applies the chunk's own definitions
\dgrpart{1}{XVIII-d}{Part 4 of 5, manuscript pages 316--429}{XVIII-d\enspace\textperiodcentered\enspace Chapter XVIII, part 4 of 5\enspace\textperiodcentered\enspace manuscript pages 316--429}{Chapter XVIII \textperiodcentered\ \textup{XVIII-d}}
\dgrchunkbegin
% DER_XVIII-d: the definitions of the chunk's preamble, applied inside its group
\let\printedpage\relax\newcommand{\printedpage}[1]{\marginpar{\footnotesize p.~#1}}
\let\msp\relax\newcommand{\msp}[1]{\marginpar{\footnotesize [#1]}}
\let\circled\relax\newcommand{\circled}[1]{\tikz[baseline=(c.base),commutative diagrams/math mode=false]\node[draw,circle,inner sep=1pt](c){#1};}
\let\ovalbox\relax\newcommand{\ovalbox}[1]{\tikz[baseline=(c.base)]\node[draw,rounded corners=6pt,inner sep=2pt](c){#1};}
\let\Ob\relax\newcommand{\Ob}{\operatorname{Ob}}
\let\card\relax\newcommand{\card}{\operatorname{card}}
\let\Ind\relax\newcommand{\Ind}{\operatorname{Ind}}
\let\Ord\relax\newcommand{\Ord}{\operatorname{Ord}}
\let\cardred\relax\newcommand{\cardred}{\mathop{\textup{card\,réd}}\nolimits}
\let\cercle\relax\newcommand{\cercle}[1]{\tikz[baseline=(c.base)]\node[draw,circle,inner sep=1pt](c){\textup{#1}};}
\let\Fl\relax\newcommand{\Fl}{\operatorname{Fl}}
\let\Hom\relax\newcommand{\Hom}{\operatorname{Hom}}
\let\id\relax\newcommand{\id}{\mathrm{id}}
\let\limind\relax\newcommand{\limind}{\varinjlim}
\let\defeq\relax\newcommand{\defeq}{\overset{\text{déf}}{=}}
\let\uHom\relax\newcommand{\uHom}{\underline{\mathrm{Hom}}}
\let\End\relax\newcommand{\End}{\operatorname{End}}
\let\edbr\relax\newcommand{\edbr}[1]{\textup{\texttt{#1}}}
\let\uFl\relax\newcommand{\uFl}{\underline{\mathrm{Fl}}}
\let\uIs\relax\newcommand{\uIs}{\underline{\mathrm{Is}}}
\let\uMon\relax\newcommand{\uMon}{\underline{\mathrm{Mon}}}
\let\uEp\relax\newcommand{\uEp}{\underline{\textup{Ép}}}
\let\uGamma\relax\newcommand{\uGamma}{\underline{\Gamma}}
\let\Ens\relax\newcommand{\Ens}{\mathrm{Ens}}
\let\pr\relax\newcommand{\pr}{\mathrm{pr}}
\let\diag\relax\newcommand{\diag}{\mathrm{diag}\,}
\let\cocart\relax\newcommand{\cocart}{\mathrm{cocart}}
\let\cart\relax\newcommand{\cart}{\mathrm{cart}}
\let\cc\relax\newcommand{\cc}{\mathrm{c}}
\let\bDelta\relax\newcommand{\bDelta}{\boldsymbol{\Delta}}
\let\cE\relax\newcommand{\cE}{\mathcal{E}}
\let\cC\relax\newcommand{\cC}{\mathcal{C}}
\let\cM\relax\newcommand{\cM}{\mathcal{M}}

% Reader layout only; mathematical readings of the edition are retained.

% ---- from der-ch18_p121-125.tex
\printedpage{119}
% head: A. Grothendieck : Dérivateurs, Ch. XVIII, Édition des Universités de Montpellier II et Paris VII
satisfaisant aux contraintes $(**)$, $(***)$, $(****)$. On va faire la construction en trois étapes.

\circled{1} Départ :
\[
i_0 \;=\; i\ .
\]
Comme $J'$ est cofinal dans $J$, existe $j_0\in\Ob J'$ et
\[
\varphi(i_0)\xrightarrow{u_0} j_0\ .
\]

\msp{page 316}

\circled{2} Supposons que les $i_\beta$, $j_\beta$ soient construits pour $\beta<\alpha$, ainsi que les $\tau_{\beta',\beta}:i_\beta\longrightarrow i_{\beta'}$, $\tau'_{\beta',\beta}:j_\beta\longrightarrow j_{\beta'}$ pour $\beta,\beta'<\alpha$, enfin les $u_\beta:\varphi(i_\beta)\longrightarrow j_\beta$, $v_\beta:j_\beta\longrightarrow\varphi(i_{\beta+1})$ pour $\beta<\alpha$ resp. pour $\beta+1<\alpha$, de façon à satisfaire aux contraintes $(**)$, $(****)$, i.e. commutativités dans
\[
\left\{
\begin{array}{l}
\left.
\begin{array}{ll}
\circled{1} &
\begin{tikzcd}[ampersand replacement=\&]
\varphi(i_\beta) \arrow[r, "u_\beta"] \arrow[rr, bend right=15, "\varphi(\tau_{\beta+1,\beta})"'] \& j_\beta \arrow[r, "v_\beta"] \& \varphi(i_{\beta+1})
\end{tikzcd}
\\[5ex]
\circled{2} &
\begin{tikzcd}[ampersand replacement=\&]
j_\beta \arrow[r, "v_\beta"] \arrow[rr, dotted, bend right=15, "\tau'_{\beta+1,\beta}"'] \& \varphi(i_{\beta+1}) \arrow[r, "u_{\beta+1}"] \& j_{\beta+1}
\end{tikzcd}
\end{array}
\right\}\quad \textit{si}\ \beta+1<\alpha
\\[8ex]
\begin{array}{ll}
\circled{3} &
\begin{tikzcd}[ampersand replacement=\&]
\varphi(i_\beta) \arrow[r, "u_\beta"] \arrow[d, "\varphi(\tau_{\beta',\beta})"'] \& j_\beta \arrow[d, "\tau'_{\beta',\beta}"] \\
\varphi(i_{\beta'}) \arrow[r, "u_{\beta'}"] \& j_{\beta'}
\end{tikzcd}
\qquad\qquad \text{si } \beta<\beta'<\alpha
\end{array}
\\[10ex]
\begin{array}{ll}
\circled{4} &
\begin{tikzcd}[ampersand replacement=\&]
j_\beta \arrow[r, "v_\beta"] \arrow[d, "\tau'_{\beta',\beta}"'] \& \varphi(i_{\beta+1}) \arrow[d, "\varphi(\tau_{\beta'+1,\beta+1})"] \\
j_{\beta'} \arrow[r, "v_{\beta'}"] \& \varphi(i_{\beta'+1})
\end{tikzcd}
\qquad\qquad \text{si } \beta<\beta' \text{ et } \beta'+1<\alpha\ ,
\end{array}
\end{array}
\right.
\]

\msp{page 317}

\dgrblock{dgrb1079}{The size convention and the purpose of the zigzag}
\begin{dcomment}{exposition}{p316}{The size convention and the purpose of the zigzag}
Grothendieck is trying to recognize a large object by the small pieces of a chosen presentation. The difficulty is that two cofinal presentations need not have any pieces in common. The construction begun on the preceding manuscript pages supplies a common piece by alternately passing from one presentation to the other, and then taking a limit of the resulting process. This is the technical engine of the membership criterion in 4.13.22.

We first fix the translation of the cardinal notation. In this chapter a category is \emph{grande devant} the infinite cardinal $\pi$ when every diagram having at most $\pi$ arrows admits a cocone. In present terminology it is $\lambda$-filtered, where $\lambda=\pi^+$. Accordingly, a $\pi$-accessible object is $\lambda$-presentable: its represented functor preserves $\lambda$-filtered colimits. A useful cardinal for $\mathcal M$ says that $\mathcal M$ is generated under these colimits by its essentially small full subcategory $\mathcal M_\pi$. We shall retain $\pi$ in quotations and use $\lambda$ when the regular-cardinal convention makes a statement easier to compare with modern literature. Replacing $\pi$ by a regular cardinal with the same name would lose the distinction between diagrams of size $\leq\pi$ and diagrams of size $<\pi$.

Here $\varphi:I\longrightarrow J$ is final, $I$ is filtered, and $J'\subset J$ is a full, isomorphism-closed, final subcategory. The desired conclusion is that the strict inverse image $I'=\varphi^{-1}(J')$ is final in $I$. It will be enough to reach an object of $I'$ from each prescribed $i\in I$: a full subcategory of a filtered category with this property is itself filtered and final. This familiar criterion is the one recorded in \citet[Lemma 4.19.3]{StacksFiltered}; throughout these pages cofinality has this ordinary categorical meaning.

Finality alone is insufficient. Let $J=\mathbb N$ with its usual order, take $I$ to be the even natural numbers, and take $J'$ to be the odd natural numbers. Both inclusions are final, but $I'$ is empty. The additional hypothesis in 4.13.17 is therefore doing real work: there is a nonzero limit ordinal $\theta$ such that the relevant $\theta$-chains have colimits, those colimits are preserved by $\varphi$ and $J'\hookrightarrow J$, and the shorter initial segments admit enough cocones. The even--odd example has no colimit for its unbounded countable chain.

Starting at $i$, the displayed maps have the form
\[
 \varphi(i_\beta)\xrightarrow{u_\beta}j_\beta
 \xrightarrow{v_\beta}\varphi(i_{\beta+1}).
\]
They do not assert that any intermediate $\varphi(i_\beta)$ belongs to $J'$. Instead, the two composites are required to be the transition maps of the two chains. At the final limit stage the chains consequently have isomorphic colimits in $J$. One colimit comes from $I$, the other belongs to $J'$; isomorphism-closure then puts the former in $J'$. This explains both the alternation and the insistence on a strictly full subcategory.
\end{dcomment}


il faut construire
% foot: Version 17 décembre 2025

\printedpage{120}
% head: A. Grothendieck : Dérivateurs, Ch. XVIII, Édition des Universités de Montpellier II et Paris VII
\begin{itemize}[leftmargin=*]
\item[1)] $i_\alpha\in\Ob I$,\quad et un homomorphisme $(i_\beta)_{\beta\in I_\alpha}\longrightarrow i_\alpha$, i.e. des $\tau_{\alpha,\beta}$ : $i_\beta\longrightarrow i_\alpha$ $(\beta\in I_\alpha)$, compatibles aux translations $\tau_{\beta',\beta}$ pour $\beta<\beta'<\alpha$.
\item[2)] $j_\alpha\in\Ob J'$,\quad et un homomorphisme $(j_\beta)_{\beta\in I_\alpha}\longrightarrow j_\alpha$, i.e. des $\tau'_{\alpha,\beta}$ : $j_\beta\longrightarrow j_\alpha$ $(\beta\in I_\alpha)$, compatibles aux translations $\tau'_{\beta',\beta}$ pour $\beta<\beta'<\alpha$.
\item[3)] Une flèche
\[
u_\alpha:\varphi(i_\alpha)\longrightarrow j_\alpha\ .
\]
\item[4)] Si $\alpha$ n'est pas ordinal limite, i.e. $\alpha=\alpha_0+1$, une flèche
\[
v_{\alpha_0}:j_{\alpha_0}\longrightarrow\varphi(i_{\alpha_0+1})=\varphi(i_\alpha)\ ;
\]
\end{itemize}

et ceci de façon à satisfaire aux contraintes \circled{1}\,\circled{2}\,\circled{3}\,\circled{4} quand on y suppose $\beta+1\leqslant\alpha$ (au lieu de $\beta'<\alpha$) dans \circled{1} et \circled{2}, et $\beta'\leqslant\alpha$ (au lieu de $\beta'<\alpha$) dans \circled{3}, et $\beta'+1\leqslant\alpha$ (au lieu de $\beta'+1<\alpha$) dans \circled{4}. Donc pour les contraintes $1^\circ$, $2^\circ$, $4^\circ$, il n'y a en tenir compte \textit{que} si $\alpha$ n'est pas ordinal limite, donc $\alpha=\alpha_0+1$, et en y faisant % [sic] il n'y a en tenir -> il n'y a à en tenir, the preposition « à » appears to be missing
\[
\beta' \;=\; \alpha_0\ .
\]
Quant à la contrainte $3^\circ$, il faut en
\msp{page 318}
tenir compte pour $\beta'=\alpha$. Donc au total, il faut exiger
\[
(3_\alpha)(\beta)\qquad\qquad
\begin{tikzcd}
\varphi(i_\beta) \arrow[r, "u_\beta"] \arrow[d, "\varphi(\circled{$\tau_{\alpha,\beta}$})"'] & j_\beta \arrow[d, "\tau'_{\alpha,\beta}" {draw, circle, inner sep=1pt}] \\
\varphi(i_\alpha) \arrow[r, "u_\alpha" {draw, circle, inner sep=1pt}] & j_\beta % [sic] j_\beta -> j_\alpha, the target of \tau'_{\alpha,\beta} and u_\alpha is j_\alpha
\end{tikzcd}
\qquad \text{commutatif pour } \beta<\alpha\ (\footnotemark[140]),
\]
\footnotetext[140]{\textbf{N.B.} Il suffit de l'exiger pour $\beta=\alpha_0$, dans le cas $\alpha=\alpha_0+1$, cf. ci-dessous, p. 323.}
(par restriction sur $\alpha$), et si $\alpha=\alpha_0+1$ est ordinal non limite, il faut exiger de plus
\[
(4_\alpha)(\beta)\qquad\qquad
\begin{tikzcd}
j_\beta \arrow[r, "v_\beta"] \arrow[d, "\tau'_{\alpha_0,\beta}"'] & \varphi(i_{\beta+1}) \arrow[d, "\varphi(\circled{$\tau_{\alpha,\beta+1}$})"] \\
j_{\alpha_0} \arrow[r, "v_{\alpha_0}" {draw, circle, inner sep=1pt}] & \varphi(i_{\alpha_0+1})=\varphi(i_\alpha)
\end{tikzcd}
\qquad \beta<\alpha_0
\]
\[
(1_\alpha)\qquad\qquad
\begin{tikzcd}
\varphi(i_{\alpha_0}) \arrow[r, "u_{\alpha_0}"] \arrow[rr, bend right=15, "\varphi(\circled{$\tau_{\alpha,\alpha_0}$})"'] & j_{\alpha_0} \arrow[r, "v_{\alpha_0}" {draw, circle, inner sep=1pt}] & \varphi(i_{\alpha_0+1})=\varphi(i_\alpha)
\end{tikzcd}
\]
% foot: Version 17 décembre 2025

\printedpage{121}
% head: A. Grothendieck : Dérivateurs, Ch. XVIII, Édition des Universités de Montpellier II et Paris VII
\[
(2_\alpha)\qquad\qquad
\begin{tikzcd}
j_{\alpha_0} \arrow[r, "v_{\alpha_0}" {draw, circle, inner sep=1pt}] \arrow[rr, bend right=15, "\tau'_{\alpha,\alpha_0}"' {draw, circle, inner sep=1pt}] & \varphi(i_{\alpha_0+1})=\varphi(i_\alpha) \arrow[r, "u_\alpha" {draw, circle, inner sep=1pt}] & j_\alpha\ .
\end{tikzcd}
\]

Dans ces compatibilités, j'ai cerclé les flèches $\tau_{\alpha,\beta}$, $\tau'_{\alpha,\beta}$ $(\beta<\alpha)$, $\Bigl(\ \tau_{\alpha,\beta+1}\ (\beta+1<\alpha)\ \Bigr)$, $v_{\alpha_0}$, $\Bigl(\ \tau_{\alpha,\alpha_0}\ \Bigr)$, $u_\alpha$, $\Bigl(\ \tau'_{\alpha,\alpha_0}\ \Bigr)$, qui ont été introduites dans les donnés 1) 2) 3) 4), page 317. % [sic] donnés -> données, the noun « données » is meant

\msp{page 319}

On va distinguer donc les deux cas.

\textit{2 A)} $\alpha=\alpha_0+1$.

Pour construire $i_\alpha$ et les $\tau_{\alpha,\beta}:i_\beta\longrightarrow i_\alpha$, il suffit de construire $i_\alpha$ et
\[
\circled{$\tau_{\alpha,\alpha_0}$}\ :\ i_{\alpha_0}\ \longrightarrow\ i_\alpha=i_{\alpha_0+1}\ .
\]
De même pour $j_\alpha$ et les $\tau'_{\alpha,\beta}$, il suffit de construire
\[
\circled{$\tau'_{\alpha,\alpha_0}$}\ :\ j_{\alpha_0}\ \longrightarrow\ j_\alpha\ .
\]
Il faut de plus se donner
\[
j_{\alpha_0}\ \xrightarrow{\circled{$v_{\alpha_0}$}}\ \varphi(\underbrace{\,i_\alpha\,}_{=i_{\alpha_0+1}})\ \xrightarrow{\circled{$u_\alpha$}} j_\alpha\ ,
\]
donc en tout deux objets $i_\alpha$, $j_\alpha$ (dans $I$ et dans $J'$ respectivement) et quatre flèches $\tau_{\alpha,\alpha_0}$, $\tau'_{\alpha,\alpha_0}$, $v_{\alpha_0}$, $u_\alpha$. Il faut quatre compatibilités (page 318), où les deux dernières $3_\alpha(\beta)$, $4_\alpha(\beta)$ dépendent d'un paramètre $\beta<\alpha$, i.e. $\beta\leqslant\alpha_0$, ou de $\beta<\alpha_0$.

\msp{page 320 biffé}

\msp{page 321}

On va d'abord s'occuper de deux premières compatibilités $1_\alpha$, $2_\alpha$, en termes de la donnée de
\[
i_{\alpha_0}\ \in\ \Ob I\ ,\qquad j_{\alpha_0}\ \in\ \Ob J'\ ,\qquad u_{\alpha_0}\ :\ \varphi(i_{\alpha_0})\ \longrightarrow\ j_{\alpha_0}\ .
\]
Les deux compatibilités voulues pour les deux objets $i_\alpha$, $j_\alpha$ et les quatre flèches $\tau_{\alpha,\alpha_0}$, $\tau'_{\alpha,\alpha_0}$, $v_{\alpha_0}$, $u_\alpha$ sont indiquées dans le diagramme
\[
(*)\qquad\qquad
\begin{tikzcd}
\varphi(i_{\alpha_0}) \arrow[r, "u_{\alpha_0}"] \arrow[rr, bend right=15, "\varphi(\tau_{\alpha,\alpha_0})"'] & j_{\alpha_0} \arrow[r, "v_{\alpha_0}\ ?"] \arrow[rr, bend left=20, "\tau'_{\alpha,\alpha_0}"] & \varphi(\underbrace{\,i_\alpha\,}_{?}) \arrow[r, "u_\alpha\ ?"] & j_\alpha\ .
\end{tikzcd}
\]

On note que comme $\varphi$ est cofinal, $\exists\ i_\alpha=i\in\Ob I$ et $j_{\alpha_0}\xrightarrow{v_{\alpha_0}}\varphi(i)$. Mais on veut que le composé
\[
w\ :\ \varphi(i_{\alpha_0})\ \xrightarrow{u_{\alpha_0}}\ j_{\alpha_0}\ \xrightarrow{v_{\alpha_0}}\ \varphi(i)
\]
soit de la forme $\varphi(\tau)$, où $\tau=\tau_{\alpha,\alpha_0}:i_{\alpha_0}\longrightarrow i_\alpha$ est une flèche dans $I$. Comme $\varphi$ est un foncteur cofinal de catégories filtrantes, on sait
% foot: Version 17 décembre 2025
\printedpage{122}
% head: A. Grothendieck : Dérivateurs, Ch. XVIII, Édition des Universités de Montpellier II et Paris VII
\msp{page 322}
qu'on peut trouver $i\xrightarrow{\rho}i'$ dans $I$, tel que la flèche composée $\varphi(i_{\alpha_0})\xrightarrow{w}\varphi(i)\xrightarrow{\varphi(\rho)}\varphi(i')$ soit de la forme $\varphi(\tau)$. Donc quitte à remplacer $i$ par $i'$, et $v_{\alpha_0}$ par le composé $v'_{\alpha_0}$ : $j_{\alpha_0}\xrightarrow{v_{\alpha_0}}\varphi(i)\xrightarrow{\varphi(\rho)}\varphi(i')$, on peut supposer que $v_{\alpha_0}u_{\alpha_0}$ est de la forme $\varphi(\tau)$. Cela nous donne donc en principe l'objet $i=i_\alpha$, et les flèches $v_{\alpha_0}:j_{\alpha_0}\longrightarrow\varphi(i)$ et $\tau_{\alpha,\alpha_0}:i_{\alpha_0}\longrightarrow i_\alpha$, de façon à satisfaire à la première compatibilité dans $(*)$. Il reste à construire $j_\alpha$, $\tau'_{\alpha,\alpha_0}$, $u_\alpha$ de façon à satisfaire à la deuxième compatibilité $(*)$. Or $J'$ étant cofinal dans $J$, il existe $j_\alpha$ dans $J'$ et une flèche $u_\alpha:\varphi(i_\alpha)\longrightarrow j_\alpha$. D'autre part, comme $J'\longrightarrow J$ est pleinement fidèle, la flèche composée $j_{\alpha_0}\xrightarrow{v_{\alpha_0}}\varphi(i_\alpha)\xrightarrow{u_\alpha}j_\alpha$ dans $J$, entre objets de $J'$, provient d'une unique flèche $\tau'_{\alpha,\alpha_0}:j_{\alpha_0}\longrightarrow j_\alpha$ dans $J'$.

\msp{page 323}

\dgrblock{dgrb1080}{Finality supplies eventual lifting, not fullness}
\begin{dcomment}{exposition}{p322}{Finality supplies eventual lifting, not fullness}
The successor construction needs to turn a map between images into the image of a map, after moving farther along the indexing category. Grothendieck invokes this several times. Writing out the small lemma makes clear why fullness of $\varphi$ is unnecessary.

Suppose $I$ is filtered and $\varphi:I\to J$ is final. Given $w:\varphi(a)\to\varphi(b)$, there exist arrows $r:b\to c$ and $t:a\to c$ in $I$ such that
\[
                  \varphi(r)w=\varphi(t).
\]
Indeed, the objects $(a,1_{\varphi(a)})$ and $(b,w)$ of the comma category $(\varphi(a)\downarrow\varphi)$ are connected by a finite zigzag. Forgetting to $I$, choose a cocone on that finite zigzag. After applying $\varphi$, the equations defining its comma morphisms say that every composite from $\varphi(a)$ to the common target is the same. The two endpoints give $t$ and $r$. This is a direct use of the connected-comma definition of finality, rather than a lifting property for individual arrows; compare \citet[Definition 4.17.1]{StacksCofinal}.

Apply the lemma after first choosing $j_{\alpha_0}\to\varphi(b)$ by finality. Its composite with $u_{\alpha_0}$ then becomes $\varphi(t)$ after a further transition. This constructs $i_\alpha$, $v_{\alpha_0}$ and $\tau_{\alpha,\alpha_0}$ with
\[
v_{\alpha_0}u_{\alpha_0}=\varphi(\tau_{\alpha,\alpha_0}).
\]
Finality of $J'$ now supplies $u_\alpha:\varphi(i_\alpha)\to j_\alpha$ with $j_\alpha\in J'$. Fullness of $J'\hookrightarrow J$ permits us to define its transition arrow by
\[
\tau'_{\alpha,\alpha_0}=u_\alpha v_{\alpha_0}.
\]
Only this last passage uses fullness. It is useful to separate it from the earlier use of finality, since the proof would otherwise appear to assume that $\varphi$ were fully faithful.
\end{dcomment}


Ainsi, on a construit les deux objets et les quatre flèches voulues, de façon à satisfaire aux contraintes $1_\alpha$, $2_\alpha$ de la page 318. Les contraintes $3_\alpha(\beta)$, $4_\alpha(\beta)$ sont-elles alors satisfaites automatiquement ?
\[
\begin{tikzcd}
\varphi(i_\beta) \arrow[r, "u_\beta"] \arrow[d] & j_\beta \arrow[d] \\
\varphi(i_{\alpha_0}) \arrow[r, "u_{\alpha_0}"] \arrow[d, "\varphi(\tau_{\alpha,\alpha_0})"'] & j_{\alpha_0} \arrow[d, "\tau'_{\alpha,\alpha_0}"] \\
\varphi(i_\alpha) \arrow[r, "u_\alpha"] & j_\alpha
\end{tikzcd}
\qquad\qquad \beta<\alpha, \text{ i.e. } \beta\leqslant\alpha_0\ .
\]
Le carré supérieur est commutatif par hypothèse de récurrence, pour que le carré total (composé) soit commutatif, il suffit que le carré inférieur le soit. On a donc une double-flèche $(\tau'_{\alpha,\alpha_0}u_{\alpha_0},u_\alpha\varphi(\tau_{\alpha,\alpha_0})):\varphi(i_\alpha)\rightrightarrows j_\alpha$, elle n'est peut-être pas triviale, mais on sait, $J'$ étant filtrante, qu'on peut l'égaliser par une flèche $j_\alpha\longrightarrow\bar{j}_\alpha$ dans $J'$. Donc quitte à remplacer $j_\alpha$ par % [sic] \varphi(i_\alpha) -> \varphi(i_{\alpha_0}), both composites of the double arrow start at \varphi(i_{\alpha_0})
\msp{page 324}
$\bar{j}_\alpha$, $u_\alpha$ par le composé $\bar{u}_\alpha:\varphi(i_0)\longrightarrow j_\alpha\longrightarrow\bar{j}_\alpha$ et $\tau'_{\alpha,\alpha_0}$ par le composé $\bar{\tau}'_{\alpha,\alpha_0}:j_{\alpha_0}\longrightarrow j_\alpha\longrightarrow\bar{j}_\alpha$ (sans toucher à $i_\alpha$, $v_{\alpha_0}$, $\tau_{\alpha_0,\alpha}$), on peut obtenir la commutativité, d'où $3_\alpha(\beta)$. % [sic] \varphi(i_0) -> \varphi(i_\alpha), u_\alpha has source \varphi(i_\alpha); [sic] \tau_{\alpha_0,\alpha} -> \tau_{\alpha,\alpha_0}, the arrow is written \tau_{\alpha,\alpha_0} everywhere else

Reste la contrainte $4_\alpha(\beta)$. Pour ceci, considérons le diagramme composite
\[
\begin{tikzcd}
j_\beta \arrow[r, "v_\beta"] \arrow[d, "\tau'_{\alpha_0,\beta}"'] & \varphi(i_{\beta+1}) \arrow[d, "\varphi(\tau_{\alpha,\beta+1})"] \\
j_{\alpha_0} \arrow[r, "v_{\alpha_0}"] & \varphi(i_{\alpha_0+1})=\varphi(i_\alpha)\ .
\end{tikzcd}
\]

Si $\alpha_0$ n'est pas un ordinal limite, mais de la forme $\alpha_1+1$, on peut encore se ramener au cas où $\beta=\alpha_1$, en procédant comme ci-dessus. Mais on peut procéder dans tous les cas, sans un tel artifice (\,!), en considérant les double-flèches
\[
(*)\qquad\qquad
\begin{tikzcd}
j_\beta \arrow[r, shift left, "\varphi(\tau_{\alpha,\beta+1})v_\beta"] \arrow[r, shift right, "v_{\alpha_0}\tau'_{\alpha_0,\beta}"'] & \varphi(i_\alpha)
\end{tikzcd}
\]
% foot: Version 17 décembre 2025
\printedpage{123}
% head: A. Grothendieck : Dérivateurs, Ch. XVIII, Édition des Universités de Montpellier II et Paris VII
déduites du carré précédent. Comme
\msp{page 325}
l'ensemble de ces double-flèches est de cardinal $\card\alpha<\pi$, et que $J$ est grande devant $\pi$ (il suffirait qu'il le soit pour les cardinaux $\pi'<\pi$), on peut trouver une flèche $\varphi(i_\alpha)\xrightarrow{u}j$ dans $J$ qui égalise toutes ces double-flèches. Comme $\varphi$ est cofinal, on peut supposer que $j$ est de la forme $\varphi(i)$. Comme $\varphi$ est un foncteur cofinal de catégories filtrantes, il existe alors $i\xrightarrow{\rho}i'$ dans $I$, tel que le composé $\varphi(i_\alpha)\xrightarrow{u}\varphi(i)\xrightarrow{\varphi(\rho)}\varphi(i')$ soit de la forme $\varphi(\sigma)$, pour $\sigma:i_\alpha\longrightarrow i'$. Au total, on a trouvé $\sigma:i_\alpha\longrightarrow i'$, tel que $\varphi(\sigma)$ $(=\varphi(\rho)u):\varphi(i_\alpha)\longrightarrow\varphi(i')$ égalise toutes les double-flèches envisagées dans $(*)$, p. 324. Donc il s'impose de remplacer $i_\alpha$ par $\bar{i}_\alpha=i'$, $v_{\alpha_0}$ par le composé $j_{\alpha_0}\longrightarrow\varphi(i_\alpha)\xrightarrow{\varphi(\sigma)}\varphi(\bar{i}_\alpha)$, $\tau_{\alpha,\alpha_0}$
\msp{page 326}
par le composé $\bar{\tau}_{\alpha,\alpha_0}:i_{\alpha_0}\xrightarrow{\tau_{\alpha,\alpha_0}}i_\alpha\xrightarrow{\sigma}\bar{i}_\alpha$.

Pour résumer, voici la recette pour construire, d'abord \ovalbox{$i_\alpha$, $v_{\alpha_0}$, $\tau_{\alpha,\alpha_0}$,} et \textit{ensuite} \ovalbox{$j_\alpha$, $u_\alpha$, $\tau'_{\alpha,\alpha_0}$}.

On procède d'abord comme page 321, 322 pour construire le premier triple, de façon à satisfaire à la première compatibilité $(*)$ (p. 321), i.e. à $1_\alpha$. On “corrige” ensuite ce choix de $i_\alpha$, $v_{\alpha_0}$, $\tau_{\alpha,\alpha_0}$, en remplaçant $i_\alpha$ par $\bar{i}_\alpha$ pour $\sigma:i_\alpha\longrightarrow\bar{i}_\alpha$, $v_{\alpha_0}$ par $\bar{v}_{\alpha_0}=\varphi(\sigma)v_{\alpha_0}$, et $\tau_{\alpha,\alpha_0}$ par $\bar{\tau}_{\alpha,\alpha_0}=\sigma\tau_{\alpha,\alpha_0}$, de sorte à obtenir les compatibilités $4_\alpha(\beta)$ pour $\beta<\alpha_0$, comme on vient de l'indiquer. On ne touchera plus à ces choix de $i_\alpha$, $v_{\alpha_0}$, $\tau_{\alpha,\alpha_0}$, dans la suite de la construction. On construit $(j_\alpha,u_\alpha,\tau'_{\alpha,\alpha_0})$ de façon à satisfaire à la deuxième compatibilité $(*)$ (p. 321), comme
\msp{page 327}
indiqué page 322. Ensuite, on corrige ce choix, en remplaçant $j_\alpha$ par $\bar{j}_\alpha$, avec $\rho:j_\alpha\longrightarrow\bar{j}_\alpha$, et $u_\alpha$ par $\bar{u}_\alpha=\rho u_\alpha$, $\tau'_{\alpha,\alpha_0}$ par $\bar{\tau}'_{\alpha,\alpha_0}=\rho\tau'_{\alpha,\alpha_0}$, de façon à satisfaire aux contraintes $3_\alpha(\beta)$ de la page 318.

Cela achève la construction dans le cas où $\alpha=\alpha_0+1$.

\textit{Cas 2 B : $\alpha$ est ordinal limite.}

Il faut donc construire les données 1) 2) 3) de la page 317 (à l'exclusion de la donnée 4), de façon à satisfaire à la seule contrainte $3_\alpha(\beta)$.

On commence à construire $i_\alpha$ et les $\tau_{\beta,\alpha}$ dans $I$. C'est possible, car $I_\alpha=\{\beta\mid\beta<\alpha\}$ est de cardinal $\pi_\alpha<\pi$ – il suffit donc de supposer $I$ grande devant les cardinaux $\pi'<\pi$.

On construit ensuite de même $J_\alpha$ dans $J'$ et les $\tau'_{\beta,\alpha}$ $(\beta<\alpha)$, en utilisant le % [sic] J_\alpha -> j_\alpha, the object of J' is meant
\msp{page 328}
fait que $J'$ est grande devant $\card\alpha=\card I_\alpha$. On doit de plus construire une flèche
\[
u_\alpha\ :\ \varphi(i_\alpha)\ \longrightarrow\ j_\alpha\ .
\]
C'est possible, car $J'$ étant cofinale dans $J$, il existe $j'$ dans $J'$ et $\varphi(i_\alpha)\longrightarrow J'$, et $J'$ étant filtrante, il existe $\bar{j}_\alpha$ et des flèches % [sic] \varphi(i_\alpha)\longrightarrow J' -> \varphi(i_\alpha)\longrightarrow j', the arrow goes to the object j'
$\begin{tikzcd}[ampersand replacement=\&, baseline=(\tikzcdmatrixname-1-1.base)]
j_\alpha \arrow[rr] \& \& j_\alpha \\
\& j' \arrow[ur] \&
\end{tikzcd}$, donc quitte à remplacer $j_\alpha$ par $\bar{j}_\alpha$, les $\tau'_{\alpha,\beta}$ par des $\bar{\tau}'_{\alpha,\beta}$, on trouve bien une flèche $u_\alpha:\varphi(i_\alpha)\longrightarrow j_\alpha$. Il faut ensuite s'occuper des % [sic] target j_\alpha in the small diagram -> \bar{j}_\alpha, the two arrows go to the object \bar{j}_\alpha just introduced
% foot: Version 17 décembre 2025

\dgrblock{dgrb1081}{The successor compatibilities are already forced}
\begin{dcomment}{development}{p327}{The successor compatibilities are already forced}
Grothendieck checks the four displayed compatibilities by successive corrections. His objective is sound, but once the two adjacent composites have been fixed, the remaining squares at a successor stage already commute. The following calculation shortens the construction without changing its hypotheses.

\textbf{Successor simplification.}\ours\ Let $\alpha=\delta+1$, assume all preceding transition maps and $u$'s satisfy the stated compatibilities, and choose
\[
 v_\delta u_\delta=\varphi(\tau_{\alpha,\delta}),
 \qquad
 u_\alpha v_\delta=\tau'_{\alpha,\delta}.
\]
Extend the transition maps to earlier indices by composition. Then both $(3_\alpha)(\beta)$ and $(4_\alpha)(\beta)$ hold automatically.

For the first square, the old naturality relation for $\beta<\delta$, and the identity relation when $\beta=\delta$, give
\[
\begin{aligned}
\tau'_{\alpha,\beta}u_\beta
 &=u_\alpha v_\delta\tau'_{\delta,\beta}u_\beta\\
 &=u_\alpha v_\delta u_\delta\varphi(\tau_{\delta,\beta})\\
 &=u_\alpha\varphi(\tau_{\alpha,\beta}).
\end{aligned}
\]
For the second, let $\beta<\delta$. Since $\beta+1\leq\delta$, the earlier second compatibility and naturality give
\[
\begin{aligned}
v_\delta\tau'_{\delta,\beta}
 &=v_\delta\tau'_{\delta,\beta+1}
                         u_{\beta+1}v_\beta\\
 &=v_\delta u_\delta\varphi(\tau_{\delta,\beta+1})v_\beta\\
 &=\varphi(\tau_{\alpha,\beta+1})v_\beta .
\end{aligned}
\]
The same equations cover $\beta+1=\delta$, using the identity transition. This proves the assertion.

Thus the extra equalizations at manuscript pages 323--327 can be omitted at successor stages. They are not a counterexample to the construction: applying them to an already equal pair causes no harm, provided the later choices are made in the order specified on page 326. In particular, the phrase suggesting that the pair on page 323 may differ should not be read as evidence of an unresolved obstruction.

At a genuine limit stage there is no preceding last object, and a simultaneous equalization still has a purpose. One chooses cocones on the two earlier chains, adds a map $\varphi(i_\alpha)\to j_\alpha$ after enlarging $j_\alpha$, and equalizes the resulting family of parallel pairs. The arrows have sources in $J$, not necessarily in $J'$, so this equalization is first performed in $J$; finality of $J'$ then moves its target back into $J'$. The required filteredness of $J$ follows from that of $I$ and finality of $\varphi$.

We checked the cofinality and eventual-equality lemmas in \citet{StacksCofinal,StacksFiltered}, and Grothendieck's complete construction. We did not locate this particular simplification as a separately stated result. The marking records our proof in this commentary, without a claim of priority.
\end{dcomment}


% ---- from der-ch18_p126-130.tex
\printedpage{124}
% head: A. Grothendieck : Dérivateurs, Ch. XVIII, Édition des Universités de Montpellier II et Paris VII

compatibilités $3_\alpha(\beta)$, donc des double-flèches
\[
\begin{tikzcd}
\varphi(i_\beta) \arrow[r, shift left, "\tau'_{\alpha,\beta}u_\beta"] \arrow[r, shift right, "u_\alpha\varphi(\tau_{\alpha,\beta})"'] & j_\alpha \; .
\end{tikzcd}
\]
Comme l'ensemble de ces double-flèches est de cardinal $\card\alpha < \pi$, et que $J'$ est grande devant ledit cardinal, on
\msp{page 329}
trouve donc une flèche
\[
\rho : j_\alpha \longrightarrow j
\]
dans $J$ qui les égalise toutes. Comme $J'$ est cofinal, on peut supposer $j$ dans $J'$, soit $\bar{j}_\alpha$. Donc quitte à remplacer $j_\alpha$ par $\bar{j}_\alpha$, et les $\tau'_{\alpha,\beta}$ par $\bar{\tau}'_{\alpha,\beta} = \rho\tau'_{\alpha,\beta}$, enfin $u_\alpha : \varphi(i_\alpha) \longrightarrow j_\alpha$ par $\bar{u}_\alpha = \rho u_\alpha : \varphi(i_\alpha) \longrightarrow \bar{j}_\alpha$, on satisfait à toutes les contraintes voulues.

Je résume le résultat (``$\pm$ tautologique''!) prouvé dans les pages 313-329.

% a vertical change bar is printed in the left margin alongside the statement from "Lemme-clef" to "suivantes :"
\textbf{Lemme-clef 4.13.17.} \textit{Considérons un carré cartésien dans} Cat
\[
\begin{tikzcd}
I \arrow[r, "\varphi"] & J \\
I' \arrow[u, hook] \arrow[r] & J' \arrow[u, hook, "\psi"'] \; .
\end{tikzcd}
\]

\textit{On fait les hypothèses suivantes :}

\msp{page 330}

\begin{enumerate}[label={}]
\item[\textit{1)}] \textit{$\varphi$ et $\psi$ sont cofinaux, $\psi$ est l'inclusion d'une sous-catégorie strictement pleine $J'$ de $J$, $I$ est filtrant (donc $J$ et $J'$ aussi).}
\end{enumerate}

\textit{On cherche des conditions pour que l'inclusion de la sous-catégorie strictement pleine $I' \longrightarrow I$ soit cofinale (auquel cas on voit tout de suite que $I'$ est filtrante, et $I' \longrightarrow J'$ est aussi cofinal). Pour ceci, on suppose qu'il existe un \textup{ordinal limite} $\alpha_0$ ayant les propriétés supplémentaires suivantes :}

\begin{enumerate}[label={}]
\item[\textit{2)}] \textit{$I$ et $J'$ sont stables par $\varinjlim$ de type $I_{\alpha_0}$ $(= \{\beta \mid \beta \ \textit{ordinal}\ < \alpha_0\})$, et les foncteurs $\varphi$, $\psi$ commutent à ces limites.}
\item[\textit{3)}] \textit{Pour tout $\alpha \in I_{\alpha_0}$, $I$ est grande devant $\card\alpha$. (Donc, $\varphi$ étant cofinal, $J$ est également grande devant $\card\alpha$, et de même, $\psi$ étant pleinement fidèle, et cofinal, $J'$ est grande devant $\card\alpha$.)}
\end{enumerate}

\textit{Sous ces conditions, $I'$ est cofinal dans $I$ (donc aussi grande devant $\card\alpha$, pour tout $\alpha \in I_{\alpha_0}$).}

La démonstration est celle donnée ci-dessus. On n'a pas eu à utiliser que l'ordinal $\alpha_0 = \alpha_\pi$ était de cette espèce particulière, seulement que c'était un cardinal limite.

\printedpage{125}
% head: A. Grothendieck : Dérivateurs, Ch. XVIII, Édition des Universités de Montpellier II et Paris VII

\msp{page 331}

\dgrblock{dgrb1082}{Taking the final colimit completes the cofinality proof}
\begin{dcomment}{exposition}{p330}{Taking the final colimit completes the cofinality proof}
The statement of 4.13.17 collects exactly the ingredients needed to finish the alternating construction. It is useful to make the last passage explicit, because it is the point where approximate membership becomes actual membership.

Write $\theta$ for the limit ordinal in that lemma. For $\beta<\theta$, the preceding construction gives compatible chains $(i_\beta)$ in $I$ and $(j_\beta)$ in $J'$, with $i_0$ prescribed. At each limit stage $\alpha<\theta$ there are at most $\lvert\alpha\rvert$ objects and compatibility equations to accommodate. The filteredness assumption for this cardinal supplies their cocones and simultaneous equalizations. At successor stages the preceding block gives all the compatibilities. Transfinite recursion therefore constructs the entire two chains.

Take their colimits,
\[
 \bar i=\varinjlim_{\beta<\theta}i_\beta\quad\hbox{in }I,
 \qquad
 \bar j=\varinjlim_{\beta<\theta}j_\beta\quad\hbox{in }J'.
\]
Both exist by hypothesis. The functors to $J$ preserve them. The maps $u_\beta$ induce $U:\varphi(\bar i)\to\bar j$. The maps $v_\beta$ induce $V:\bar j\to\varphi(\bar i)$: compatibility $(4)$ is precisely what is needed to define this map. The composite $VU$ agrees on $\varphi(i_\beta)$ with the transition to $\varphi(i_{\beta+1})$ and hence with the colimit structure map. Therefore $VU=1$. Likewise $UV=1$ by compatibility $(2)$. There is always a next index because $\theta$ is a limit ordinal; the successor indices are cofinal in $\theta$.

We have proved $\varphi(\bar i)\simeq\bar j$. Since $J'$ is closed under isomorphisms, $\bar i\in I'$, and the first colimit map gives $i_0\to\bar i$. As $i_0$ was arbitrary and $I'$ is full in the filtered category $I$, the required finality follows. A full final subcategory of a category admitting cocones on diagrams of a specified size again admits such cocones: form a cocone in the ambient category, then move its vertex into the subcategory. This also proves the final parenthesis of 4.13.17.

There is a similarly direct verification that $I'\to J'$ is final. The composite $I'\to I\to J$ is final. For each $j'\in J'$, its comma category with that composite is identical to $(j'\downarrow (I'\to J'))$, because $J'\subset J$ is full. Its nonemptiness and connectedness give the assertion.

A cardinal remark in the subsequent application requires correction. The initial ordinal of cardinality $\pi$ is not large before every smaller cardinal when $\pi$ is singular: a cofinal subset of size $\operatorname{cf}(\pi)<\pi$ has no upper bound. There is a repair which does not enlarge $\pi$: use instead the initial ordinal $\theta=(\pi'_0)^+$. Since $\pi>\pi'_0$, its cardinality is at most $\pi$, it is large before $\pi'_0$, and the hypotheses of adaptedness and $\mathrm L(\pi'_0,\pi)$ provide exactly the $\theta$-colimits needed. Every shorter segment has size at most $\pi'_0$, so the remaining hypotheses of the lemma hold. If one insists on using the initial ordinal of $\pi$, it suffices to assume $\operatorname{cf}(\pi)>\pi'_0$, for example by imposing $\pi^{\pi'_0}=\pi$. The stronger sentence on manuscript page 336 is false at singular cardinals, but the argument has the smaller-ordinal repair just given.
\end{dcomment}


Ce lemme implique, comme on l'a vu à la page 312, l'implication b $\Longrightarrow$ d dans le diagramme d'implications $(*)$ page 310, lequel est donc à présent prouvé.

Pour achever de prouver le théorème 4.13.15, il reste à prouver les implications
\[
\begin{tikzcd}[row sep=small]
\mathrm{c}' \arrow[dr, Rightarrow] & & \\
& \mathrm{a}' \arrow[r, Rightarrow] & \mathrm{a} \; . \\
\mathrm{b} \arrow[ur, Rightarrow] & &
\end{tikzcd}
\]

Démontrons donc (b ou c$'$) $\Longrightarrow$ a$'$. Il reste à prouver que $\underline{S}$ est $\pi$-accessible dans $\mathcal{M}$, i.e. que le foncteur pleinement fidèle
\[
\Ind_\pi(\underline{S}(\pi)) \longrightarrow \mathcal{M} \tag*{$(*)$}
\]
a pour image essentielle $\underline{S}$. Donc il faut prouver que si $x \in \Ob\mathcal{M}$, alors
\[
x \in \Ob\underline{S} \quad \overset{(?)}{\Longleftrightarrow} \quad
\begin{minipage}[c]{0.55\textwidth}
$x$ est dans l'image essentielle de $\Ind_\pi(\underline{S}(\pi))$, i.e. $x$ est isomorphe à $\varinjlim\limits_{i\in I} x_i$, les $x_i$ dans $\Ob\underline{S}(\pi)$, $I$ grande devant $\pi$.
\end{minipage}
\tag*{$(*)$}
\]

\msp{page 332}

$\Bigg($ Supposons d'abord que $\underline{S}$ satisfasse b) (\footnote[141]{Maladroit, on peut traiter les deux cas b) et c) en même temps, cf. bas de page 333.}). Donc on sait que
\[
x \in \Ob\underline{S} \qquad \Longleftrightarrow \qquad \underline{S}/x \text{ est cofinale dans } \mathcal{M}_\pi/x \; . \tag*{$(**)$}
\]

Nous voulons déduire l'équivalence $(*)$ de l'équivalence $(**)$.

Prouvons $\Longrightarrow$ dans $(*)$, i.e. supposons $x \in \Ob\underline{S}$, et prouvons que $x$ est dans l'image essentielle de $\Ind_\pi\underline{S}(\pi)$. En vertu de $(**)$, $\underline{S}(\pi)/x$ est une sous-catégorie (strictement pleine) cofinale dans $\mathcal{M}_\pi/x$, donc comme $\mathcal{M}_\pi/x$ est grande devant $\pi$, $\underline{S}(\pi)/x$ l'est aussi. Comme on sait que
\[
x \ = \ \varinjlim_{i\in\mathcal{M}_\pi/x} x_i
\]
($\pi$ étant cardinal utile pour $\mathcal{M}$), il s'ensuit par cofinalité qu'on a aussi
\[
x \ = \ \varinjlim_{i\in\underline{S}(\pi)/x} x_i \; ,
\]
et comme $\underline{S}(\pi)/x$ est grande devant $\pi$, cela implique que $x$ appartient à l'image essentielle de $\Ind_\pi(\underline{S}(\pi))$.

Prouvons $\Longleftarrow$ dans $(*)$, i.e. supposons que $x$ soit dans l'image essentielle de $\Ind_\pi(\underline{S}(\pi))$,
\printedpage{126}
% head: A. Grothendieck : Dérivateurs, Ch. XVIII, Édition des Universités de Montpellier II et Paris VII
\msp{page 333}
\[
x \ = \ \varinjlim_{i\in I} x_i \; , \qquad\quad \text{les } x_i \text{ dans } \underline{S}(\pi),\ I \text{ grande devant } \pi.
\]

Alors le foncteur composé $i \longmapsto x_i$,
\[
I \ \longrightarrow \ \underline{S}(\pi)/x \ \hookrightarrow \ \mathcal{M}_\pi/x
\]
est cofinal (cf. 4.4.2, page 79), donc $\underline{S}(\pi)/x \longrightarrow \mathcal{M}_\pi/x$ satisfait la condition F 1 (puisque le composé le satisfait), et comme $\underline{S}(\pi)/x \longrightarrow \mathcal{M}_\pi/x$ est pleinement fidèle, cela implique déjà que ce foncteur est cofinal. Donc par $(**)$ on en conclut $x \in \Ob\underline{S}$, q.e.d.

Supposons maintenant que $S$ satisfasse c$'$ (au lieu de b), et prouvons encore l'équivalence $(*)$. Pour ceci, écrivons
\[
\square \qquad x \ = \ \varinjlim_{i\in I} x_i \qquad\quad
\left\{
\begin{minipage}[c]{0.5\textwidth}
les $x_i \in \Ob\mathcal{M}_\pi$, $I$ \textit{ensemble ordonné}, $(x_i)_{i\in I}$ est $\pi_0$-$\pi$-adaptée (\footnotemark[142]).
\end{minipage}
\right.
\]
\footnotetext[142]{\textbf{N.B.} Ça existe, grâce à 4.13.14.}

ce qui est possible puisque $\pi$ est utile pour $\mathcal{M}$ $\Bigg)$. Définissons $I_{\underline{S}}$ comme dans l'énoncé de c), donc comme image inverse de $\underline{S}(\pi)/x \hookrightarrow \mathcal{M}_\pi/x$ par $I \longrightarrow \mathcal{M}_\pi$. Par hypothèse d$'$) \texttt{[plutôt c$'$)]}, on a
\[
x \in \Ob\underline{S} \qquad \Longleftrightarrow \qquad I_{\underline{S}} \text{ est cofinal dans } I \; . \tag*{$(**')$}
\]

En utilisant $(**')$, prouvons encore $(*)$ (page 331). En fait, on utilisera seulement qu'on a
\msp{page 334}
l'isomorphisme $\square$, avec $I$ grande devant $\pi$, et qu'on a l'équivalence $(**')$, sans avoir à utiliser que $I$ est ordonnée. Donc la démonstration vaut aussi dans le cas c).

Prouvons $\Longrightarrow$ dans $(*)$, i.e. supposons $x \in \Ob\underline{S}$, prouvons que $x$ est dans l'image essentielle de $\Ind_\pi(\underline{S}(\pi))$. En vertu de $(**')$, $I_{\underline{S}}$ est cofinale dans $I$, et comme $I_{\underline{S}} \longrightarrow I$ est une inclusion pleinement fidèle, il s'ensuit que $I_{\underline{S}}$ est également grande devant $\pi$, d'autre part on aura par cofinalité
\[
x \ = \ \varinjlim_{i\in I_{\underline{S}}} x_i \; ,
\]
où cette fois les $x_i$ sont dans $\underline{S}(\pi) = \underline{S}\cap\mathcal{M}_\pi$, d'où résulte que $x$ est dans l'image essentielle de $\Ind_\pi(\underline{S}(\pi))$.

Prouvons $\Longleftarrow$ dans $(*)$, i.e. supposons que l'on ait
\[
x \ = \ \varinjlim_{j\in J} x_j \; , \qquad\quad \text{les } x_j \text{ dans } \underline{S}(\pi),\ J \text{ grande devant } \pi \; .
\]

\printedpage{127}
% head: A. Grothendieck : Dérivateurs, Ch. XVIII, Édition des Universités de Montpellier II et Paris VII

\msp{page 335}

\dgrblock{dgrb1083}{A divisible group recognized through countable subgroups}
\begin{dcomment}{examples}{p334}{A divisible group recognized through countable subgroups}
The criterion is easiest to see when the small pieces are actual subobjects. Consider $\mathcal M=\mathrm{Ab}$ and $\pi=\aleph_0$, so that the modern regular cardinal is $\lambda=\aleph_1$. The objects of $\mathcal M_\pi$ are the countable abelian groups. A countable group has a presentation with countably many generators and relations, and a map from it into an $\aleph_1$-filtered colimit factors at one stage, with all relations and all equalities satisfied at one further stage. Conversely every group is the $\aleph_1$-filtered union of its countable subgroups; presentability then makes its identity factor through one countable subgroup, forcing the group itself to be countable.

Let $\mathcal D\subset\mathrm{Ab}$ be the full subcategory of divisible groups. For an abelian group $X$, put
\[
 I_X=\{A\leq X:\lvert A\rvert\leq\aleph_0\},
\]
ordered by inclusion. This poset is $\aleph_1$-filtered: a countable family has a countable subgroup containing its union. It has colimits of countable chains, given by unions, and the inclusion into $\mathrm{Ab}$ preserves them. Its colimit is $X$.

If $X$ is divisible, every countable subgroup $A_0$ is contained in a countable divisible subgroup. Having chosen $A_n$, choose for each $a\in A_n$ and each positive integer $m$ one $m$th root of $a$ in $X$, and let $A_{n+1}$ be the subgroup generated by $A_n$ and these roots. There are countably many choices, so $A_{n+1}$ remains countable. The union $D=\bigcup_{n<\omega}A_n$ is divisible: an element belongs to some $A_n$ and each required root belongs to $A_{n+1}$. Thus the countable divisible subgroups form a cofinal subposet of $I_X$.

Conversely, if that subposet is cofinal, any $x\in X$ lies in a countable divisible subgroup, by applying cofinality to $\langle x\rangle$. Every equation $my=x$ with $m\geq1$ then has a solution in $X$. Hence
\[
 X\text{ divisible}\quad\Longleftrightarrow\quad
 \{A\in I_X:A\text{ divisible}\}\text{ cofinal in }I_X.
\]
Divisibility is preserved by filtered colimits: represent an element at a stage, choose a root there, and take its image. The countable divisible groups form a set of isomorphism classes and are $\aleph_1$-presentable in $\mathcal D$, since the inclusion preserves these colimits. The preceding cofinal presentation therefore proves that $\mathcal D$ is $\aleph_1$-accessible and accessibly embedded, with the indicated presentable objects. This establishes all the ingredients of the example, without invoking the classification of divisible groups.

The adaptedness condition cannot be discarded even here. The sequence
\[
 \mathbb Z\xrightarrow{\;2\;}\mathbb Z
 \xrightarrow{\;3\;}\mathbb Z
 \xrightarrow{\;4\;}\cdots
\]
has colimit $\mathbb Q$, using the maps $m\mapsto m/(n+1)!$ at stage $n$. None of its stages is divisible. Its indexing poset $\omega$ is filtered but is not $\aleph_1$-filtered, and does not contain the colimit of its own unbounded countable chain. The countable-subgroup presentation has exactly the additional closure which permits the membership criterion to work.
\end{dcomment}


En invoquant à nouveau 4.4.2, on trouve que $J \longrightarrow \mathcal{M}_\pi/x$ est cofinal, et comme ce foncteur se factorise par le foncteur d'inclusion pleinement fidèle $\underline{S}(\pi)/x \hookrightarrow \mathcal{M}_\pi/x$, il s'ensuit comme tantôt que ce dernier est cofinal. Nous allons admettre pour un moment que ceci implique que l'image inverse $I_{\underline{S}}$ de cette sous-catégorie strictement pleine, par le foncteur également cofinal $I \longrightarrow \mathcal{M}_\pi/x$, est cofinal dans $I$. Ceci étant, l'équivalence $(**')$ (p. 333) implique $x \in \Ob\underline{S}$, ce qu'on voulait.

Donc il reste à prouver que si
\[
J' \ = \ \underline{S}(\pi)/x \hookrightarrow \mathcal{M}_\pi/x \ = \ J
\]
est cofinal, alors $I' = I_{\underline{S}} \longrightarrow I$ l'est aussi. Il semble qu'il faille bel et bien à nouveau invoquer les conditions de stabilité $\mathrm{L}(\pi_0,\pi)$ faites dans d$'$) \texttt{[plutôt c$'$)]} ce qui n'avait pas été nécessaire pour le cas c) et
\msp{page 336}
le lemme-clef (un peu) vache 4.13.17, où on prendra (comme dans la démonstration écrite
avant l'énoncé de cet lemme) % [sic] cet lemme -> ce lemme, demonstrative adjective "cet" before a consonant
$\alpha = \alpha_\pi \overset{\text{déf}}{=}$ plus petit ordinal $\alpha$ tel que $\card\alpha = \pi$. Donc $I_\alpha$ est de cardinal $\pi$, et grande devant tout cardinal $\pi' < \pi$, et en particulier devant $\pi'_0$. L'hypothèse $2^\circ$) de 4.13.17 est vérifiée pour $I$, par définition d'un système inductif $\pi_0$-$\pi$-adapté, et pour $J'$, grâce à l'hypothèse $\mathrm{L}(\pi_0,\pi)$ sur $\underline{S}(\pi)$. De même, l'hypothèse $3^\circ$) ($I$ grande devant tout cardinal $< \pi$) est satisfaite par définition des systèmes $\pi_0$-$\pi$-adaptés -- $I$ est même grande devant $\pi$. L'hypothèse $1^\circ$) a déjà été vue.

On a donc prouvé (par voie circulaire)
\[
\begin{tikzcd}[row sep=small]
& & & \mathrm{c}' & \\
\mathrm{a} \arrow[r, Rightarrow] & \Bigg(\mathrm{a}' \arrow[r, Leftrightarrow] & \mathrm{c} \arrow[ur, Leftrightarrow] \arrow[dr, Rightarrow] & & \Bigg)\, , \\ % Verified against PDF leaf 9: the oblique equivalence has both arrowheads.
& & & \mathrm{b} &
\end{tikzcd}
\]

\texttt{[et aussi b $\Longrightarrow$ c,]} en ayant à invoquer le lemme-clef vache pour les \textit{deux} implications
\[
\mathrm{a}' \ \Longrightarrow \ \mathrm{c} \qquad\quad \text{et} \qquad\quad \mathrm{c}' \ \Longrightarrow \ \mathrm{a}'
\]
(donc pour deux implications \textit{en sens inverses}, chose étrange \dots).

\msp{page 337}

Il reste seulement encore à prouver l'implication
\[
\mathrm{a}' \ \Longrightarrow \ \mathrm{a} \; ,
\]
ce qui est essentiellement sorital.

Cela achève la démonstration du critère d'accessibilité 4.13.15 d'une sous-catégorie pleine $\underline{S}$ de $\mathcal{M}$.

On va donner une variante un peu décantée de 4.13.15, qui (avec des hypothèses un peu plus faibles) capte l'essentiel de sa substance. L'énoncé 4.13.15 en est une conséquence quasi immédiate :

\printedpage{128}
% head: A. Grothendieck : Dérivateurs, Ch. XVIII, Édition des Universités de Montpellier II et Paris VII

\msp{page 338}

\textbf{Corollaire 4.13.18} (de 4.13.15). \textit{Soient $\mathcal{M}$ une catégorie accessible, $\pi$ un cardinal utile pour $\mathcal{M}$, $\alpha_0$ un \textup{ordinal limite} tel que l'on ait $\card\alpha_0 \leqslant \pi$. Soit $I_{\alpha_0} = \{\alpha \mid \alpha < \alpha_0\}$, et supposons que $\mathcal{M}$ soit stable par $\varinjlim$ de type $I_{\alpha_0}$.}

\textit{Soit alors $\underline{S}$ une sous-catégorie strictement pleine de $\mathcal{M}$, $\underline{S}(\pi) = \underline{S}\cap\mathcal{M}_\pi$. Désignons par $L(\alpha_0)$ la condition : $\underline{S}(\pi)$ est stable dans $\mathcal{M}_\pi$ (ou ce qui revient au même, dans $\mathcal{M}$) par $\varinjlim$ de type $I_{\alpha_0}$. Alors les conditions suivantes sont équivalentes.}

\begin{enumerate}[label={}]
\item[\textit{a)}] \textit{$\underline{S}$ est $\pi$-accessible dans $\mathcal{M}$, et est stable dans $\mathcal{M}$ par $\varinjlim$ de type $I_{\alpha_0}$.}
\item[\textit{a$'$)}] \textit{$\underline{S}$ est $\pi$-accessible et $\underline{S}(\pi)$ satisfait $L(\alpha_0)$.}
\item[\textit{b)}] \textit{$\underline{S}(\pi)$ satisfait $L(\alpha_0)$, et pour tout $x \in \Ob\mathcal{M}$, on a ceci : $x \in \Ob\underline{S}$ si et seulement si la sous-catégorie strictement pleine $\underline{S}(\pi)/x$ de $\mathcal{M}(\pi)/x$ est cofinale dans $\mathcal{M}(\pi)/x$} \texttt{[lire $\mathcal{M}_\pi/x$]}.
\item[\textit{c)}] \textit{$\underline{S}(\pi)$ satisfait $L(\alpha_0)$, de plus pour tout $x \in \Ob\mathcal{M}$, et toute représentation}
\[
x \ \simeq \ \varinjlim_{i\in I} x_i
\]
\msp{page 339}
\textit{qui est ``$\alpha_0$-$\pi$-adaptée'', i.e. telle que}
\begin{enumerate}[label={}]
\item[\textit{$1^\circ$)}] \textit{$I$ équivalente à une petite catégorie, et grande devant $\pi$,}
\item[\textit{$2^\circ$)}] \textit{les $\varinjlim$ de type $I_{\alpha_0}$ existent dans $I$,}
\item[\textit{$3^\circ$)}] \textit{le foncteur $I \longrightarrow \mathcal{M}$, $i \longmapsto x_i$, commute aux $\varinjlim$ de type $I_{\alpha_0}$, et enfin}
\item[\textit{$4^\circ$)}] \textit{les $x_i \in \Ob\mathcal{M}_\pi$}
\end{enumerate}
\textit{(comparer 4.13.11), on a ceci (où on désigne par $I_{\underline{S}}$ la sous-catégorie strictement pleine de $I$ formée des $i$ tels que $x_i \in \Ob\underline{S}$) : on a $x \in \Ob\underline{S}$ si et seulement si $I_{\underline{S}}$ est cofinale dans $I$.}
\item[\textit{c$'$)}] \textit{(\textup{Si} $I_{\alpha_0}$ est grande devant $\pi_0$, et $\pi^{\pi_0} = \pi$, pour pouvoir invoquer 4.13.14, cf. p. 335.) Comme c), mais en supposant $I$ ordonnée.}
\end{enumerate}
(\footnote[143]{\textbf{N.B.} Si on ne suppose par $I_{\alpha_0}$ grande devant $\pi_0$ et % [sic] suppose par -> suppose pas, negation "ne ... pas"
$\pi^{\pi_0} = \pi$, alors toutes les conditions de a) à c) (a, a$'$, b, c) sont équivalentes et \textit{impliquent} c, mais ne lui sont peut-être pas équivalentes.}).

\dgrblock{dgrb1084}{The canonical presentation and an arbitrary adapted one}
\begin{dcomment}{exposition}{p338}{The canonical presentation and an arbitrary adapted one}
Grothendieck's criterion compares two ways of saying that the small objects of $\underline S$ suffice to build $x$. In the canonical presentation, this means that $\underline S(\pi)/x$ is final in $\mathcal M_\pi/x$. In a chosen presentation it means that the indices already lying in $\underline S$ are final. The first description is intrinsic; the second can be used in calculations.

For a useful $\pi$, the canonical diagram $\mathcal M_\pi/x\to\mathcal M$ is $\lambda$-filtered and has colimit $x$, where $\lambda=\pi^+$. If $\underline S(\pi)/x$ is final, it supplies a $\lambda$-filtered presentation of $x$ using objects of $\underline S(\pi)$, so $x$ lies in the essential image of $\operatorname{Ind}_\lambda(\underline S(\pi))$. Conversely, a presentation $x=\varinjlim_i x_i$ by such objects gives a final functor $I\to\mathcal M_\pi/x$. For example, an arrow $a\to x$ with $a$ presentable factors at a stage, and equality of two such factorizations holds at a later stage; these are the two required comma-category conditions. Its factorization through the full subcategory $\underline S(\pi)/x$ makes that subcategory final.

This equivalence alone does not permit pulling the subcategory back along every final presentation. That was precisely the failure in the even--odd example. Conditions on $\Phi$, and on the preservation of its colimits by the indexing diagram, supply the hypotheses of 4.13.17. They transfer finality to
\[
 I_{\underline S}
 =\{i\in I:x_i\in\underline S\}.
\]
The converse direction is then immediate from the closure of $\underline S$ under the filtered colimits in question.

Here being $\pi$-accessible \emph{in $\mathcal M$} carries more information than abstract accessibility of $\underline S$. It identifies $\underline S$ with the full image of the Ind-completion of $\underline S\cap\mathcal M_\pi$ and ensures that the inclusion preserves the chosen filtered colimits. In particular, its small objects are measured in the ambient category. This distinction is made explicit in 4.13.26. It prevents one from treating any abstractly accessible full subcategory as if it automatically had the canonical presentations needed here.

The optional restriction to ordered indexing categories introduces a separate existence problem. The exponentiation hypotheses in c$'$ are used to construct sufficiently closed ordered presentations of controlled cardinality. They are not part of the elementary argument with the canonical slice. Thus the text correctly distinguishes the unconditional equivalence of a), a$'$), b), c) from the stronger assertion involving c$'$; the missing prime in footnote 143 is a typographical slip in the edition.
\end{dcomment}


\textsc{Démonstration}. Elle se fait suivant le même diagramme logique que celle de 4.13.15,
\[
\begin{tikzcd}[row sep=small]
& & & \mathrm{c}' \arrow[dr, Rightarrow, "\textbf{!}"] & \nearrow\ \begin{tabular}[b]{@{}r@{}}\footnotesize ici, il faut supposer\\ \footnotesize $I_{\alpha_0}$ grande devant $\pi_0$\\ \footnotesize et $\pi^{\pi_0} = \pi$.\end{tabular} & \\
\mathrm{a} \arrow[r, Rightarrow] & \mathrm{a}' \arrow[r, Rightarrow, "\textbf{!}"] & \mathrm{c} \arrow[ur, Rightarrow] \arrow[dr, Rightarrow] & & \mathrm{a}' \arrow[r, Rightarrow] & \mathrm{a} \; , \\
& & & \mathrm{b} \arrow[ur, Rightarrow] & &
\end{tikzcd}
\]
où on a marqué du signe\,! les deux implications qui sont ``vaches'', en ce sens qu'elles font appel au lemme-clef 4.13.17 (pages 329-330), qui se prouve par une construction transfinie un peu vissée. La démonstration des sept implications du diagramme est essentiellement
celle donné dans les pages précédentes 311-337, % [sic] celle donné -> celle donnée, agreement with feminine "celle"

% ---- from der-ch18_p131-135.tex
\setcounter{footnote}{143}

\printedpage{129}
% head: A. Grothendieck : Dérivateurs, Ch. XVIII, Édition des Universités de Montpellier II et Paris VII
% foot: Version 17 décembre 2025

\msp{page 340}
à l'exception de a$'$ $\Longrightarrow$ a, qui est un peu plus simple dans le cas actuel. On va la dégager sous forme de lemme un peu plus général. (\textbf{N.B.} La démonstration de a$'$ $\Longrightarrow$ a dans 4.13.15 avait été omise comme “soritale”\,!, cf. page 337\,\dots)

\textbf{Lemme 4.13.19.} \textit{Soient $\mathcal{M}$ une catégorie accessible dans $\mathcal{M}$, $\pi$ un cardinal utile pour $\mathcal{M}$ et tel que $\underline{S}$ soit $\pi$-accessible, i.e. $\underline{S}$ est l'image essentielle de $\Ind_\pi(\underline{S}(\pi))$ par le foncteur pleinement fidèle canonique $\Ind_\pi(\underline{S}(\pi))\longrightarrow\mathcal{M}$, où $\underline{S}(\pi)=\underline{S}\cap\mathcal{M}_\pi$. On suppose qu'il existe un ordinal limite $\alpha$, tel que $\card\alpha\leqslant\pi$ et tel que $\mathcal{M}_\pi$ et $\underline{S}(\pi)$ $(=\underline{S}_\pi)$ soient stables par $\varinjlim$ de type $I_\alpha$ soient stables par $\varinjlim$ de type $I_\alpha$, et que le foncteur d'inclusion $\underline{S}(\pi)\longrightarrow\mathcal{M}_\pi$ y commute. Soit $\Phi$ un ensemble de catégories telles que $\forall\ I\in\Phi$, on ait $\cardred I\leqslant\pi$, et supposons que $\mathcal{M}_\pi$ soit stable par $\varinjlim$ de type $\Phi$, et que $\underline{S}(\pi)$ soit stable dans $\mathcal{M}_\pi$ par lesdites limites (ou ce qui revient au même, qu'il soit stable en lui-même par ces $\varinjlim$\,, et que le foncteur d'inclusion y commute). Alors $\mathcal{M}$ et $\underline{S}$ sont stables par $\varinjlim$ de type $\Phi$, et le foncteur d'inclusion $\underline{S}\longrightarrow\mathcal{M}$ y commute, i.e. $\mathcal{M}$ est stable par ces limites, et $\underline{S}$ stable dans $\mathcal{M}$ par ces mêmes limites.} % [sic] « soient stables par lim de type I_alpha » printed twice -> once, repetition

\msp{page 341}
\textsc{Démonstration}. On applique 4.11.1 (p. 267-268) séparément à $\mathcal{M}\simeq\Ind_\pi\mathcal{M}_\pi$ et à $\underline{S}\simeq\Ind_\pi\underline{S}(\pi)$, pour conclure que $\mathcal{M}$ et $\underline{S}$ sont stables séparément par $\varinjlim$ de type $\Phi$. Il reste à voir que le foncteur $\underline{S}\longrightarrow\mathcal{M}$ y commute. Cela se voit comme dans la démonstration de 4.11.1, on obtient plus généralement ceci :

\textbf{Lemme 4.13.20} (cor. à 4.11.1). \textit{Soient $f_0 : C\longrightarrow C'$ un foncteur entre catégories équivalentes à des petites catégories, $\pi$ un cardinal infini, $f:\mathcal{M}\stackrel{\textup{déf}}{=}\Ind_\pi C\longrightarrow\mathcal{M}'\stackrel{\textup{déf}}{=}\Ind_\pi C'$ l'extension canonique $\Ind_\pi(f_0)$ de $f_0$. Soient $\Phi$ et $\alpha$ comme dans 4.11.1, et supposons que $I_\alpha\in\Phi$ et que $f_0$ commute aux $\varinjlim$ de type $\Phi$. Alors $f$ commute à ces mêmes limites inductives.} (\footnote{Il vaut mieux prendre des foncteurs $f_0 : C\longrightarrow\mathcal{N}$, et $\mathcal{N}$ [?] que $f_0$ commute aux $\varinjlim$ de type $\Phi$, donc il en est de même de son extension $\Ind_\pi(C,\mathcal{N})\longrightarrow\mathcal{N}$. (\textbf{N.B.} On suppose que $\mathcal{N}$ satisfait $\mathrm{L}_\pi$.)}).

\textbf{Corollaire 4.13.21.} \textit{Sous les conditions de 4.13.20, supposons que $\Phi$ contienne les catégories de cardinal $\leqslant\pi$ (resp. les catégories filtrantes de cardinal $\leqslant\pi$, resp. les catégories filtrantes grandes devant $\pi_0$ de cardinal $\leqslant\pi$, resp. les catégories discrètes de cardinal $\leqslant\pi$). Alors $f$ commute aux petites $\varinjlim$ (resp. aux petites $\varinjlim$ filtrantes, resp. aux petites $\varinjlim$ filtrantes grandes devant $\pi_0$, resp. aux petites sommes). (Réciproque claire et “laissée au lecteur”.)}

Se voit comme 4.4.14 et 4.4.14.1 (p. 97, 98).

\msp{page 342}
Voici enfin, il me semble, la forme définitive (\,?) de 4.13.15, 4.13.18.

\dgrblock{dgrb1085}{Why a colimit calculation on small objects extends}
\begin{dcomment}{exposition}{p341}{Why a colimit calculation on small objects extends}
The passage from a$'$ to a) requires more than saying that large objects are filtered colimits of small ones. A diagram of large objects also has arrows and relations; the small presentations must realize them simultaneously. The earlier representation theorem, used in 4.11.1, supplies this compatibility under the limit-ordinal hypothesis.

Let $C$ and $C'$ be the essentially small categories of 4.13.20 and put $\lambda=\pi^+$. For a diagram $D:J\to\operatorname{Ind}_\lambda(C)$ of one of the allowed shapes, its simultaneous presentation is a functor
\[
                       X:J\times A\longrightarrow C
\]
with $A$ $\lambda$-filtered and $D(j)\simeq\varinjlim_{a\in A}X(j,a)$ naturally in $j$. The assumptions ensure that the $J$-colimits of $X(-,a)$ exist in $C$ and represent the corresponding colimits in its Ind-completion. Consequently
\[
\varinjlim_{j\in J}D(j)
 \simeq
\varinjlim_{a\in A}\left(\varinjlim_{j\in J}X(j,a)\right).
\]
This interchange uses only the universal property of colimits: both sides represent the set of cocones from the two-variable diagram to a prescribed target. It is not an interchange of a limit and a colimit.

If $f=\operatorname{Ind}_\lambda(f_0)$ and $f_0$ preserves these $J$-colimits, then
\[
\begin{aligned}
f\!\left(\varinjlim_jD(j)\right)
 &\simeq\varinjlim_a f_0\!\left(\varinjlim_jX(j,a)\right)\\
 &\simeq\varinjlim_a\varinjlim_jf_0X(j,a)\\
 &\simeq\varinjlim_j f(D(j)).
\end{aligned}
\]
All comparisons are natural, since they are induced by the same cocones. Apply this to the inclusion $\underline S(\pi)\hookrightarrow\mathcal M_\pi$ to obtain 4.13.19, including preservation by $\underline S\hookrightarrow\mathcal M$. This fills in the proof that page 337 calls sorital, while retaining the nontrivial simultaneous-presentation input rather than silently assuming it.

For the arbitrary-small-colimit assertion of 4.13.21, one can also see the size reduction directly. A small category is the directed union of its subcategories generated by at most $\pi$ arrows. The poset of these subcategories is $\lambda$-filtered; each has at most $\pi$ arrows because $\pi$ is infinite. The colimit of a diagram over the whole category is the colimit, over this poset, of its restricted colimits. Preservation of the bounded colimits and of $\lambda$-filtered colimits therefore gives preservation of all small colimits. The versions with filtered shapes use filtered envelopes of the small subcategories; the later construction in 4.13.37--39 makes that step explicit.

For the variant restricted to categories large before $\pi_0$, that bounded-envelope argument is verified below with the sufficient additional hypothesis $\pi^{\pi_0}=\pi$. The bare formulation of 4.13.21 does not display this hypothesis. We have not established that its one ordinal-colimit assumption suffices to remove it, so we do not certify that particular fixed-rank inference in its full printed generality. This qualification concerns the third variant, not the explicit interchange calculation above.

The underlying extension principle already occurs in \citet[Expos\'e I, 8.7.2--8.7.3]{SGA4I1972}: a functor on $C$ extends uniquely, up to canonical isomorphism, to a functor on $\operatorname{Ind}(C)$ preserving filtered colimits. Here the additional work controls the cardinal and the extra colimit shapes. It belongs to the same mathematical setting as the accessible-category theory of \citet{MakkaiPare1989}, which predates these manuscript pages; it should not be described as a later discovery of that theory.
\end{dcomment}


\textbf{Théorème 4.13.22} (\footnote{cf. variante 4.13.36, page 358.}). \textit{Soient $\mathcal{M}$ une catégorie accessible, $\pi$ un cardinal infini, utile pour $\mathcal{M}$, et $\Phi$ un ensemble de catégories satisfaisant les deux conditions :}
\begin{itemize}[leftmargin=2.5em]
\item[\textit{1°)}] $J\in\Phi$ \textit{implique} $\cardred J\leqslant\pi$.
\item[\textit{2°)}] \textit{Il existe un} ordinal limite $\alpha_0$ \textit{avec} $I_{\alpha_0}\in\Phi$.
\end{itemize}
\textit{On suppose $\mathcal{M}$ stable par $\varinjlim$ de type $\Phi$. Soient $\underline{S}$ une sous-catégorie} strictement pleine \textit{de $\mathcal{M}$, et $\underline{S}(\pi)=\underline{S}\cap\mathcal{M}_\pi$. Désignons par $L(\varphi)$ la condition sur $\underline{S}(\pi)$ : $\underline{S}(\pi)$ est stable dans} % [sic] L(\varphi) -> L(\Phi), the condition is later written L(\Phi)
\printedpage{130}
% head: A. Grothendieck : Dérivateurs, Ch. XVIII, Édition des Universités de Montpellier II et Paris VII
% foot: Version 17 décembre 2025
\textit{$\mathcal{M}$ (où ce qui revient au même, dans $\mathcal{M}_\pi$) par $\varinjlim$ de type $\Phi$. Ceci posé, les conditions suivantes sur $\underline{S}$ sont équivalentes :} % [sic] où ce qui revient -> ou ce qui revient, conjunction "ou"
\begin{itemize}[leftmargin=2.5em]
\item[\cercle{a}] \textit{$\underline{S}$ est $\pi$-accessible dans $\mathcal{M}$, et stable dans $\mathcal{M}$ par $\varinjlim$ de type $\Phi$.}
\item[a$'$)] \textit{$\underline{S}$ est $\pi$-accessible dans $\mathcal{M}$, et $\underline{S}(\pi)$ satisfait à $L(\Phi)$.}
\end{itemize}
\msp{page 343}
\begin{itemize}[leftmargin=2.5em]
\item[b)] \textit{$\underline{S}(\pi)$ satisfait à $L(\Phi)$, et pour tout $x\in\Ob\mathcal{M}$, on a ceci : $x\in\Ob\underline{S}$ si et seulement si la sous-catégorie strictement pleine $\underline{S}(\pi)/x$ de $\mathcal{M}(\pi)/x$ est cofinale dans $\mathcal{M}(\pi)/x$} \texttt{[lire} $\mathcal{M}_\pi/x$\texttt{]}.
\item[\cercle{c}] \textit{$\underline{S}(\pi)$ satisfait à $L(\Phi)$, et pour tout $x\in\Ob\mathcal{M}$ et toute représentation $x\simeq\varinjlim\limits_{i\in I}x_i$ qui est $\Phi$-$\pi$-adaptée (cf. définition plus bas), on a ceci : $x\in\Ob\underline{S}$ si et seulement si $I_{\underline{S}}$ est cofinale dans $I$, où $I_{\underline{S}}$ est la sous-catégorie strictement pleine de $I$ formée des $i\in\Ob I$ tels que $x_i\in\Ob\underline{S}$ (ou encore, $x_i\in\Ob\underline{S}(\pi)$).} \textbf{N.B.} \textit{On dit qu'un système inductif $(x_i)_{i\in I}$ dans $\mathcal{M}$ est $\Phi$-$\pi$-adapté si}
\begin{itemize}[leftmargin=2em]
\item[\textit{1)}] \textit{$I$ équivalente à une petite catégorie, et $I$ grande devant $\pi$,}
\item[\textit{2)}] \textit{les $\varinjlim$ de type $\Phi$ existent dans $I$,}
\item[\textit{3)}] \textit{le foncteur $i\longmapsto x_i$ de $I$ dans $\mathcal{M}$ commute aux $\varinjlim$ de type $\Phi$, et enfin}
\item[\textit{4)}] \textit{les $x_i\in\Ob\mathcal{M}_\pi$.}
\end{itemize}
\item[c$'$)] \textit{(Si $\pi^{\pi_0}=\pi$, où $\pi_0$ est un cardinal $\leqslant\pi$ utile pour $\mathcal{M}$, et si $J\in\Phi$ implique que $J$ est grande devant $\pi_0$.) Comme c), avec $I$} ordonnée.
\end{itemize}
(\footnote{\textbf{N.B.} Quand les conditions restrictives ci-contre \texttt{[parenthèse dans c$'$)]} ne sont pas satisfaites, alors a, a$'$, b, c sont encore équivalentes, et impliquent c$'$, mais ne sont peut-être pas équivalentes à c$'$).},\footnote{\textbf{N.B.} Il faudrait réénoncer les conditions annexes [?] (Si \dots), cf. p. 359-360.}).

\textbf{Définition 4.13.23} (\footnote{\textbf{N.B.} Il y aurait intérêt à remonter cette définition (\textit{sans} b), dès avant 4.13.15.}). Soient $\mathcal{M}$ une catégorie accessible, $\pi$ un cardinal, $\Phi$ un ensemble de catégories, satisfaisant les deux conditions :
\begin{itemize}[leftmargin=2.5em]
\item[a)] $J\in\Phi\Longrightarrow\cardred J\leqslant\pi$.
\item[b)] Il existe un ordinal limite $\alpha_0$, tel que $I_{\alpha_0}\in\Phi$ (\footnote{Si c'est pour poser une définition, on peut se passer de la condition b).}).
\end{itemize}
\msp{page 344}
Soit $S$ une partie de $\Ob\mathcal{M}$, $\underline{S}$ la sous-catégorie pleine qu'elle définit. On dit que $S$ est une \textit{partie $\Phi$-$\pi$-accessible de} $\Ob\mathcal{M}$, ou que $\underline{S}$ est une \textit{sous-catégorie $\Phi$-$\pi$-accessible dans} $\mathcal{M}$, si les conditions suivantes sont satisfaites.
\begin{itemize}[leftmargin=2.5em]
\item[1°)] $\pi$ est utile pour $\mathcal{M}$, $\mathcal{M}$ est stable par $\varinjlim$ de type $\Phi$ (\footnote{Condition préliminaire sur $\mathcal{M}$.}).
\item[2°)] $\underline{S}$ est une sous-catégorie $\pi$-accessible dans $\mathcal{M}$.
\item[3°)] $\underline{S}$ est stable dans $\mathcal{M}$ par $\varinjlim$ de type $\Phi$.
\end{itemize}
(\footnote{\textbf{N.B.} Quand $\Phi=\{I_{\alpha_0}\}$, on dit aussi que $\underline{S}$ est $\alpha_0$-$\pi$-accessible dans $\mathcal{M}$. Si $\Phi$ est formé des catégories de cardinal $\leqslant\pi$ et grandes devant $\pi_0$, on dit que $\underline{S}$ est $\pi_0$-$\pi$-accessible dans $\mathcal{M}$.}).

\textbf{N.B.} Quand la condition préliminaire 1°) \textit{sur} $\mathcal{M}$ est satisfaite, la validité des conditions 2°) 3°) équivaut donc à la condition que $\underline{S}$ soit \textit{strictement} pleine dans $\mathcal{M}$, et que de
\printedpage{131}
% head: A. Grothendieck : Dérivateurs, Ch. XVIII, Édition des Universités de Montpellier II et Paris VII
% foot: Version 17 décembre 2025
plus $\underline{S}$ satisfasse la condition a) du théorème 4.13.22. D'autre part, cette condition est équivalente à chacune des conditions a$'$, b, c de ce même théorème, et aussi à la dernière condition c$'$, sous réserve qu'on ait $\pi^{\pi_0}=\pi$ (pour un cardinal $\pi_0<\pi$ utile pour $\mathcal{M}$) et que $J\in\Phi$ implique que $J$ est grande devant $\pi_0$.

Conceptuellement, c'est la définition donnée dans 4.13.23, i.e. la condition \cercle{a} de 4.13.22, qui paraît la plus claire, ou sa variante immédiate a$'$. Mais du point de vue pratique, du travail avec les parties accessibles, c'est par propriété \cercle{c} qui paraît la plus commode. Le poids principal du théorème 4.13.22 est bien
\msp{page 345}
l'implication a $\Longrightarrow$ c (via a $\Longrightarrow$ a$'$ $\overset{!}{\Longrightarrow}$ c). Donc il faut surtout voir ce théorème comme donnant un critère remarquable d'appartenance à $\Ob\underline{S}$ d'un objet $x$ de $\mathcal{M}$ qui est présenté par un système inductif $\Phi$-$\pi$-adapté :
\[ x\ \in\ \Ob\underline{S}\qquad\quad\Longleftrightarrow\qquad\quad I_{\underline{S}}\ \text{cofinal}\ , \tag*{$(*)$} \]
et dans cette équivalence, c'est l'implication $\Longrightarrow$ qui est la partie non soritale, et remarquable.

\textbf{Proposition 4.13.24.} \textit{Soient $\mathcal{M}$, $\pi$, $\Phi$ comme dans 4.13.23} (\footnote{\textbf{N.B.} La condition $I_{\alpha_0}\in\Phi$ n'est pas utile.})\textit{, on suppose que la condition 1° de 4.13.23 est satisfaite ($\pi$ utile pour $\mathcal{M}$, $\mathcal{M}$ stable par $\varinjlim$ de type $\Phi$). Soit ${\Sigma'}^{\Phi}_{\pi}$ l'ensemble des sous-catégories $\underline{S}$ de $\mathcal{M}$ qui sont $\Phi$-$\pi$-accessibles dans $\mathcal{M}$, et $\Sigma^{\Phi}_{\pi}$ l'ensemble des sous-catégories $\underline{S}(\pi)$ de $\mathcal{M}_\pi$ qui sont strictement pleines, et stables dans $\mathcal{M}_\pi$ par $\varinjlim$ de type $\Phi$. (Ces sous-catégories sont forcément karoubiennes, car stables par $\varinjlim$ de type $I_{\alpha_0}$, donc aussi par facteurs directs.) Alors les bijections (4.13.4.1) (page 288) induisent des bijections inverses l'une à l'autre}
\[
\begin{tikzcd}
\Sigma^{\Phi}_{\pi} \arrow[r, shift left] & {\Sigma'}^{\Phi}_{\pi} \arrow[l, shift left]
\end{tikzcd}\ .
\]

\dgrblock{dgrb1086}{Retracts are essential in the small-subcategory correspondence}
\begin{dcomment}{development}{p345:Ces sous-catégories sont forcément karoubiennes}{Retracts are essential in the small-subcategory correspondence}
Grothendieck wants to replace a large accessible subcategory by its small part, so that operations on the former can be performed inside an essentially small category. Proposition 4.13.24 gives this correspondence when the small part is closed under the stated colimits and under retracts. The parenthetical observation supplies retract closure from an ordinal colimit. Its footnote, however, says that the ordinal hypothesis is unnecessary, without replacing it by a retract hypothesis.

\textbf{Counterexample to the footnote as stated.}\ours\ Take $\mathcal M=\mathrm{Set}$, $\pi=\aleph_0$ and $\Phi=\varnothing$. Let $C$ be the strictly full subcategory of $\mathcal M_\pi$ consisting of the two-element sets. It satisfies all the stated $\Phi$-closure requirements, vacuously. Nevertheless the filtered-colimit closure of $C$ contains a singleton, so taking its small part does not recover $C$.

For completeness, fix a two-element set $B$ and an idempotent constant map $e:B\to B$. Index a diagram by the ordinal $\omega_1$, put $B$ at each index, and put $e$ on every strict transition and the identity on identity transitions. This is a functor because $e^2=e$. The ordinal $\omega_1$ is $\aleph_1$-filtered: every countable set of its indices is bounded. Its colimit in sets is the image of $e$, a singleton. The singleton is also $\aleph_1$-presentable, hence belongs to the small part of the Ind-image, but not to $C$. This proves the failure of the asserted inverse correspondence under the footnote's weakened hypotheses.

The exact repair is to require $C$ to be closed under retracts, or to retain a nonzero limit-ordinal shape among the closure assumptions. To verify the latter option, take any idempotent $e:x\to x$ in $C$ and form the constant-object diagram along that ordinal, with $e$ on strict transitions. If its colimit is $q:x\to y$, the compatible maps $e:x\to x$ induce $s:y\to x$ with $sq=e$. Since $qe=q$ and the colimit maps jointly determine a morphism out of $y$, one has $qs=1_y$. Thus $y$ splits $e$, and closure under this colimit puts its splitting object in $C$.

Conversely, when $C\subset\mathcal M_\pi$ is full and retract closed, every $\lambda$-presentable object of its Ind-image is a retract of an object of $C$: factor its identity through one stage of a $\lambda$-filtered presentation. Hence the small part of that image is exactly $C$. Combined with 4.13.19, this gives the stated correspondence with the intended additional-colimit closure.

The general role of idempotent completion in Ind-constructions is standard. Our contribution here is the explicit diagnosis and counterexample for the particular footnote, checked against the edition and the surrounding propositions. We have not found a published discussion of this manuscript slip.
\end{dcomment}


\msp{page 346}
% marge : trait manuscrit (reproduit en image) en regard de la Proposition 4.13.25
\textbf{Proposition 4.13.25.} \textit{Soient $\mathcal{M}$, $\pi$, $\Phi$ comme dans 4.13.23, avec $I_{\alpha_\pi}\in\Phi$ ($\alpha_\pi$ étant le plus petit ordinal de cardinal $\Phi$} \texttt{[plutôt} $\pi$\texttt{]}\textit{). Alors l'ensemble ${\Sigma'}^{\Phi}_{\pi}$ des sous-catégories de $\mathcal{M}$ $\Phi$-$\pi$-accessibles dans $\mathcal{M}$ est stable par intersection de familles de cardinal $\leqslant\pi$} (\footnote{Si $\pi$ n'est pas cardinal limite. Sinon, il faut dire : familles de cardinal $<\pi$. Cf. 4.13.10.3, où je donne un énoncé rectifié légèrement plus général.}).

(\textbf{N.B.} La bijection $\Sigma^{\Phi}_{\pi}\xrightarrow{\ \sim\ }{\Sigma'}^{\Phi}_{\pi}$ de 4.13.24 respecte les relations d'ordre naturelles d'inclusion, donc transforme intersections en intersections.)

Se voit comme 4.13.7 (qui est contenu dans l'énoncé actuel), en utilisant le lemme 4.13.10.

\textbf{4.13.26.} Soient $\mathcal{M}$, $\pi$, $\Phi$ comme dans 4.13.23, et soit $\underline{S}$ une sous-catégorie \textit{strictement pleine} de $\mathcal{M}$, $f:\underline{S}\longrightarrow\mathcal{M}$ le foncteur d'inclusion. La condition que $\underline{S}$ soit $\Phi$-$\pi$-accessible dans $\mathcal{M}$ se décompose en :
\begin{itemize}[leftmargin=2.5em]
\item[1)] $\pi$ est utile pour $\mathcal{M}$ \textit{et} pour $\underline{S}$.
\item[2)] $\mathcal{M}$ et $\underline{S}$ sont stables par $\varinjlim$ de type $\Phi$.
\item[3)] $f$ commute auxdites limites.
\item[4)] $f$ est $\pi$-accessible.
\item[5)] $f(\underline{S}_\pi)\subset\mathcal{M}_\pi$.
\end{itemize}

\printedpage{132}
% head: A. Grothendieck : Dérivateurs, Ch. XVIII, Édition des Universités de Montpellier II et Paris VII
% foot: Version 17 décembre 2025
\msp{page 347}
Soit alors $\pi'>\pi$ un autre cardinal, quand peut-on assurer que si $\underline{S}$ est $\Phi$-$\pi$-accessible dans $\mathcal{M}$, il soit aussi $\Phi$-$\pi'$-accessible\,? Les conditions 2) 3) 4) sont visiblement stables par passage de $\pi$ à $\pi'$. Restent les conditions 1) et 5). On sait que si $\pi'^{\pi}=\pi'$, alors la condition 1) reste valable pour $\pi'$. Plus généralement, si $\pi_0$ est un cardinal $\leqslant\pi$ tel que $\pi_0$ soit utile pour $\mathcal{M}$ et pour $\underline{S}$, alors si $\pi'>\pi_0$ est tel que $\pi'^{\pi_0}=\pi'$, alors $\pi'$ est aussi utile pour $\mathcal{M}$ et pour $\underline{S}$. Si on admet l'HCG (hypothèse du continu généralisée), il en est ainsi pour tout $\pi'>\pi_0$, donc pour tout $\pi'>\pi$.

\dgrblock{dgrb1087}{Changing the rank without changing the condition}
\begin{dcomment}{exposition}{p346}{Changing the rank without changing the condition}
Grothendieck wants a calculus of conditions on objects which remains usable when several constructions have to be carried out together. An accessible subcategory comes with a size bound, whereas an intersection, an inverse image, or a category of diagrams will usually require a larger common bound. The point of 4.13.26--27 is to arrange that enlargement once, including its effect on the compact generators, before using the calculus.

There are two sizes in the notation. Here ``grande devant $\pi$'' means that diagrams with at most $\pi$ arrows have cocones; in the usual convention this is $\pi^+$-filteredness. Thus $\mathcal M_\pi$ is the subcategory of $\pi^+$-presentable objects, and a useful $\pi$ supplies an essentially small family of these objects generating $\mathcal M$ under $\pi^+$-filtered colimits. The five conditions on the inclusion $f:\underline S\hookrightarrow\mathcal M$ separate preservation of colimits from preservation of presentability. The former becomes easier as the filteredness bound increases. The latter is a genuine uniformity requirement: sufficiently small generators of the source must have sufficiently small images in the target. It is this requirement which accounts for the cardinal of $\underline S_{\pi_0}$ in 4.13.27.

The explicit enlargement $\pi'=2^c$ is particularly convenient. If $c$ bounds $\pi$, the generating data, and $\pi_0$, then
\[
 (\pi')^{\pi_0}=(2^c)^{\pi_0}
 =2^{c\cdot\pi_0}=2^c=\pi'.
\]
The argument therefore provides arbitrarily large usable bounds without a continuum hypothesis. In modern accounts this belongs to uniformization: accessible functors preserve presentable objects at suitable sufficiently large ranks, not at an arbitrarily prescribed larger rank. The distinction is visible in the rank comparison recalled in \citet[Section 2]{BrookeTaylorRosicky2017}.

The proposed GCH shortcut requires a correction. Even under GCH, $\lambda>\mu$ need not imply $\lambda^\mu=\lambda$. Take $\mu=\aleph_0$ and $\lambda=\aleph_\omega$. To see the obstruction directly, partition a set of cardinal $\lambda$ into countably many pieces $A_n$ of cardinal less than $\lambda$. If all functions $\mathbf N\to\lambda$ could be enumerated as $(f_\alpha)_{\alpha<\lambda}$, choose $g(n)$ outside $\{f_\alpha(n):\alpha\in A_n\}$. This is possible because $|A_n|<\lambda$. For every $\alpha$, the coordinate corresponding to the piece containing $\alpha$ gives $g(n)\ne f_\alpha(n)$, contradicting the enumeration. Hence $\lambda^{\aleph_0}>\lambda$.

A correct GCH version adds $\operatorname{cf}(\lambda)>\mu$. Every map $\mu\to\lambda$ then has bounded image. For an infinite $\nu<\lambda$, GCH gives $\nu^\mu\leq 2^{\max(\nu,\mu)}\leq\lambda$, so the union over all bounds still has cardinal at most $\lambda$; constant functions give the reverse inequality. This repairs the arithmetic statement. It does not require us to reject the category-theoretic conclusions obtained by choosing the larger bound freely: the power-of-two construction already supplies what their proofs need.
\end{dcomment}


Reste la condition 5), pour laquelle une réponse est donnée par 4.12.1 (p. 273). Si $\pi_0\leqslant\pi$ est choisi de façon que $f$ soit $\pi_0$-accessible (p.ex. $\pi_0=\pi$), et si
\[
\begin{aligned}
\pi'_1 &= \cardred\underline{S}_{\pi_0} \qquad\quad (\,\leqslant\cardred\underline{S}(\pi)\,)\\
\pi''_1 &= \sup(\pi'_1,\pi)\ ,
\end{aligned}
\]
on voit que si $\pi'\geqslant\pi''_1$ (i.e. $\pi'\geqslant\pi'_1,\pi$) est tel que $\pi'^{\pi_0}=\pi'$, alors on a bien
\[ f(\underline{S}_{\pi'})\ \subset\ \mathcal{M}_{\pi'}\ . \]

\msp{page 348}
Si donc de plus $\pi'>\pi\geqslant\pi_0$ (comme supposé au début), et si on suppose l'HCG, alors la condition $\pi'^{\pi_0}=\pi'$ est automatique, et il reste la seule condition $\pi'\geqslant\pi''_1$, i.e. en définitive $\pi'\geqslant\cardred\underline{S}_{\pi_0}$. Pour résumer :

\textbf{Proposition 4.13.27.} \textit{Sous les conditions de 4.13.23 (sans condition b) sur $\Phi$), soit $\underline{S}$ une sous-catégorie de $\mathcal{M}$ qui est $\Phi$-$\pi$-adaptée. Soit $\pi_0\leqslant\pi$ un cardinal qui est utile pour $\underline{S}$ et pour $\mathcal{M}$, et tel que $f:\underline{S}\longrightarrow\mathcal{M}$ soit $\pi_0$-accessible, i.e. tel que $\underline{S}$ soit stable dans $\mathcal{M}$ par $\varinjlim$ grandes devant $\pi_0$ (on peut p.ex. prendre $\pi_0=\pi$). Alors pour tout cardinal $\pi'>\pi$ tel que}
\begin{itemize}[leftmargin=2.5em]
\item[\textit{a)}] \textit{$\pi'^{\pi_0}=\pi'$ et}
\item[\textit{b)}] \textit{$\pi'\geqslant\cardred\underline{S}_{\pi_0}$,}
\end{itemize}
\textit{$\underline{S}$ est $\Phi$-$\pi'$-accessible dans $\mathcal{M}$. Donc si $\pi'>\pi$ est tel que $\pi'^{\pi}=\pi'$, $\pi'>\cardred\underline{S}_\pi$, alors $\underline{S}$ est $\Phi$-$\pi'$-accessible. (Il suffit p.ex. de prendre $\pi'=2^c$, avec $c\geqslant\sup(\pi,\cardred\underline{S}(\pi))$.) Si on admet l'HCG, il suffit donc qu'on ait $\pi'>\sup(\pi,\cardred\underline{S}(\pi))$.}

\msp{page 349}
\textbf{Proposition 4.13.28} (\footnote{À rejeter plus bas, après 4.13.31, qui est la première application importante de 4.13.22, alors que 4.13.28 se prouve sans avoir à utiliser 4.13.22 ou ses variantes.}). \textit{Soient $\mathcal{M}$, $\pi$, $\Phi$ comme dans 4.13.23. Supposons de plus que $J\in\Phi$ implique que $J$ est} filtrante. \textit{Alors l'ensemble $\tilde{\Sigma}^{\Phi}_{\pi}$ des sous-ensembles $\Phi$-$\pi$-accessibles de $\mathcal{M}$ est stable par réunions finies. Plus généralement, soit $\lambda$ un cardinal $<\pi$ tel que $\forall\ J\in\Phi$, $J$ soit grande devant $\lambda$. Alors l'ensemble $\tilde{\Sigma}^{\Phi}_{\pi}$ est stable par réunions $(S_r)_{r\in A}$, indexées par des ensembles $A$ de cardinal $\leqslant\lambda$.}

\textbf{Corollaire 4.13.29.} \textit{L'ensemble des parties accessibles de} $\Ob\mathcal{M}$ \textit{est stable par “petites” réunions -- i.e. si $(S_r)_{r\in A}$ est une petite famille de parties accessibles de} $\Ob\mathcal{M}$, \textit{alors $\bigcup_{r\in A}S_r$ est aussi accessible.}

La démonstration de 4.13.28 repose sur le lemme suivant (je n'explicite par la démonstration) : % [sic] par la démonstration -> pas la démonstration, sense of the sentence

\printedpage{133}
% head: A. Grothendieck : Dérivateurs, Ch. XVIII, Édition des Universités de Montpellier II et Paris VII
% foot: Version 17 décembre 2025
\msp{page 350}
\textbf{Lemme 4.13.30.} \textit{Soit $I$ une catégorie grande devant $\pi$, et $(J_r)_{r\in A}$ une famille des sous-catégories pleines telles que $\card A\leqslant\pi$, telle que $\bigcup_{r\in A}\Ob J_r$ soit cofinal dans $\Ord(I)$, i.e. pour tout $i\in\Ob I$ $\exists\ j\in\bigcup_{r\in A}\Ob J_r$, et une flèche $i\longrightarrow j$. Alors il existe $r\in A$ tel que $J_r$ soit cofinal dans $I$.}

\textbf{Théorème 4.13.31.} \textit{Soient $\mathcal{M}$, $\mathcal{N}$ des catégories accessibles, $f:\mathcal{M}\longrightarrow\mathcal{N}$ un foncteur accessible, $T\subset\Ob\mathcal{N}$ une partie de $\mathcal{N}$, $\underline{T}\subset\mathcal{N}$ la sous-catégorie strictement pleine correspondante de $\mathcal{N}$, $S=f^{-1}(T)$, $\underline{S}=f^{-1}(\underline{T})$ la sous-catégorie pleine de $\mathcal{M}$ associée. Si $T$ est accessible, alors $S$ est accessible\,; en d'autres termes, si $\underline{T}$ est accessible dans $\mathcal{N}$, $\underline{S}$ est accessible dans $\mathcal{M}$.}

Cela résulte de l'énoncé plus précis :

\textbf{Corollaire 4.13.32.} \textit{Soient $\pi_0$ et $\pi>\pi_0$ des cardinaux ayant les propriétés suivantes :}

\msp{page 351}
\begin{itemize}[leftmargin=2.5em]
\item[\textit{a)}] \textit{$\underline{T}$ est $\pi_0$-$\pi$-accessible dans $\mathcal{N}$ (4.13.23).}
\item[\textit{b)}] \textit{$\pi_0$ est utile pour $\mathcal{M}$ (il l'est pour $\mathcal{N}$, à cause de a)).}
\item[\textit{c)}] \textit{$f(\mathcal{M}_\pi)\subset\mathcal{N}_\pi$.}
\item[\textit{d)}] \textit{$\mathcal{M}$ satisfait à $L_{\pi_0}$, et $f$ est $\pi_0$-accessible.}
\end{itemize}
\textit{(L'existence d'un tel couple $(\pi_0,\pi)$ résulte aussitôt de 4.13.16 et de 4.12.1.) Alors $\underline{S}$ est $\pi_0$-$\pi$-accessible dans $\mathcal{N}$.}

\dgrblock{dgrb1088}{A calculus of object conditions}
\begin{dcomment}{exposition}{p350}{A calculus of object conditions}
Grothendieck's accessible parts are intended to behave like manageable predicates on the objects of a large category. The two operations now needed are disjunction and substitution: a small union combines alternative conditions, while an inverse image expresses a condition after applying a functor. These are not statements merely about the number of objects. They say that the resulting full subcategory again has compact generators and that its sufficiently filtered colimits are computed in the ambient category.

The omitted proof of Lemma 4.13.30 explains why disjunction works. Let $I$ be $\pi^+$-filtered and suppose that no $J_r$ is objectwise cofinal. Choose $i_r\in I$ for which no arrow $i_r\to j$ exists with $j\in J_r$. Since there are at most $\pi$ indices $r$, choose a common upper bound $i$ of the $i_r$. Cofinality of the union supplies an arrow $i\to j$ with $j\in J_s$ for some $s$. Its composite with $i_s\to i$ contradicts the choice of $i_s$. Consequently some $J_s$ is objectwise cofinal.

For a full subcategory of a filtered category, objectwise cofinality supplies the usual categorical finality. Any finite diagram in $J_s$ first has a cocone in $I$; an arrow from its vertex to an object of $J_s$ gives a cocone in $J_s$, because the subcategory is full. The same argument in each comma category $(i\downarrow J_s)$ supplies a nonempty filtered category. This proves finality and justifies using $J_s$ to compute colimits. Applied to the indices at which an adapted presentation belongs to a given member of a family, the lemma says that a sufficiently filtered colimit lying in their union can be computed wholly inside one member. Enlarging the common cardinal to bound a small family gives Corollary 4.13.29.

The inverse-image argument uses the stronger cofinal-presentation criterion established earlier. For $x\in f^{-1}(\underline T)$, present $x$ by the compact objects over it. The image presentation consists of compact objects of $\mathcal N$, once the common bound has been chosen to make $f$ preserve presentability. Since its colimit lies in $\underline T$, the criterion supplies cofinally many stages already lying in $\underline T$. Their preimages are the desired compact generators of $f^{-1}(\underline T)$. The argument proves essential surjectivity of
\[
 \Ind_\pi\bigl(f^{-1}(\underline T)\cap\mathcal M_\pi\bigr)
 \longrightarrow f^{-1}(\underline T).
\]
Full faithfulness follows from compactness in $\mathcal M$.

There is a small rank omission in the quantitative formulation 4.13.32. Its proof uses that $\pi$ is useful for $\mathcal M$, whereas condition b) literally says that $\pi_0$ is useful. The safe formulation includes the former condition, obtained, for example, by imposing $\pi^{\pi_0}=\pi$. This is harmless for the qualitative theorem, where the common rank is free to increase. The conclusion of the corollary also prints $\mathcal N$ where it must say $\mathcal M$; that is a source slip, not an error of transcription.

The present categorical formulation packages the same mechanism as a pseudopullback. Its objects are triples $(x,t,\theta)$ with $\theta:f(x)\xrightarrow{\sim}t$ and $t\in\underline T$. Because $\underline T$ is replete and full, forgetting $t,\theta$ is an equivalence with the strict full inverse image. The accessible pseudopullback theorem therefore applies. This is one instance of the limit theorem of \citet[Theorem 5.1.6]{MakkaiPare1989}; its precise relation to products, inserters and equifiers is discussed in \citet[Introduction and Section 5]{Positselski2024Limits}. The repleteness assumption is what permits passage from an isomorphism to literal membership, and should remain visible in either proof.
\end{dcomment}


Comme $f$ commute aux $\varinjlim$ grandes devant $\pi_0$ en vertu de d), et que $\underline{T}$ est stable par ces limites par a) (cf. 4.13.23, 3°), il en est donc de même de $\underline{S}$, i.e. le 3° de la définition 4.13.23 est satisfait. D'autre part, la condition 1° ($\pi$ utile pour $\mathcal{M}$, $\mathcal{M}$ stable par $\varinjlim$ grandes devant $\pi_0$ et de cardinal $\leqslant\pi$) est vérifiée par b) et d). Donc il reste à vérifier la condition 2°, i.e. que $\underline{S}$ est $\pi$-accessible. Posons
\[
\begin{aligned}
\underline{T}(\pi) &= \underline{T}\cap\mathcal{N}_\pi\\
\underline{S}(\pi) &= \underline{S}\cap\mathcal{M}_\pi\ \ \underbrace{\ \subset\ }_{\text{grâce à c)}}\ \ f^{-1}(\underline{T}(\pi))\ .
\end{aligned}
\]
Comme $\underline{S}$ est stable dans $\mathcal{M}$ par $\varinjlim$ grandes devant $\pi_0$ (lesquelles existent dans $\mathcal{M}$ par d), et a fortiori par $\varinjlim$ grandes devant $\pi$, on a un foncteur canonique pleinement fidèle
\[ \Ind_\pi\underline{S}(\pi)\ \longrightarrow\ \underline{S}\ . \]

\msp{page 352}
D'ailleurs, $\underline{T}$ étant strictement pleine dans $\mathcal{N}$, $\underline{S}$ est strictement pleine dans $\mathcal{M}$, et il reste à vérifier que le foncteur précédent est essentiellement surjectif. Soit donc $x\in\Ob\underline{S}$. On aura, grâce à b),
\[ x\ =\ \varinjlim_{I}x_i\ ,\qquad\quad x_i\in\mathcal{M}_\pi,\ I=\mathcal{M}_\pi/x\ . \]
D'ailleurs, le système inductif précédent est $\pi_0$-$\pi$-adapté, i.e. (p. 339) :
\begin{itemize}[leftmargin=2.5em]
\item[1)] $I$ équivalente à une petite catégorie, $I$ grande devant $\pi$ (résulte du fait que $\pi$ est utile pour $\mathcal{M}$, cf.\,\cercle{b}\,).
\end{itemize}

% ---- from der-ch18_p136-140.tex
\printedpage{134}
% head: A. Grothendieck : Dérivateurs, Ch. XVIII, Édition des Universités de Montpellier II et Paris VII
% foot: Version 17 décembre 2025
\begin{itemize}[leftmargin=2.5em]
\item[2)] Les limites inductives grandes devant $\pi_0$ est de cardinal $\leqslant \pi$ existent dans $I$ (résulte de $\mathrm{L}_{\pi_0}$ sur $\mathcal{M}$ (cf.\,\cercle{d})). % [sic] est de cardinal -> et de cardinal, coordination ("grandes devant π0 et de cardinal ≤ π")
\item[3)] Le foncteur $\xi : I \longrightarrow \mathcal{M}$ y commute.
\item[4)] Les $x_i$ sont dans $\mathcal{M}_\pi$.
\end{itemize}
Comme $f$ commute aux $\limind$ grandes devant $\pi_0$ (cf.\,\cercle{1}), a fortiori à celles grandes devant $\pi$, on a donc
\[ f(x) \;\simeq\; \varinjlim_{I} f(x_i) \ , \]
je dis que le système inductif $(f(x_i))_{i\in I}$ dans $\mathcal{N}$ est $\pi_0$-$\pi$-adapté. Les conditions 1), 2) sont \textit{ne varietur} (ne concernent que $I$), le foncteur $I \longrightarrow \mathcal{N}$ est le composé $I \xrightarrow{\xi} \mathcal{M} \xrightarrow{f} \mathcal{N}$, donc commute aux $\limind$ grandes devant $\pi_0$ et de cardinal $\leqslant \pi$ (\footnote[155]{cette commutation devrait avoir un nom, condition $\mathrm{L}(\pi_0,\pi)$ sur un foncteur.}), puisqu'il en est ainsi de $\xi$ et de $f$, enfin les $f(x_i)$ sont dans $\mathcal{N}_\pi$ puisque par \cercle{c} $f(\mathcal{M}_\pi) \subset \mathcal{N}_\pi$.

\msp{page 353}

Par ailleurs $f(x) \in \Ob \underline{T}$ (puisque $x \in \Ob \underline{S}$, où $\underline{S} \defeq f^{-1}(\underline{T})$). Donc on peut appliquer 4.13.22 a $\Longrightarrow$ c à la situation $(\mathcal{N}, \pi_0, \pi, \underline{S})$ et au système inductif $(f(x_i))_{i\in I}$, pour conclure que la sous-catégorie $I'$ de $I$ formée des $i$ tels que $f(x_i) \in \Ob \underline{S}$ \texttt{[plutôt Ob}~$\underline{T}$\texttt{]}, i.e. tels que $x_i \in \underline{S}$, donc $x_i \in \underline{S}(\pi)$, est \textit{cofinale} dans $I$. Donc on a
\[ x \;\simeq\; \varinjlim_{i\in I'} x_i \ , \]
avec cette fois les $x_i$ dans $\underline{S}(\pi)$, et $I'$ grande devant $\pi$ (car cofinale et pleine dans $I$, qui est grande devant $\pi$), ce qui prouve que $x$ est dans l'image essentielle de $\Ind_\pi(\underline{S}(\pi))$, q.e.d.

\textbf{Question 4.13.33.} On vient de voir que si $f : \mathcal{M} \longrightarrow \mathcal{N}$ est un foncteur accessible entre catégories accessibles, alors l'image inverse d'une partie accessible de $\Ob \mathcal{N}$ est une partie accessible de $\Ob \mathcal{M}$. Je conjecture que l'énoncé similaire, non pas pour les images directes, mais pour les \textit{images essentielles}, est
\msp{page 354}
vraie, i.e. que pour toute sous-catégorie $\underline{S}$ de $\mathcal{M}$ accessible dans $\mathcal{M}$, son image essentielle dans $\mathcal{N}$ est accessible dans $\mathcal{N}$ (\footnote[156]{Faux, même si on prend la $\natural$-clôture de l'image essentielle (“image $\natural$-essentielle”).}). Dans cette question, on peut supposer que $\underline{S} = \mathcal{M}$, et la question revient à celle-ci : tout foncteur accessible $f : \mathcal{M} \longrightarrow \mathcal{N}$ entre catégories accessibles se factorise-t-il en un foncteur (accessible) \textit{essentiellement surjectif}, et un foncteur accessible pleinement fidèle $\mathcal{M} \longrightarrow \underline{T} \longrightarrow \mathcal{N}$ (où $\underline{T}$ est elle-même accessible) (\footnote[157]{Il faut se borner tout au moins à essentiellement $\natural$-surjectif – mais même cela est trop optimiste.})~? Si oui, cela me paraîtrait peut-être le résultat le plus profond (peut-être le seul véritablement profond) en théorie des catégories. % [sic] vraie -> vrai, agreement with « l'énoncé similaire »

\dgrblock{dgrb1089}{Why forgetting a retraction destroys the image theorem}
\begin{dcomment}{examples}{p353}{Why forgetting a retraction destroys the image theorem}
The hope behind 4.13.33 is that an accessible construction might be followed by forgetting all its choices, while its essential image still retained a compact description. The footnotes reject that hope, including the attempted repair by closing under retracts. The following example makes both failures explicit. It is the split-monomorphism example of \citet[Example 3.5(1), author manuscript]{Rosicky2009Combinatorial}; the retract-closed formulation appears in \citet[Remark 3.3(2), accepted version]{BrookeTaylorRosicky2017}.

Let $R$ be the small category with two objects $a,b$, arrows $i:a\to b$ and $p:b\to a$, and relation $pi=1_a$. The category $\operatorname{Pos}^{R}$ consists of split inclusions of posets with chosen retractions and maps preserving both pieces of data. Restriction along the walking arrow gives
\[
 F:\operatorname{Pos}^{R}\longrightarrow
       \operatorname{Pos}^{[1]},\qquad (i,p)\longmapsto i.
\]
These are accessible categories and $F$ is accessible. One can check this without a general model theorem. A poset is the filtered union of its finite subposets. A finite collection of elements in an $R$-diagram generates a finite subdiagram, since $R$ has only finitely many morphisms. The inherited orders are finite, and all operations and their equations restrict. Such finite diagrams are finitely presentable: maps and the finitely many order relations factor, and are eventually identified, in a filtered colimit. The same argument applies to the arrow category. Filtered colimits are computed pointwise, so restriction preserves them.

The replete full image $\mathcal L$ consists of split monomorphisms of posets, with every commutative square allowed as a morphism. Fix any infinite regular cardinal $\lambda$. For $\alpha<\lambda$, put
\[
 A_\alpha=\{\beta:\beta\leq\alpha\},\qquad
 B_\alpha=A_\alpha\amalg\{\top\},
\]
where $\top$ lies above every element, and let $j_\alpha:A_\alpha\hookrightarrow B_\alpha$ be the inclusion. It splits by sending $\top$ to $\alpha$. The inclusions of initial segments, together with $\top\mapsto\top$, define a $\lambda$-chain in $\mathcal L$. Its index ordinal is $\lambda$-filtered. In the arrow category its colimit is
\[
 j:\lambda\hookrightarrow\lambda\amalg\{\top\}.
\]
A monotone retraction would have to send $\top$ to an ordinal dominating every $\beta<\lambda$. No such ordinal belongs to $\lambda$. Thus $j$ does not split. Since this construction works for every regular $\lambda$, the full image is not accessibly embedded.

Closing the image under retracts changes nothing. Suppose an arrow $j:A\to B$ is a retract of a split arrow $i:X\to Y$, with arrow-category maps $(u,v):j\to i$ and $(r,s):i\to j$ satisfying $ru=1_A$, $sv=1_B$. If $p:Y\to X$ retracts $i$, then
\[
 (rpv)j=rp(iu)=ru=1_A.
\]
Hence $j$ already splits. This proves directly the stronger rejection recorded in footnote 156.

The reason the diagram above does not lift to $\operatorname{Pos}^{R}$ is concrete: the chosen retraction of the new top changes from $\alpha$ to $\beta$ as the index increases. The squares preserve the inclusions but not their chosen retractions. Fullness of the essential image permits exactly these additional morphisms. In contrast, an $\Ind$-completion of compact objects embeds fully faithfully before any choices are forgotten; this is why the image used in 4.13.34 behaves well.
\end{dcomment}


Voici une question qui me paraît très voisine. Soient $\mathcal{M}$ une catégorie accessible, $A$ une sous-catégories pleine petite. Supposons que $\mathcal{M}$ satisfasse à $\mathrm{L}_{\pi_0}$ (p.ex. $\pi_0$ utile pour $\mathcal{M}$), et soit $\underline{S}$ la plus petite sous-catégorie strictement pleine de $\mathcal{M}$ contenant $A$, et stable par limites grandes devant $\pi_0$. Je conjecture que $\underline{S}$ est accessible dans $\mathcal{M}$. % [sic] sous-catégories -> sous-catégorie, singular after « une »

\msp{page 355}

De même, si $\Phi$ est un petit ensemble de petites catégories, tel que $\mathcal{M}$ soit stable par $\limind$ de type $\Phi$, on peut prendre l'enveloppe $\underline{S}$ de $A$ pour les $\limind$ grandes devant $\pi_0$ \textit{et} les
\printedpage{135}
% head: A. Grothendieck : Dérivateurs, Ch. XVIII, Édition des Universités de Montpellier II et Paris VII
% foot: Version 17 décembre 2025
$\limind$ de type $\Phi$, et je conjecture encore que $\underline{S}$ est accessible. Dans cette direction, on a du moins le résultat très partiel suivant.

\textbf{Proposition 4.13.34.} \textit{Soit $\mathcal{M}$ une catégorie accessible stable par petites $\limind$ filtrantes, et soit $\underline{S}$ une sous-catégorie strictement pleine de $\mathcal{M}$. Conditions équivalentes :}
\begin{itemize}[leftmargin=2.5em]
\item[\textit{a)}] \textit{$\underline{S}$ est accessible dans $\mathcal{M}$, et stable dans $\mathcal{M}$ par petites $\limind$ filtrantes.}
\item[\textit{b)}] \textit{Il existe une petite sous-catégorie pleine $A$ de $\mathcal{M}$, telle que $\underline{S}$ soit la plus petite sous-catégorie strictement pleine de $\mathcal{M}$ contenant $A$ et stable par petites $\limind$ filtrantes.}
\end{itemize}

\msp{page 356}

\dgrblock{dgrb1090}{The small-envelope conjectures at a fixed threshold}
\begin{dcomment}{development}{p355}{The small-envelope conjectures at a fixed threshold}
Grothendieck separates the false assertion about arbitrary essential images from a more controlled construction: start with a small collection of objects and close it under specified colimits already available in the ambient category. His proof of 4.13.34 treats all filtered colimits. The same method, combined with bounded completion of index categories, answers the two preceding envelope questions at a fixed filteredness threshold. We give the argument because the role of the size bound is otherwise easy to miss.

\textbf{Proposition (small colimit envelopes).}\ours\
Let $\kappa$ be an infinite regular cardinal. Let $\mathcal M$ be accessible, admitting all small $\kappa$-filtered colimits and the colimits of a small family $\Phi$ of small categories. If $A\subset\mathcal M$ is small, then its smallest replete full envelope closed under these colimits is accessible and accessibly embedded in $\mathcal M$.

We use the $\Ind$ extension result proved in 4.13.19: a small replete subcategory of $\mathcal M_\pi$, closed under a family of colimits of size at most $\pi$ including a nonzero limit-ordinal shape, has a fully faithfully embedded $\Ind_\pi$-completion closed under the same colimits. The closure under a limit-ordinal shape supplies the retract closure needed in applying that result.

\textit{Proof.} Choose an infinite cardinal $\pi$ which is useful for $\mathcal M$, with
\[
 \kappa\leq\pi,\qquad \pi^{<\kappa}=\pi,
 \qquad A\subset\mathcal M_\pi,
 \qquad |\operatorname{Fl}I|\leq\pi\quad(I\in\Phi).
\]
Such a bound exists. Start above a usefulness bound for $\mathcal M$, above the presentability ranks of the small family $A$, and above the sizes of the shapes; then take a sufficiently large power $2^c$. The exponentiation identity gives the displayed closure condition and the usual enlargement of usefulness.

Let $\Phi_\pi$ consist of the shapes in $\Phi$ and one representative of each $\kappa$-filtered category having at most $\pi$ arrows. There is a set of such representatives. The subcategory $\mathcal M_\pi$ is closed under these colimits: a colimit with at most $\pi$ arrows of $\pi^+$-presentable objects remains $\pi^+$-presentable, because $\pi^+$-filtered colimits of sets commute with limits indexed by at most $\pi$ arrows. Let $B$ be the smallest replete full subcategory of $\mathcal M_\pi$ containing $A$ and closed under $\Phi_\pi$-colimits. It is essentially small, being contained in $\mathcal M_\pi$.

We make its retract closure explicit. The ordinal $\kappa$ is $\kappa$-filtered and has at most $\pi$ arrows. For an idempotent $e:b\to b$, the diagram constant at $b$ with transition $e$ for every strict inequality has a colimit which splits $e$. Indeed, if $q:b\to c$ is its colimit map, then $qe=q$, while the compatible maps $e:b\to b$ induce $u:c\to b$ with $uq=e$; the colimit universal property gives $qu=1_c$. Thus $c\in B$, and repleteness includes every ambient retract of $b$.

Consequently the canonical functor $\Ind_\pi B\to\mathcal M$ is fully faithful. Write $\mathcal S'$ for its replete image. It is accessible, its $\pi^+$-filtered colimits are computed in $\mathcal M$, and 4.13.19 shows that it is closed there under $\Phi_\pi$-colimits. In particular it is already closed under the additional shapes $\Phi$. It remains to pass from bounded $\kappa$-filtered diagrams to arbitrary small ones.

Here is the index-category argument with its cardinal estimate. If $I$ is small and $\kappa$-filtered, every subcategory $K\subset I$ with at most $\pi$ arrows is contained in a $\kappa$-filtered subcategory $\overline K\subset I$ with at most $\pi$ arrows. For each diagram with fewer than $\kappa$ arrows in the current subcategory, choose a cocone in $I$, adjoin its arrows and vertex, and close under finite composition. There are at most $\pi^{<\kappa}=\pi$ diagrams to treat, including their possible small shapes, so each stage still has at most $\pi$ arrows. Iterate through $\kappa$ stages, taking unions at limit stages. Since $\kappa$ is regular, every diagram with fewer than $\kappa$ arrows in the final union is contained in one earlier stage and acquires its cocone at the next. The final union is the required $\overline K$. The empty diagram is included, so nonemptiness is also assured.

Let $P$ be the poset of all $\kappa$-filtered subcategories $J\subset I$ with at most $\pi$ arrows. It is a set and is $\pi^+$-filtered. Indeed, the union of at most $\pi$ such subcategories generates a subcategory with at most $\pi$ arrows, and the preceding construction supplies an upper bound in $P$. Every object and arrow of $I$ belongs to some member of $P$. These are actual subcategories, not full subcategories; taking full closure could add too many arrows.

Now let $D:I\to\mathcal S'$ be any small $\kappa$-filtered diagram. Each inner colimit below exists in $\mathcal S'$ by bounded closure, and the outer one exists by $\pi^+$-accessibility:
\[
 y=\varinjlim_{J\in P}\left(\varinjlim_{j\in J}D(j)\right).
\]
Both are computed in $\mathcal M$. A compatible family of cocones on the restrictions $D|_J$ gives a unique cocone on $D$: to define its component at $i$, choose a $J$ containing $i$; the result is independent of this choice because two such $J$ have a common upper bound. Naturality follows by choosing a $J$ containing the arrow in question. Conversely any cocone on $D$ restricts compatibly. Hence $y$ is the colimit of $D$, and $\mathcal S'$ is closed under all small $\kappa$-filtered colimits in $\mathcal M$.

Finally let $\mathcal S$ denote the envelope in the proposition. Every object of $B$ belongs to $\mathcal S$ by its construction. Every $\pi^+$-filtered category is $\kappa$-filtered, so every object of $\mathcal S'$, being a $\pi^+$-filtered colimit of objects of $B$, belongs to $\mathcal S$. The reverse inclusion follows from the closure properties just proved for $\mathcal S'$. Therefore $\mathcal S=\mathcal S'$, proving the assertion.\hfill$\square$

For the manuscript's two questions take $\kappa=\pi_0^+$, first with $\Phi$ empty and then with its specified additional shapes. The proof follows Grothendieck's own $\Ind$ and bounded-colimit method. We checked the cited primary papers on accessible constructions and images, together with the publisher descriptions of \citet{MakkaiPare1989} and \citet{AdamekRosicky1994}; we did not locate this exact envelope formulation in the material inspected. The marker records the proof developed for this commentary, without asserting priority for the theorem.
\end{dcomment}


On a a) $\Longrightarrow$ b), car si $\underline{S}$ est $\pi$-accessible dans $\mathcal{M}$, il est immédiat (sans autre condition de stabilité sur $\underline{S}$) que $\underline{S}$ est la plus petite sous-catégorie strictement pleine de $\mathcal{M}$ contenant $\underline{S}(\pi) \defeq \underline{S}\cap\mathcal{M}_\pi$ et stable par $\limind$ grandes devant $\pi$. Si donc on sait déjà que $\underline{S}$ est même stable par \textit{toutes} $\limind$ filtrantes petites, il s'ensuit que $\underline{S}$ est la plus petite sous-catégorie de $\mathcal{M}$ ayant cette propriété de stabilité plus forte, et contenant $A$.

Prouvons b $\Longrightarrow$ a. Soit $\pi$ un cardinal utile pour $\mathcal{M}$, et tel que $A \subset \mathcal{M}_\pi$. Soit $B$ la plus petite sous-catégorie strictement pleine de $\mathcal{M}$ contenant $A$ et stable par $\limind$ filtrantes de cardinal $\leqslant \pi$. On a donc $B \hookrightarrow \mathcal{M}_\pi$ et $B \subset \underline{S}$, et $\underline{S}$ est aussi la L-enveloppe de $B$. D'ailleurs $B$ est stable par
\msp{page 357}
facteurs directs. Considérons alors l'image essentielle $\underline{S}'$ de $\Ind_\pi(B)$ dans $\mathcal{M}$. On a $\underline{S}' \subset \underline{S}$, et tout revient à prouver que $\underline{S}' = \underline{S}$, d'où résultera en effet que $\underline{S}$ est $\pi$-accessible dans $\mathcal{M}$, a fortiori accessible dans $\mathcal{M}$, d'où a). Pour prouver que $\underline{S}' = \underline{S}$, par définition de $\underline{S}$ comme la L-enveloppe de $A$, il suffit de voir que $\underline{S}'$ est stable par petites $\limind$ filtrantes. Or cela résulte du fait que $B \subset \mathcal{M}_\pi$ est stable dans $\mathcal{M}_\pi$ par $\limind$ filtrantes de cardinal $\leqslant \pi$, et de 4.13.19, qui implique tout au moins, comme $B = \underline{S}' \cap \mathcal{M}_\pi$ (en vertu de 4.13.4) que $\underline{S}'$ est stable par $\limind$ filtrantes \textit{de cardinal} $\leqslant \pi$. Qu'il soit même stable par petites $\limind$ filtrante quelconques, résulte de ceci, et du fait qu'il l'est par $\limind$ grandes devant $\pi$. (Cf. 4.13.21 \dots) % [sic] filtrante quelconques -> filtrantes quelconques, plural agreement

\msp{page 358}

\textbf{Définition 4.13.35.} Soient $\mathcal{M}$ une catégorie accessible, stable par petites $\limind$ filtrantes, $S$ une partie de $\Ob \mathcal{M}$, $\underline{S}$ la sous-catégorie pleine correspondante dans $\mathcal{M}$. On dit que $S$ est \textit{L-accessible}, ou que $\underline{S}$ est L-accessible dans $\mathcal{M}$, si $\underline{S}$ est accessible dans $\mathcal{M}$, et de plus stable par petites $\limind$ filtrantes (\footnote[158]{\textbf{N.B.} L'ensemble des parties L-accessibles de $\Ob \mathcal{M}$ est stable par petites intersections, et par réunions \textit{finies}. L'image inverse d'une partie L-accessible par un foncteur (\footnotemark[159]) $f : \mathcal{M} \longrightarrow \mathcal{N}$ de catégories accessibles est L-accessible.}).\footnotetext[159]{un \textit{$L_0$-foncteur} : qui commute aux petites $\limind$ filtrantes.}

\textbf{Proposition 4.13.36.} \textit{Soient $\mathcal{M}$, $S$, $\underline{S}$ comme dans 4.13.35, et soit $\pi$ un cardinal utile pour $\mathcal{M}$. Supposons $\underline{S}$ stable dans $\mathcal{M}$ par petites $\limind$ filtrantes, ou seulement $\underline{S}(\pi) \defeq \underline{S} \cap \mathcal{M}_\pi$ stable par petites $\limind$ filtrantes de cardinal $\leqslant \pi$. Conditions équivalentes :}
\begin{itemize}[leftmargin=2.5em]
\item[\textit{a)}] \textit{$\underline{S}$ est $\pi$-accessible. (Il sera automatiquement stable par petites $\limind$ filtrantes, donc L-accessible.)}
\item[\textit{b)}] \textit{Pour tout $x \in \Ob \mathcal{M}$, on a $x \in \Ob \underline{S}$ si et seulement si $\underline{S}(\pi)/x$ est cofinale dans $\mathcal{M}_\pi/x$.}
\printedpage{136}
% head: A. Grothendieck : Dérivateurs, Ch. XVIII, Édition des Universités de Montpellier II et Paris VII
% foot: Version 17 décembre 2025
\item[\textit{c)}] \textit{Pour tout $x \in \Ob \mathcal{M}$, et toute représentation $x \simeq \varinjlim\limits_{i\in I} x_i$}
\msp{page 359}
\textit{qui est L-$\pi$-adaptée (cf. page 343, c), on a ceci : $x \in \Ob S$ si et seulement si $I_{\underline{S}}$ est cofinal dans $I$.}
\item[\textit{c}$'$\textit{)}] \textit{Comme c), avec $I$ ordonnée.}
\item[\textit{c}$''$\textit{)}] \textit{(Si on a choisi, pour tout $x \in \Ob \mathcal{M}$,} une \textit{représentation $\pi$-adaptée $x \simeq \varinjlim\limits_{I} x_i$) : $x \in \Ob \underline{S} \Longleftrightarrow$ l'ensemble $I_{\underline{S}}$ est cofinal dans $I$.}
\end{itemize}
\textsc{Démonstration}. C'est essentiellement l'énoncé de 4.13.22 dans le cas particulier où $\Phi =$ ensemble des catégories \textit{filtrantes} de cardinal $\leqslant \pi$, mais où, pour simplifier, on a supposé d'avance que $\underline{S}(\pi)$ satisfait à $\mathrm{L}(\Phi)$. L'équivalence de a) b) c) est contenue dans 4.13.22, ainsi que c $\Longrightarrow$ c$'$. Mais pour c$'$ $\Longrightarrow$ a, l'hypothèse accessoire $\pi^{\pi_0} = \pi$ et que les $J \in \Phi$ seront grandes devant $\pi_0$, n'est pas satisfaite ici, et il faut un démonstration indépendante, calquée sur celle des pages 333-336. Il suffit pour cela, à l'endroit \texttt{[de]} $\square$, page 333, d'invoquer le % [sic] un démonstration -> une démonstration, gender agreement

\textbf{Lemme 4.13.37.} \textit{Soient $\mathcal{M}$ une catégorie}
\msp{page 360}
\textit{accessible stable par $\limind$ filtrantes, $\pi$ un cardinal utile pour $\mathcal{M}$. Alors pour tout $x \in \Ob \mathcal{M}$, on a une représentation}
\[ x \;=\; \varinjlim_{i\in I} x_i \]
\textit{L-$\pi$-adaptée (i.e. avec}
\begin{itemize}[leftmargin=2.5em]
\item[\textit{1)}] \textit{$I$ équivalente à une petite catégorie, et $I$ grande devant $\pi$,}
\item[\textit{2)}] \textit{$I$ stable par $\limind$ filtrantes de cardinal $\leqslant \pi$,}
\item[\textit{3)}] \textit{$I \longrightarrow \mathcal{M}$ commute à ces $\limind$ ,}
\item[\textit{d)}] \textit{$x_i \in \Ob \mathcal{M}_\pi \ \forall\, i \in I$),} % [sic] d) -> 4), the items are numbered 1), 2), 3)
\end{itemize}
\textit{avec de plus $I$} ordonné.

En effet, on part de
\[ x \;=\; \varinjlim_{j\in J} x_j \ , \qquad\quad J \;\simeq\; \underbrace{(\mathcal{M}_\pi/x)_0}_{\substack{\text{sous-catégorie réduite}\\ \text{équivalente à } \mathcal{M}_\pi/x}} \qquad . \]
On définit $I$ comme l'ensemble des sous-catégories \textit{filtrantes} $J'$ de $J$, de cardinal \texttt{[“cardinal” biffé]} $\leqslant \pi$, $I$ ordonné par inclusion, et on pose, pour $J' \in I$,
\[ x_{J'} \;=\; \varinjlim_{j\in J'} x_j \ . \]
Il est immédiat qu'on obtient ainsi un système inductif dans $\mathcal{M}$, indexé par $I$.

\msp{page 361}

D'autre part, il est clair que $I$ est petite, je dis que $I$ est stable par $\limind$ filtrantes de cardinal $\leqslant \pi$, et que le foncteur $J' \longmapsto x_{J'} : I \longrightarrow \mathcal{M}$ y commute, enfin que les $x_{J'}$ sont dans $\mathcal{M}_\pi$. Il reste à voir que $I$ est grande devant $\pi$, i.e. que si on a $(J'_r)_{r\in A}$, avec $\card A \leqslant \pi$, les $J'_r$ dans $I$, alors $\exists\, J' \in I$ qui majore les $J'_r$. Soit $K$ la sous-catégorie de $J$ engendrée par les $J'_r$, on voit facilement que $\card K \leqslant \pi$. Tout \texttt{[“tout” biffé]} revient à prouver :

\printedpage{137}
% head: A. Grothendieck : Dérivateurs, Ch. XVIII, Édition des Universités de Montpellier II et Paris VII
% foot: Version 17 décembre 2025
\textbf{Lemme 4.13.38} (\footnote[160]{\textbf{N.B.} Pour la démonstration de 4.13.37, il suffit de prouver 4.13.38 dans le cas où $J$ est ordonnée, auquel cas les contorsions de la fin sont inutiles.}). \textit{Soient $J$ une catégorie grande devant $\pi$, $K$ une sous-catégorie de $J$ de cardinal $\leqslant \pi$. Alors il existe une sous-catégorie $K'$} filtrante \textit{de $J$, contenant $K$, de cardinal $\leqslant \pi$. Si on suppose que l'on n'est pas, dans le cas où $J$ contient un objet maximal, et que $K$ contienne tous les objets maximaux de $J$,}
\msp{page 362}
\textit{on peut même prendre $K'$ telle que $\Ob K' = \Ob K \cup \{e\}$, où $e$ est un} objet final \textit{de $K'$, et $K$ est une sous-catégorie pleine de $K'$. (Donc si $K' \neq K$, $K'$ est isomorphe au cône à droite construit sur $K$, “déduit de $K$ en lui ajoutant un objet final”.)}

\dgrblock{dgrb1091}{Making the presentation continuous}
\begin{dcomment}{exposition}{p361}{Making the presentation continuous}
The ordered presentation in Lemma 4.13.37 has a purpose beyond replacing a category by a poset. Grothendieck needs the index poset to contain the bounded filtered colimits used to improve an approximation, and he needs the presentation itself to preserve them. The ordinary statement that every object is a filtered colimit of compact objects supplies neither requirement by itself.

There is a short proof of the bounded enlargement used here. In any filtered category $J$, start with a subcategory $K$ having at most the infinite cardinal $\pi$ arrows. For every pair of its objects choose a common target and arrows to it. For every parallel pair of its arrows choose an equalizing arrow. Adjoin these choices and close under finite composition; the resulting subcategory still has at most $\pi$ arrows. Repeat this construction countably many times. The union $K_\infty$ has at most $\pi$ arrows and is filtered: every finite configuration occurs in one stage, and its required witnesses occur at the next. Include one object at the start if $K$ is empty. This proves the first assertion of 4.13.38 under the weaker assumption that $J$ is ordinarily filtered. The stronger ambient filteredness is needed elsewhere, including the one-vertex variant, but not for this bounded enlargement.

Apply this to the small skeleton $J$ of $\mathcal M_\pi/x$. Let $P$ be the poset of its filtered subcategories of size at most $\pi$. Given at most $\pi$ members, generate their union under composition and apply the enlargement above. This makes $P$ $\pi^+$-filtered. A filtered union of at most $\pi$ members is itself filtered and still of size at most $\pi$; it is their least upper bound in $P$. Thus the required bounded filtered colimits in the index poset are simply unions.

For $J'\in P$, put $x_{J'}=\varinjlim_{j\in J'}x_j$. Each $x_j$ is $\pi^+$-presentable, and the diagram has at most $\pi$ arrows, so $x_{J'}$ is again $\pi^+$-presentable. The colimit comparison for a filtered union of subcategories is an isomorphism: a cocone on the union is exactly a compatible collection of cocones on its pieces. Therefore $J'\mapsto x_{J'}$ preserves the required colimits. The same description of cocones, now over all of $P$, gives $x\simeq\varinjlim_Px_{J'}$. This supplies all four conditions of an L-$\pi$-adapted presentation, including the continuity which drives the cofinal-membership criterion.

That continuity must also be read into condition c$''$ of 4.13.36, where the printed text abbreviates the phrase to ``$\pi$-adaptée''. To see what fails if it is omitted, let $\mathcal M$ be sets and $\underline S$ the full subcategory of sets of cardinal at most one. It is accessible and closed under filtered colimits. Over the ordinal $\pi^+$ take the constant object $\{0,1\}$, with every strict transition the map sending both points to $0$. This is a $\pi^+$-filtered presentation by compact objects and its colimit is a singleton, but no stage lies in $\underline S$. At a limit stage the displayed object is still a two-element set, whereas the colimit of the earlier constant-transition diagram is a singleton. Precisely the missing continuity detects the defect.
\end{dcomment}


\textsc{Démonstration}. Si $K$ admet un objet final, on prend $K' = K$. Sinon, comme $J$ est grande devant $\card \Ob K$, il existe un $j \in \Ob J$ qui majore tous les objets de $K$. Choisissons, pour tout $k \in \Ob K$, une flèche
\[ \alpha_k \;:\; k \longrightarrow j \ . \]
Pour toute flèche
\[ u \;:\; k' \longrightarrow k \]
dans $K$, considérons la double-flèche
\[ (\alpha_k \circ u, \alpha_{k'}) \;:\; k' \rightrightarrows j \ . \]
L'ensemble de ces double-flèches, de but $j$, est de cardinal $\leqslant \card \Fl K \leqslant \pi$, donc $J$ étant grande devant $\pi$, il existe une flèche $\rho : j \longrightarrow e$ dans $J$ qui égalise toutes ces double-flèches. Soit $K' \subset J$, ayant comme objets $\Ob K \cup \{e\}$, et ayant comme flèches, en plus de
\msp{page 363}
celles de $K$, les flèches $\rho\alpha_k : k \longrightarrow e$ (pour $k \in \Ob K$), et $\id_e$. \textit{Si} $e \notin \Ob K$, on voit que c'est bien une sous-catégorie, admettant $e$ comme objet final et $K$ comme sous-catégorie pleine. Si on avait $e \in \Ob K$, et s'il existe un objet $e'$ de $J$ qui majore $e$ et qui n'est pas dans $K$, on se ramène aussitôt au cas précédent, en remplaçant $e$ par $e'$.

$\Biggl($ (\footnote[161]{La construction qui suit est contenue dans l'énoncé plus général et plus précis 4.13.39.}). Donc il reste le cas où il existe un objet $e$ dans $K$ qui est un plus grand élément dans $K$, et tel que tout objet de $J$ majorant $e$ soit dans $K$. Cela implique, $J$ étant filtrante, que $e$ est un plus grand objet de $J$. (Car si $x \in J$, $\exists\, x'$ majorant commun de $x$ et $e$, donc $x'$ est dans $K$, donc majoré par $e$, qui majore donc $x$.) On voit comme tantôt qu'il existe une flèche $p : e \longrightarrow j$ dans $J$, qui égalise toutes les double-flèches $k \rightrightarrows e$ de but $e$ \textit{dans}~$K$. Quitte à composer avec $j \longrightarrow e$, on peut supposer $j = e$. On a donc un endomorphisme $p$ de $e$ dans $J$, tel que $pu = pv$ pour toute double-flèche $k \overset{u}{\underset{v}{\rightrightarrows}} e$ dans $K$, de but $e$.

\msp{page 364}

Pour $k \in \Ob K$, soit
\[ f_k \;:\; k \longrightarrow e \]
la valeur commune des $pu$, pour $u \in \Hom_K(k,e)$. On voit tout de suite que si on a $k' \xrightarrow{v} k$ dans $K$, alors
\[ f_{k'} \;=\; f_k \circ v \ . \]
\printedpage{138}
% head: A. Grothendieck : Dérivateurs, Ch. XVIII, Édition des Universités de Montpellier II et Paris VII
% foot: Version 17 décembre 2025
En particulier, prenant $k' = e$, on trouve
\[ f_k \circ v \;=\; f_e \,(=\ p\ ) \qquad\quad \text{si } v : e \longrightarrow k \ . \]
Soit donc $K'$ la sous-catégorie de $J$ ayant mêmes objets que $K$, et comme flèches, en plus de celles de $K$, les flèches composées $k \xrightarrow{u} e \xrightarrow{p^r} e \xrightarrow{v} k'$, avec $u$, $v$ dans $K$, et $r \geqslant 1$, i.e. les flèches composées $k \xrightarrow{f_k} e \xrightarrow{p^{r'}} e \xrightarrow{v} k'$, avec $k, k' \in \Ob K$, $v \in \Fl K$, $r' \geqslant 0$. Il faut voir que c'est stable par composition (pour avoir bien une sous-catégorie). C'est clair pour $vu$ si $u$ ou $v$ sont dans $\Fl K$. Si \texttt{[ni]} $u$ ni $v$ ne le sont, on doit prendre un composé
\[
\begin{tikzcd}[column sep=large]
k \arrow[r, "u_1"] & e \arrow[r, "p^r", "r\geqslant 1"'] & e \arrow[r, "u_2"] \arrow[rr, bend right=20] & k' \arrow[r, "v_1"] & e \arrow[r, "p^s", "s\geqslant 1"'] & e \arrow[r, "v_2"] & k'' \ ,
\end{tikzcd}
\]
avec $u_1$, $u_2$, $v_1$, $v_2$ dans $\Fl K$, $r, s \geqslant 1$. Or on a
\[ p(v_1u_2) \;=\; p \ , \qquad\quad \text{d'où} \qquad\quad p^s(v_1u_2) \;=\; p^{s-1}(pv_1u_2) \;=\; p^{s-1}p \;=\; p^s \ , \]
donc le composé est
\[ k \xrightarrow{\;u_1\;} e \xrightarrow{\;p^{r+s}\;} e \xrightarrow{\;v_2\;} k'' \ , \]
qui
\msp{page 365}
est donc bien dans $\Fl K'$.

Ainsi, $K'$ est une sous-catégorie de $J$, mais est-elle filtrante ? La condition PS 1$'$ est évidente, puisque $e$ est un plus grand élément de $J$, reste la condition PS 2 de la double-flèche. On voit aisément que cette condition est satisfaite dans $K'$, si et seulement si elle est satisfaite pour toute double-flèche
\[ e \overset{p^r}{\underset{p^s}{\rightrightarrows}} e \qquad\quad \text{avec } r, s \geqslant 0 \ . \]
(En effet, on se ramène aussitôt aux double-flèches de but $e$, et on note que toute flèche de $K'$ de but $e$ s'écrit comme un composé $k \xrightarrow{u} e \xrightarrow{p^r} e \xrightarrow{v} e$, avec $r \geqslant 0$ et $u$, $v$ des flèches dans $K'$. En composant avec $e \xrightarrow{p} e$, on trouve alors, comme $pv = pu$, $k \xrightarrow{u} e \xrightarrow{p^{r+1}} e$, soit $k \xrightarrow{f_k} e \xrightarrow{p^r} e$.) Mais on voit aussitôt que la possibilité d'égaliser $(p^r, p^s)$ dans $K'$ équivaut à l'existence de $t \geqslant 0$ tel que $p^{r+t} = p^{s+t}$. Donc $K'$ est filtrante si et seulement si $r \longmapsto p^r$ est une suite stationnaire (faire $r = 0$, $s = 1$, d'où $p^r = p^{r+1}$ pour $r$ grand). Ce n'est
\msp{page 366}
pas forcément le cas.

Posons donc $K = K_0$, $K' = K_1$, et construisons par induction une suite croissante de sous-catégories $K_i$ ($i \in \mathbf{N}$), avec $\Ob K_i = \Ob J$, et (pour $i \geqslant 1$) $K_i$ étant déduites de $K_{i-1}$ en lui rajoutant une flèche $p_i : e \longrightarrow e$, comme ci-dessus. Ainsi toute double-flèche dans un $K_i$ peut s'égaliser dans $K_{i+1}$. Prenons $K_\infty = \limind K_i$, $K_\infty$ est une sous-catégorie \textit{filtrante} de $J$, contenant $K$ et de cardinal $\leqslant \pi$, q.e.d. $\Biggr)$

On peut préciser cette construction ainsi :

% ---- from der-ch18_p141-145.tex
\printedpage{139}
% head: A. Grothendieck : Dérivateurs, Ch. XVIII, Édition des Universités de Montpellier II et Paris VII
% foot: Version 17 décembre 2025

\noindent\textbf{Corollaire 4.13.39.} {\itshape Sous les conditions de 4.13.38, supposons que $K$ contienne un objet $e$ maximal dans $J$. Soit $\Sigma$ l'ensemble des endomorphismes $p$ de $e$ dans $J$, tels que pour toute double-flèche $k \overset{u}{\underset{v}{\rightrightarrows}} e$ \textup{dans} $K$, on ait $pu=pv$. Alors $\Sigma$ est stable par composition, et il existe une partie non vide $\Sigma_\infty$ \edbr{[}$=\Sigma_0$\edbr{]} de $\Sigma$, telle que pour $p,p'\in\Sigma_0$, existe
\msp{page 367}
$q\in\Sigma_0$ tel que $qp=qp'$ (i.e. telle que la sous-catégorie $J_0$ de $J$ dont l'ensemble des objets est $\{e\}$, et l'ensemble des flèches $\{\id_e\}\cup\Sigma_\infty$, soit filtrante). Si $\Sigma_0$ est choisie ainsi, soit $\mathcal{F}$ l'ensemble réunion de $\Fl K$ et de l'ensemble des flèches composées
\[
k \xrightarrow{\ u\ } e \xrightarrow{\ p\ } e \xrightarrow{\ v\ } k'\ , \tag*{$(*)$}
\]
avec $u,v\in\Fl J$, $p\in\Sigma_\infty$. Alors $\mathcal{F}$ est l'ensemble des flèches d'une sous-catégorie $K'$ de $J$, ayant mêmes objets que $J$, et admettant $e$ pour plus grand élément. De plus, $K'$ contient $K$ et est filtrante.}

\noindent\textsc{Démonstration}. Que $\Sigma$ soit stable par composition est clair. Pour construire $\Sigma_\infty$ ayant les propriétés énoncées, on note que, $J$ étant grande devant $\pi$, $\Sigma$ est non vide (\footnote[162]{cf. argument plus bas pour construction de $p_i$.}). On construit par récurrence une suite d'éléments
\[
p_0\ ,\ p_1\ ,\ \ldots\ ,\ p_n\ ,\ \ldots
\]
de $\Sigma$, en prenant pour $p_0$ un élément quelconque de $\Sigma$, et pour $i\geqslant 1$, prenant $p_i\in\Sigma$ tel que pour toute double-flèche $(p,p')$ du sous-monoïde $M_{i-1}$ de $\End_J(e)$ engendré par $(p_j)_{0\leqslant j\leqslant i-1}$, on ait $p_ip=p_ip'$. Un tel $p_i$ existe, car prenant l'ensemble des double-flèches de la forme précédente $(p,p')$, avec $p,p'\in M_{i-1}$, et de celles de la forme $(u,v) : k \rightrightarrows e$, avec $u,v\in\Fl J$, l'ensemble de
\msp{page 368}
toutes ces double-flèches est de cardinal $\leqslant\pi$ ($M_{i-1}$ étant dénombrable, donc de cardinal $\leqslant\pi$), donc il existe une flèche $e\longrightarrow j$ dans $J$ qui les égalise toutes ($J$ étant grande devant $\pi$). Quitte à composer avec $j\longrightarrow e$ ($e$ étant objet maximal de $J$), on trouve le $p_i : e\longrightarrow e$ cherché.

Soit alors $\Sigma_0$ le sous-ensemble multiplicatif de $\Sigma$ engendré par les $p_i$ ($i\in\mathbf{N}$), limite inductive des ensembles similaires $\Sigma_i$ engendrés par les $p_j$ pour $j\in[0,i]$. Si $p$, $p'$ sont dans $\Sigma_\infty$, il sont dans un même $\Sigma_i$, donc on aura $p_{i+1}p=p_{i+1}p'$, ce qui prouve que $\Sigma_\infty$ a la propriété filtrante exigée. % [sic] il sont -> ils sont, agreement of the verb with the plural subject

Prouvons maintenant que $\mathcal{F}$ est l'ensemble de flèches d'une sous-catégorie $K'$ de $J$, ayant $\Ob K$ comme ensemble d'objets. Il reste à établir la stabilité par composition $vu$. Mais si $u$ et $v$ sont dans $\Fl K$, ou si l'un des deux est dans $\Fl K$, l'autre non (donc de la forme $(*)$), cette stabilité est évidente. Reste à voir le cas où ni \edbr{[}$u$\edbr{]} ni $v$ sont dans $\Fl K$, donc sont tous deux de la forme $(*)$ :
\[
k \xrightarrow{\ u\ } e \xrightarrow{\ p\ } e \xrightarrow{\ v\ } k' \xrightarrow{\ u'\ } e \xrightarrow{\ p'\ } e \xrightarrow{\ v'\ } k''\ ,
\]
avec $u,v,u',v'\in\Fl K$, $p,p'\in\Sigma_\infty$. Comme $(u'v,\id_e)$ est une double-flèche de
\printedpage{140}
% head: A. Grothendieck : Dérivateurs, Ch. XVIII, Édition des Universités de Montpellier II et Paris VII
% foot: Version 17 décembre 2025
\msp{page 369}
$K$ de but $e$, on a
\[
p'u'v \;=\; p'\id_e \;=\; p'\ ,
\]
donc le composé en question s'écrit aussi
\[
k \xrightarrow{\ u\ } e \xrightarrow{\ p'p\ } e \xrightarrow{\ v'\ } k''\ ,
\]
où $p'p\in\Sigma_\infty$ puisque $\Sigma_\infty$ est stable par composition, donc il est dans $\mathcal{F}$, q.e.d.

Il est clair que $K'$ contient $K$, et ayant mêmes objets que $K$ qui admet $e$ comme plus grand élément, que $e$ est un plus grand élément de $K'$. Reste à prouver que $K'$ est filtrante. Comme $K'$ a un plus grand élément, il reste à voir que toute double-flèche $k\rightrightarrows l$ dans $K'$ peut être égalisée dans $K'$. Quitte à composer avec $l\longrightarrow e$, on peut supposer $l=e$. Or on voit que toute flèche $k\longrightarrow e$ dans $K'$ est de la forme
\[
k \xrightarrow{\ u\ } e \xrightarrow{\ p\ } e \xrightarrow{\ v\ } e\ ,
\]
avec $u,v\in\Fl K$ et $p\in\Sigma_\infty\cup\{\id_e\}$. Quitte à composer avec $p_0\in\Sigma_\infty$ fixé, cette flèche devient (puisque $p_0v=p_0\id_e=p_0$)
\[
k \xrightarrow{\ u\ } e \xrightarrow{\ p_0p\ } e\ , \qquad\quad \text{i.e. } p_0pu\ .
\]
Pour une flèche $v'p'u' : k\longrightarrow e$, on trouve $p_0p'u'$. Soit alors $q\in\Sigma_0$ tel que $q(p_0p)=q(p_0p')$, soit $r$ ($q$ existe puisque $p_0p,p_0p'\in\Sigma_\infty$). Donc $(q(p_0pu),q(p_0p'u'))=(ru,ru')$, et on a $ru=ru'$ puisque $r\in\Sigma_\infty$. Au total, $qp_0$ égalise les deux flèches $vpu$ et $v'p'u'$, q.e.d.

\msp{page 370}

\dgrblock{dgrb1092}{The exceptional vertex and its endomorphisms}
\begin{dcomment}{exposition}{p369}{The exceptional vertex and its endomorphisms}
Grothendieck pursues the one-vertex construction because it describes how much enlargement an approximation really needs. If a cocone vertex lies outside $K$, adding it with its compatible incoming arrows produces a terminal object immediately. If every possible vertex is already present, the obstruction lies in endomorphisms, and a weakly terminal object need not become terminal after a single equalization.

The exceptional-case proof needs one preliminary adjustment. A vertex $e$ which is maximal in the reachability preorder of $J$ receives arrows in $J$ from every object of $K$, but those arrows need not belong to $K$. Adjoin one such arrow from each object and close under composition. This preserves the object set and the bound $\pi$ on arrows. After this adjustment $e$ is weakly terminal in $K$, and the formula $f_k=pu$ is defined for every $k$. The later normal-form argument then applies exactly as written, with its local $J/K$ slips read correctly.

Here is the point of the powers of $p$. If $p$ equalizes all $K$-arrows with target $e$, then a new arrow can be written $vp^ru$ with $u,v$ in $K$. Products reduce because $p(v_1u_2)=p$ for every $K$-endomorphism $v_1u_2$ of $e$. But the powers $1,p,p^2,\ldots$ need not equalize one another. Filteredness of the resulting one-object part forces, by equalizing $1$ and $p$, an equality $p^t=p^{t+1}$ at some stage. Conversely that equality equalizes every pair of powers after further multiplication. Thus the additional iteration is mathematically necessary.

In 4.13.39 the countable sequence $p_i$ arranges this for all words already constructed. At stage $i+1$ one equalizes pairs in the monoid generated by $p_0,\ldots,p_i$, including the identity, together with the original parallel arrows of $K$ ending at $e$. There are at most $\pi$ pairs. Their simultaneous equalizer in $J$ can be returned to $e$ by maximality. Any two words occur at a finite stage, so a later $p_i$ equalizes them. The resulting semigroup gives the filtered enlargement claimed. The proof uses arrows of $K$ in its normal forms and retains $\operatorname{Ob}K$; the printed substitutions of $J$ at these places are recorded in the correction register.
\end{dcomment}


\section*{4.14\quad Critères d'accessibilité pour grosses catégories}

\noindent\textbf{Théorème 4.14.1.} {\itshape Soient $\mathcal{C}$ une catégorie équivalente à une petite catégorie, $\Phi$, $\Psi$ deux petits ensembles de catégories, équivalentes à des petites catégories. Pour toute catégorie $I$, on désigne par $I_{\bullet}$ (resp. ${}_{\bullet}I$) le cône droit (resp. gauche) sur $I$, déduit de $I$ en ajoutant un objet final (resp. initial) $e=e_I$. Soit $\Lambda_1$ (resp. $\Lambda_2$) un ensemble de couples $(I,\varphi)$ (resp. $(I,\psi)$), où $I\in\Phi$ (resp. $I\in\Psi$), et où
\[
\varphi\ :\ I \longrightarrow \mathcal{C} \qquad\quad \text{(resp. }\psi\ :\ I \longrightarrow \mathcal{C}\ \ \text{)} \tag*{\textup{(4.14.1.1)}}
\]
est un foncteur. Pour toute catégorie $\mathcal{N}$ stable (\footnote[163]{Cette propriété de stabilité n'est pas nécessaire pour la définition de $\uHom^{\Lambda_2}_{\Lambda_1}(\mathcal{C},\mathcal{N})$, et il est presque sûr que le théorème est valable sans faire cette hypothèse sur $\mathcal{N}$. Cf. 4.14.16, où on voit que c'est en tous cas inutile pour $\Psi$.}), par $\varinjlim$ de type $\Phi$ et par $\varprojlim$ de type $\Psi$ on désigne par
\[
\uHom^{\Lambda_2}_{\Lambda_1}(\mathcal{C},\mathcal{N})\ \subset\ \uHom(\mathcal{C},\mathcal{N}) \tag*{\textup{(4.14.1.2)}}
\]
la sous-catégorie strictement pleine de $\uHom(\mathcal{C},\mathcal{N})$ formée des foncteurs $f:\mathcal{C}\longrightarrow\mathcal{N}$ ayant la propriété suivante :
\begin{enumerate}[leftmargin=2.5em]
\item[\textit{a)}] $\forall\ (I,\varphi)\in\Lambda_1$, le foncteur composé
\[
I_{\bullet} \xrightarrow{\ \varphi\ } \mathcal{C} \xrightarrow{\ f\ } \mathcal{N}
\]
est $I$-exact à gauche, i.e. fait de $(f\varphi)(e_I)$ une $\varinjlim\limits_{I}$ de $f\varphi|I$.
\printedpage{141}
% head: A. Grothendieck : Dérivateurs, Ch. XVIII, Édition des Universités de Montpellier II et Paris VII
% foot: Version 17 décembre 2025
\msp{page 371}
\item[\textit{b)}] Pour tout $(I,\psi)\in\Lambda_2$, le foncteur composé
\[
{}_{\bullet}I \xrightarrow{\ \psi\ } \mathcal{C} \xrightarrow{\ f\ } \mathcal{N}
\]
est $I$-exact à droite, i.e. fait de $(f\psi)(e_I)$ une limite projective de $(f\psi)|I$.
\end{enumerate}
\noindent Avec ces notations : si $\mathcal{N}$ est une catégorie accessible, il en est de même de $\uHom^{\Lambda_2}_{\Lambda_1}(\mathcal{C},\mathcal{N})$.}

\dgrblock{dgrb1093}{Marked cones and categories of models}
\begin{dcomment}{context}{p370}{Marked cones and categories of models}
Grothendieck now turns the calculus of accessible parts into a way of constructing categories. A functor $\mathcal C\to\mathcal N$ carries raw diagram data; marking a cone or a cocone in $\mathcal C$ asks that some of those data satisfy a universal property. Section 4.14 treats the resulting categories of models uniformly. This is why the preceding theorem on inverse images matters: the condition that a canonical comparison be an isomorphism is an accessible condition on the functor.

The data in display (4.14.1.1) must include the whole marked cone. For a colimit condition one has $\varphi:I_\bullet\to\mathcal C$, not just $\varphi:I\to\mathcal C$; for a limit condition one has $\psi:{}_\bullet I\to\mathcal C$. The cone vertex is the object to which the universal property is attached. With the first of these data and an arbitrary $F:\mathcal C\to\mathcal N$, the comparison is
\[
 \varinjlim_{i\in I}F\varphi(i)\longrightarrow F\varphi(e_I).
\]
The full subcategory of functors making this map invertible expresses exactly the desired colimit condition. The dual comparison expresses a limit condition. Taking the small intersection of these conditions gives the theorem, once accessibility of the comparison functors has been established. The hypothesis that $\mathcal N$ admits the indicated colimits and limits makes those comparison functors defined on the entire ambient functor category.

This is the language of sketches: a small category or graph together with specified diagrams, cones and cocones, whose models realize the specified universal constructions. The background is the study of categories presented by small data, including \citet{GabrielUlmer1971}, and the theory of accessible categories developed by \citet{MakkaiPare1989}. The latter appeared in 1989 and is contemporaneous with this manuscript; it should not be described simply as a later verification of Grothendieck's work. The comparison here concerns the mathematical formulations, without asserting a documented line of influence.

For a concrete example, take the one-sorted finite-product theory of groups. Its objects are the finite powers $1,x,x^2,\ldots$, and its maps are tuples of group terms modulo the group identities. Mark the finite-product cones. A product-preserving functor to sets chooses a set $G$, sends multiplication, inverse and unit terms to operations on $G$, and makes the term equalities precisely the group axioms. Conversely those operations determine the functor. A natural transformation is determined by its component at $x$, and naturality for the operation terms says exactly that this component is a group homomorphism. Thus the model category is the usual category of groups. Finite presentations supply its compact generators; an arbitrary group is the filtered colimit of the finite generators-and-relations approximations mapping to it. This familiar example exhibits the full programme: the source encodes operations and equations, and the marked cones enforce the intended relation between the different arities.

The same method is useful for large source categories after restriction to compact objects, but then the chosen bound must control both the diagram shapes and the operations being imposed. This is a mathematical requirement on the restriction argument, not merely a matter of replacing a large set of formulas by a smaller notation.
\end{dcomment}


\noindent\textbf{Commentaire.} On utilisera ce théorème surtout quand on a $\Lambda_2=\emptyset$ ou $\Lambda_1=\emptyset$, auquel cas on note $\uHom_{\Lambda_1}(\mathcal{C},\mathcal{N})$, resp. $\uHom^{\Lambda_2}(\mathcal{C},\mathcal{N})$, la catégorie $\uHom^{\Lambda_2}_{\Lambda_1}(\mathcal{C},\mathcal{N})$. On appliquera le théorème surtout quand $\mathcal{C}$ elle-même est stable par $\varinjlim$ de type $\Phi$ et par $\varprojlim$ de type $\Psi$, et quand $\Lambda_1$ (resp. $\Lambda_2$) est déjà formé de foncteurs $\varphi$ (resp. $\psi$) qui sont $I$-exacts à gauche (resp. à droite). Donc $\uHom^{\Lambda_2}_{\Lambda_1}(\mathcal{C},\mathcal{N})$ est alors formée des $f:\mathcal{C}\longrightarrow\mathcal{N}$ qui commutent à certaines $\varinjlim$ de type $\Phi$, et à certaines $\varprojlim$ de type $\Psi$. En particulier, on trouve :

\msp{page 372}

\noindent\textbf{Corollaire 4.14.2.} {\itshape Soient $\mathcal{C}$ une catégorie équivalente à une petite catégorie, $\mathcal{N}$ une $\mathfrak{U}$-catégorie, $\Phi$ et $\Psi$ deux petits ensembles de catégories $I$, équivalentes à des petites catégories. On suppose $\mathcal{C}$ et $\mathcal{N}$ stables par $\varinjlim$ de type $\Phi$, par $\varprojlim$ de type $\Psi$. Soit $\uHom^{\Psi}_{\Phi}(\mathcal{C},\mathcal{N})$ la sous-catégorie strictement pleine de $\uHom(\mathcal{C},\mathcal{N})$ formée des $f:\mathcal{C}\longrightarrow\mathcal{N}$ qui commutent aux $\varinjlim$ de type $\Phi$ et aux $\varprojlim$ de type $\Psi$. Alors, si $\mathcal{N}$ est accessible, il en est de même de $\uHom^{\Psi}_{\Phi}(\mathcal{C},\mathcal{N})$.}

Voici des cas particuliers remarquables de ce corollaire :

\noindent\textbf{Corollaire 4.14.3.} {\itshape Soient $\mathcal{M}$, $\mathcal{N}$ deux catégories accessibles, et $\pi_0$ un cardinal infini, tel que $\mathcal{M}$ satisfasse à $L_{\pi_0}$. Soit $\uHom_{\pi_0}(\mathcal{M},\mathcal{N})$ la sous-catégorie strictement pleine de $\uHom(\mathcal{M},\mathcal{N})$ formée des foncteurs $\pi_0$-accessibles. Alors $\uHom_{\pi_0}(\mathcal{M},\mathcal{N})$ est accessible.}

\msp{page 373}

(Le cas où $\pi_0$ est utile pour $\mathcal{M}$ a déjà été traité.)

\noindent\textsc{Démonstration}. Dans le cas général, soit $\pi>\pi_0$ un cardinal utile pour $\mathcal{M}$, de sorte qu'on a
\[
\mathcal{M} \xleftarrow{\ \sim\ } \Ind_\pi(\mathcal{C}) \qquad\quad \text{avec } \mathcal{C}=\mathcal{M}_\pi\ .
\]
On peut supposer $\mathcal{M}=\Ind_\pi(\mathcal{C})$. Le fait que $\mathcal{M}$ satisfasse à $\mathrm{L}_{\pi_0}$ implique que $\mathcal{C}$ est stable par $\varinjlim$ grandes devant $\pi_0$, et de cardinal $\leqslant\pi$. Soit $\Phi$ l'ensemble des catégories de cette nature. $\Phi$ n'est pas petit, mais il existe $\Phi_0\subset\Phi$ petit, tel que toute $I\in\Phi$ soit isomorphe à une $I_0\in\Phi_0$, ce qui suffira pour la validité de 4.14.1 et 4.14.2 (où on fera ici $\Psi=\emptyset$). On peut considérer $\uHom_{!\,\pi_0}(\mathcal{M},\mathcal{N})$ comme une sous-catégorie strictement pleine de $\uHom_{!\,\pi}(\mathcal{M},\mathcal{N})\xrightarrow{\ \sim\ }\uHom(\mathcal{C},\mathcal{N})$ (cf. 4.12.3, p. 275). Pour un foncteur $f:\mathcal{M}\longrightarrow\mathcal{N}$ qui est $\pi$-accessible, dire qu'il est même $\pi_0$-accessible revient à dire que sa restriction à $\mathcal{C}$, i.e. son image dans $\uHom(\mathcal{C},\mathcal{N})$, commute aux $\varinjlim$ de type $\Phi$. (Cf. 4.13.20, page 341 (\footnote[164]{La référence n'est pas tout à fait adéquate. Il faudrait réunir dans un paragraphe à part tous les énoncés de ce type, dont certains utilisent le théorème de représentation indicielle.}), où on prend $\alpha=\alpha_\pi$, notant que $J_{\alpha_\pi}$ grande devant $\pi_0$ puisque $\pi>\pi_0$).

\msp{page 374}

Donc on trouve
\[
\uHom_{!\,\pi_0}(\mathcal{M},\mathcal{N}) \xrightarrow{\ \sim\ } \uHom_{\Phi}(\mathcal{C},\mathcal{N})\ ,
\]
\printedpage{142}
% head: A. Grothendieck : Dérivateurs, Ch. XVIII, Édition des Universités de Montpellier II et Paris VII
% foot: Version 17 décembre 2025
et le corollaire 4.14.2 donne la conclusion voulue.

\dgrblock{dgrb1094}{The target colimits in the functor-category argument}
\begin{dcomment}{status}{p374:et le corollaire 4.14.2 donne la conclusion voulue.}{The target colimits in the functor-category argument}
Grothendieck wants fixed-threshold accessible functors to form an accessible category even when their source is large. Restricting to a sufficiently large small subcategory of presentable objects is the right approach. It reduces preservation of all the relevant filtered colimits to preservation of a bounded family, provided the target has the colimits needed in the reduction.

Corollary 4.14.3 is stated with $L_{\pi_0}$ only on $\mathcal M$. Its proof then invokes Corollary 4.14.2 for shapes which are $\pi_0^+$-filtered and have size at most $\pi$. That corollary assumes these colimits exist in $\mathcal N$. Accessibility of $\mathcal N$ by itself gives colimits above some sufficiently large threshold, which need not be the specified $\pi_0$. Choosing $\pi$ useful for both categories repairs the equivalence with functors on $\mathcal M_\pi$, but does not supply the missing smaller-threshold colimits. Thus the printed deduction proves the assertion with the additional hypothesis $\mathcal N\models L_{\pi_0}$. The unrestricted case is not established by that deduction; it is part of the issue about colimit conditions without ambient colimit completeness explicitly raised in 4.14.16. We do not infer a counterexample to the unrestricted statement from this proof gap.

There is a separate, elementary correction to the sentence about the ordinal $\alpha_\pi$. Its index category is $\pi_0^+$-filtered when $\operatorname{cf}(\pi)>\pi_0$, not simply when $\pi>\pi_0$. A power-of-two rank as above ensures the needed cofinality. In Corollary 4.14.4 the target filtered-colimit hypothesis is included, but the restriction to $\mathcal M_\pi$ should also be made after simultaneous uniformization of the selected limit and colimit functors. Merely bounding the sizes of their shapes does not guarantee that their values on presentable diagrams remain presentable. These rank adjustments leave the intended method intact and state precisely which version its argument proves.
\end{dcomment}


\noindent\textbf{Corollaire 4.14.4.} {\itshape Soient $\mathcal{M}$, $\mathcal{N}$ deux catégories accessibles, $\Phi$ et $\Psi$ deux ensembles des catégories comme dans 4.14.2, on suppose que $\mathcal{M}$ et $\mathcal{N}$ satisfassent $L_{\pi_0}$, et \edbr{[soient]} stables par $\varinjlim$ de type $\Phi$ et par $\varprojlim$ de type $\Psi$. Soit $\uHom^{\Psi}_{\Phi,\pi_0}(\mathcal{M},\mathcal{N})$ la sous-catégorie strictement pleine de $\uHom(\mathcal{M},\mathcal{N})$ formée des foncteurs $\pi_0$-accessibles qui commutent aux $\varinjlim$ de type $\Phi$, et aux $\varprojlim$ de type $\Psi$. Alors $\uHom^{\Psi}_{\Phi,\pi_0}(\mathcal{M},\mathcal{N})$ est accessible.}

\noindent\textsc{Démonstration}. Cette fois, on choisit $\pi>\pi_0$ tel qu'on ait $\cardred I\leqslant\pi$ pour tout $I$ dans $\Phi$ ou dans $\Psi$. La catégorie envisagée s'interprète alors (posant encore $\mathcal{C}=\mathcal{M}_\pi$) comme
\msp{page 375}
\[
\uHom^{\Psi}_{\Phi,\pi_0}(\mathcal{M},\mathcal{N}) \xrightarrow{\ \sim\ } \uHom^{\Psi}_{\Phi'}(\mathcal{C},\mathcal{N})\ ,
\]
où $\Phi'$ est formée des catégories qui sont soit dans $\Phi$, ou qui sont grandes devant $\pi_0$ et de cardinal $\leqslant\pi$. Il suffit donc d'appliquer encore 4.14.2.

\noindent\textbf{Corollaire 4.14.5.} {\itshape Considérons les trois ensembles suivants :
\begin{enumerate}[leftmargin=2.5em]
\item[\textit{a)}] $\Phi_1$ est formée de toutes les catégories équivalentes à des petites catégories.
\item[\textit{b)}] $\Phi_2$ comme a), mais en restreignant aux catégories filtrantes.
\item[\textit{c)}] $\Phi_3$ formé des catégories discrètes, et de celles qui sont grandes devant $\pi_0$ ($\pi_0$ un cardinal infini donné).
\end{enumerate}
\noindent Soit $\Phi$ l'un des ensembles de catégories $\Phi_1$, $\Phi_2$, $\Phi_3$, $\Phi_2\cup\Phi_3$ ou la réunion d'un de ceux-ci et d'un petit ensemble $\Phi_0$ de catégories essentiellement petites (\footnote[165]{On pourrait faire ici un énoncé qui englobe 4.14.4 en rajoutant un
\begin{enumerate}[leftmargin=2.5em]
\item[d)] $\Phi_4$ formé des catégories grandes devant $\pi_0$,
\end{enumerate}
\noindent et prendre $\Phi$ de la forme de $\Phi_1$, $\Phi_2$, $\Phi_4$, $\Phi_3$, $\Phi_3\cup\Phi_2$, $\Phi_3\cup\Phi_4$, ou la réunion de ceux-ci et d'un $\Phi_0$.}). Soit $\Psi$ un petit ensemble de catégories. Soient $\mathcal{M}$, $\mathcal{N}$ des catégories accessibles, stables par $\varinjlim$ de type $\Phi$, et par $\varprojlim$ de type $\Psi$. Alors $\uHom^{\Psi}_{\Phi}(\mathcal{M},\mathcal{N})$ est accessible.}

\msp{page 376}

\noindent\textsc{Démonstration}. Soit $\pi'_0\geqslant\pi_0$ un cardinal utile pour $\mathcal{M}$, $\pi$ un cardinal ayant les propriétés suivantes.
\begin{enumerate}[leftmargin=2.5em]
\item[a)] $\pi\geqslant\sup(\pi'_0,\cardred\mathcal{M}_{\pi'_0})$, $\pi>\pi_0$.
\item[b)] $\pi^{\pi'_0}=\pi$.
\item[c)] $I\in\Phi_0\cup\Psi\Longrightarrow\cardred I\leqslant\pi$.
\end{enumerate}
\noindent Soit $\mathcal{C}=\mathcal{M}_\pi$, d'où $\mathcal{M}\simeq\Ind_\pi(\mathcal{C})$. Dans tous les cas envisagés, $\Phi$ contient l'ensemble des catégories essentiellement petites qui sont grandes devant $\pi_0$, a fortiori celles qui sont grandes devant $\pi$, donc la catégorie $\uHom^{\Psi}_{\Phi}$ s'interprète comme une sous-catégorie strictement pleine de
\[
\uHom(\mathcal{C},\mathcal{N})\ .
\]
Si d'autre part $\Phi_\pi$ désigne l'ensemble des $I\in\Phi$ telles que $\card I\leqslant\pi$, on trouve
\[
\uHom^{\Psi}_{\Phi}(\mathcal{M},\mathcal{N})\ \simeq\ \uHom^{\Psi}_{\Phi_\pi}(\mathcal{C},\mathcal{N})\ ,
\]
et on applique 4.14.2.

\printedpage{143}
% head: A. Grothendieck : Dérivateurs, Ch. XVIII, Édition des Universités de Montpellier II et Paris VII
% foot: Version 17 décembre 2025
\msp{page 377}

\dgrblock{dgrb1095}{A bounded sketch for a large functor category}
\begin{dcomment}{exposition}{p376}{A bounded sketch for a large functor category}
Grothendieck is trying to make preservation of universal constructions into an accessible condition. The domain may have a proper class of objects and the demand may concern all small colimits. The preceding work with cardinals serves to replace that demand by a set of tests on one essentially small category. This is why these corollaries belong to the same argument as the long cofinality constructions: a good presentation rank makes the large category recoverable from its bounded part.

Choose $\pi$ useful for $\mathcal M$ and write $\lambda=\pi^+$, so that the objects of $\mathcal M_\pi$ are $\lambda$-presentable in modern terminology. Choose a small skeleton $\mathcal A$ of that subcategory. A functor preserving $\lambda$-filtered colimits is determined, including its natural transformations, by its restriction to $\mathcal A$. Its extension is expressed by
\[
 F(X)\simeq\varinjlim_{(A\to X)\in\mathcal A/X}F(A).
\]
The choice of a common rank must ensure that the target has these filtered colimits. Preservation of additional shapes is then tested by comparison arrows associated to a set of bounded cones. One asks that these arrows become invertible. The category of candidates is therefore obtained from an accessible diagram category by inverse images of an accessible isomorphism subcategory and a small intersection. These are precisely the operations established in the preceding section.

The colimit assertion in 4.14.6 illustrates both the strength and the exact range of this method. If $\mathcal M$ and $\mathcal N$ are locally presentable, the category of cocontinuous functors between them is accessible by the bounded-sketch argument. It is also cocomplete. Given a small diagram $(F_j)$ of cocontinuous functors, define $F(X)=\varinjlim_jF_j(X)$. For a small diagram $(X_i)$ one has
\[
\begin{aligned}
 F(\varinjlim_iX_i)
 &\simeq\varinjlim_j\varinjlim_iF_j(X_i)\\
 &\simeq\varinjlim_i\varinjlim_jF_j(X_i)
 \simeq\varinjlim_iF(X_i).
\end{aligned}
\]
The interchange is just the universal property of a colimit over the product indexing category; no exactness condition on $\mathcal N$ enters. Thus this functor category is locally presentable. The first assertion of 4.14.6, with only accessibility of $\mathcal N$, is not justified by the cited corollary: the latter assumes the target colimits that the former omits. We retain that distinction in the audit instead of silently adding a hypothesis.

The tensor product in 4.14.8 uses the other side of the construction. Its objects are functors $\mathcal M^{op}\to\mathcal N$ carrying colimits in $\mathcal M$ to limits in $\mathcal N$. The relevant constraints are limit-cone constraints, and the argument on manuscript page 389 shows that these can be recognized even when other limits in the target do not exist. Here $\mathcal N$ is complete, so limits of such functors are computed pointwise: two successive limits commute. Consequently the tensor category is accessible and complete, hence locally presentable. This argument does not require the erroneous covariant representability assertion that follows.

A small calculation makes the construction concrete. Every set $X$ is the coproduct of $X$ copies of the singleton. Therefore a functor $F:\mathbf{Set}^{op}\to\mathcal N$ taking colimits to limits must satisfy $F(X)\simeq F(1)^X$. For a map $a:X\to Y$, its image sends a family indexed by $Y$ to the family indexed by $X$ obtained by composition with $a$. Conversely, for any $N\in\mathcal N$, the functor $X\mapsto N^X$ takes colimits of sets to limits: applying $\Hom(T,-)$ identifies the claim with the familiar contravariant representable functor $\Hom_{\mathbf{Set}}(-,\Hom(T,N))$. A natural transformation between these functors is determined by its component at $1$, since projections to the singleton coordinates detect it. Hence $\mathbf{Set}\boxtimes\mathcal N\simeq\mathcal N$. The unit of the tensor construction is already visible in this elementary case.
\end{dcomment}


En particulier, on trouve :

\noindent\textbf{Corollaire 4.14.6.} {\itshape Soient $\mathcal{M}$, $\mathcal{N}$ deux catégories accessibles, supposons $\mathcal{M}$ stable par petites $\varinjlim$ (i.e. $\mathcal{M}$ est une catégorie paratopos). Alors la catégorie $\uHom_{!}(\mathcal{M},\mathcal{N})$ des foncteurs $f:\mathcal{M}\longrightarrow\mathcal{N}$ qui commutent aux petites $\varinjlim$ est accessible. Si $\mathcal{N}$ est elle aussi un paratopos, $\uHom_{!}(\mathcal{M},\mathcal{N})$ est un paratopos.}

\noindent\textsc{Démonstration}. La première assertion est le cas particulier de 4.14.5 relatif au cas où $\Phi=\Phi_0$, $\Psi=\emptyset$. La deuxième assertion en résulte, puisque $\uHom_{!}(\mathcal{M},\mathcal{N})$ est visiblement stable par petites $\varinjlim$ , donc c'est un paratopos étant accessible (4.8.9, page 207).

\noindent\textbf{Corollaire 4.14.7.} {\itshape Soient $\Phi$, $\Psi$, $\mathcal{M}$ comme dans 4.14.5. Alors pour toute catégorie accessible $\mathcal{N}$, la catégorie
\msp{page 378}
\[
\uHom^{\Phi}_{\Psi}(\mathcal{M}^{\circ},\mathcal{N})\ \simeq\ (\uHom^{\Psi}_{\Phi}(\mathcal{M},\mathcal{N}^{\circ}))^{\circ}
\]
est accessible.}

\noindent\textsc{Démonstration}. On procède comme pour 4.14.5, pour expliciter
\[
\uHom^{\Psi}_{\Phi}(\mathcal{M},\mathcal{N}^{\circ}) \xrightarrow{\ \sim\ } \uHom^{\Psi}_{\Phi_\pi}(\mathcal{C},\mathcal{N}^{\circ})\ ,
\]
d'où
\[
\uHom^{\Phi}_{\Psi}(\mathcal{M}^{\circ},\mathcal{N}) \xrightarrow{\ \sim\ } \uHom^{\Phi_\pi}_{\Psi}(\mathcal{C}^{\circ},\mathcal{N})\ ,
\]
et on applique 4.14.2 à $\mathcal{C}^{\circ}$, $\Psi$, $\Phi_\pi$.

\noindent\textbf{Corollaire 4.14.8.} {\itshape Soient $\mathcal{M}$ et $\mathcal{N}$ deux catégories paratopos. Alors
\[
\mathcal{M}\boxtimes\mathcal{N}\ \simeq\ \uHom^{!}(\mathcal{M}^{\circ},\mathcal{N}) \qquad\quad \text{(cf. 4.2 et 4.3)}
\]
est un paratopos.}

En effet, en vertu de 4.14.7 (avec $\Phi=\Phi_\pi$ [?], $\Psi=\emptyset$), c'est une catégorie accessible.

\msp{page 379}

D'autre part, il est immédiat qu'elle est stable par petites $\varprojlim$ , $\mathcal{N}$ l'étant. Donc on conclut encore par 4.8.9 (page 207) que c'est une catégorie paratopos.

\noindent\textbf{Corollaire 4.14.9.} {\itshape Soient $\mathcal{M}$, $\mathcal{N}$ deux catégories paratopos, et $\uHom^{!}(\mathcal{M},\mathcal{N})$ la catégorie des foncteurs $\mathcal{M}\longrightarrow\mathcal{N}$ qui commutent aux petites $\varprojlim$ . On a alors une équivalence de catégories
\[
\begin{aligned}
\uHom^{!}(\mathcal{M},\mathcal{N}) &\ \xrightarrow{\ \sim\ }\ \uHom_{!}(\mathcal{N},\mathcal{M})^{\circ}\\
f &\ \longmapsto\ \text{adjoint à gauche de } f\ ,
\end{aligned}
\]
\noindent\textup{\textbf{!}}\ \textup{\textbf{!}} et par suite la catégorie \textup{\textbf{opposée}} à $\uHom^{!}(\mathcal{M},\mathcal{N})$ est accessible.} % Verified against PDF leaf 25: a nonverbal marginal curved mark; see the correction register.

Ceci résulte du

\noindent\textbf{Lemme 4.14.10} (pour mémoire). {\itshape Soient $\mathcal{M}$ \edbr{[une]} catégorie pseudo-topos, $f:\mathcal{M}\longrightarrow\mathcal{N}$ un foncteur. Pour que $f$ admette un adjoint à gauche (resp. à droite), il faut et il suffit que $f$ commute aux petites $\varprojlim$ (resp. aux petites $\varinjlim$ ).}

% ---- from der-ch18_p146-150.tex
\printedpage{144}
% head: A. Grothendieck : Dérivateurs, Ch. XVIII, Édition des Universités de Montpellier II et Paris VII
% foot: Version 17 décembre 2025

\msp{page 380}
Ceci équivaut au cas particulier :

\textbf{Corollaire 4.14.11.} \textit{Soit $\mathcal{M}$ un paratopos, et $f$ un foncteur covariant (resp. contravariant) de $\mathcal{M}$ dans \textup{Ens}. Pour que $f$ soit représentable, il faut et il suffit que $f$ commute aux $\varprojlim$ .}

Ça a été vu -- le cas des contrafoncteurs étant le plus évident (cf. 4.1 -- il suffit que $\mathcal{M}$ soit un paratopos, et même que $\mathcal{M}$ admette une petite sous-catégorie pleine génératrice, cf. SGA 4 I 8.12.7 (pour l'énoncé dual). Dans les cas d'un foncteur covariant, la difficulté provient du fait que $\mathcal{M}$ n'admet pas forcément de petite sous-catégorie cogénératrice. Il faut utiliser l'existence d'une filtration cardinale (généralisée), cf. loc. cit. 8.3.12.) % [sic] Dans les cas -> Dans le cas, singular "un foncteur covariant"

Après tous ces beaux corollaires, on est mûr pour s'attaquer à la démonstration de 4.14.1 !

\dgrblock{dgrb1096}{Continuity does not supply a solution set}
\begin{dcomment}{examples}{p380:Après tous ces beaux corollaires}{Continuity does not supply a solution set}
Grothendieck wants the familiar passage between left and right adjoints to turn the preceding accessibility theorem into its dual. The formal passage is sound, but its input in 4.14.9--4.14.11 is too broad. A functor preserving all small limits need not have a left adjoint, even when its source and target are locally presentable. This is a genuine mathematical error in the printed statement, not a defect introduced by the transcription. The source's marginal mark beside 4.14.9 is retained without assigning it a meaning beyond what can be read on the page.

The obstruction was already available in the literature being cited: \citet[Expos\'e~I, Example~8.12.9]{SGA4I1972} attributes a continuous nonrepresentable functor on groups to H.~Bass. Grothendieck himself had recalled this example and the need for an extra cogeneration hypothesis at the beginning of the chapter \citep[ms.~50--52]{Grothendieck1990XVIII}. We give an explicit version so that the missing hypothesis has a visible effect. For every infinite cardinal $\kappa$ choose a simple group $S_\kappa$ of cardinality $\kappa$; one may take the group of even permutations of a set of that size having finite support. Simplicity follows from the finite alternating groups: a nonidentity element of a normal subgroup lies in a finite alternating subgroup on at least five letters, and its normal closure there contains a $3$-cycle; conjugation gives all $3$-cycles, which generate the finitary alternating group.

Define a covariant functor on groups by
\[
 B(G)=\prod_{\kappa\text{ infinite}}\Hom(S_\kappa,G).
\]
This notation denotes a set, despite its apparent indexing by all infinite cardinals. A nontrivial homomorphism from a simple group is injective, so the factor indexed by $\kappa>|G|$ is a singleton. One can therefore define $B(G)$ by the product over $\kappa\leq|G|$, with trivial coordinates understood elsewhere. Postcomposition defines $B$ on homomorphisms and respects these omitted coordinates.

To verify continuity, take any small diagram $(G_i)$ with limit $G$. There is a cardinal bounding the sizes of $G$ and of every $G_i$. Above that cardinal all coordinates are trivial. Below it, each functor $\Hom(S_\kappa,-)$ preserves limits, and a set-indexed product of such functors does so as well. Thus the canonical map
\[
 B(\varprojlim_iG_i)\longrightarrow\varprojlim_iB(G_i)
\]
is bijective. This checks all small shapes, including products with arbitrarily large small indexing sets.

To see why representability fails, suppose that an isomorphism $\Hom(H,-)\simeq B$ existed. By Yoneda it would be determined by an element $b\in B(H)$. Choose $\kappa>|H|$. The $\kappa$-coordinate of $b$ is trivial, and therefore every element obtained from $b$ by postcomposition has trivial $\kappa$-coordinate. But $B(S_\kappa)$ contains the family whose $\kappa$-coordinate is $\id_{S_\kappa}$ and whose other coordinates are trivial. That family cannot lie in the image of $\Hom(H,S_\kappa)$, a contradiction. Nor can $B$ have a left adjoint, since an adjunction would give $B(G)\simeq\Hom(L(1),G)$.

The correct version of 4.14.9 uses accessible continuous functors. Between locally presentable categories these are exactly the right adjoints, and taking left adjoints gives
\[
 \operatorname{Fun}^{R}(\mathcal M,\mathcal N)
 \simeq
 \operatorname{Fun}^{L}(\mathcal N,\mathcal M)^{op}.
\]
The reversal on natural transformations is the mate correspondence. Accessibility supplies the solution-set requirement in the adjoint functor theorem. For a precise modern statement, including a bound on the sizes of limits that need to be tested once accessibility has been fixed, see \citet[Theorems~2.9 and~2.11]{Tendas2022Continuity}.

The other half of 4.14.10 remains valid, even with its stated pseudo-topos hypothesis: a cocontinuous functor from a cocomplete locally small category with a small family generating by strict epimorphisms to a locally small category has a right adjoint. This is the dual of SGA~4~I~8.12.7--8.12.8. Locally presentable categories are a special case. Correspondingly, a contravariant set-valued functor taking colimits to limits is representable. There is no license to reverse this theorem without also reversing its size hypotheses. The reference printed as SGA~4~I~8.3.12 points toward 8.13.2, whose prorepresentability criterion explicitly assumes accessibility of the functor.
\end{dcomment}


\dgrblock{dgrb1097}{Testing universal arrows on small generators}
\begin{dcomment}{exposition}{p380:Après tous ces beaux corollaires}{Testing universal arrows on small generators}
The proof that now begins returns to the part of the argument that works without the representability detour. Grothendieck encodes a universal property by a comparison arrow and then tests whether that arrow is invertible. This makes the preservation condition a property of an object in an accessible arrow category, where the inverse-image theorem can be used.

For a prescribed cocone, its comparison arrow is
\[
 \varinjlim_{i\in I}F\varphi(i)\longrightarrow F\varphi(e_I).
\]
Its accessibility follows directly from pointwise colimits and interchange of colimits. For a prescribed cone, the comparison has the opposite orientation,
\[
 F\varphi(e'_I)\longrightarrow\varprojlim_{i\in I}F\varphi(i),
\]
and a cardinal bound enters. Fix a regular $\lambda$ larger than the number of arrows of $I$, and a small strong generating family $\mathcal A$ of $\lambda$-presentable objects in $\mathcal N$. If $(D_j)$ is a $\lambda$-filtered diagram of $I$-diagrams, then for $A\in\mathcal A$ the comparison between $\operatorname{colim}_j\lim_iD_j(i)$ and $\lim_i\operatorname{colim}_jD_j(i)$ becomes
\[
 \varinjlim_j\varprojlim_i\Hom(A,D_j(i))
 \longrightarrow
 \varprojlim_i\varinjlim_j\Hom(A,D_j(i)).
\]
It is a bijection because $\lambda$-filtered colimits of sets commute with limits having fewer than $\lambda$ arrows. The strong generators detect the original isomorphism. This explains both the size hypothesis and the correction of the printed $\varinjlim_I$ to $\varprojlim_I$ on manuscript page 383.

The isomorphism subcategory is especially transparent. The functor $X\mapsto\id_X$ is fully faithful from $\mathcal N$ to $\mathcal N^{[1]}$, preserves the available filtered colimits, and has as its replete essential image the invertible arrows. Thus this subcategory is accessible. For monomorphisms one uses the same small family differently: $u$ is monic precisely when every $\Hom(A,u)$ is injective. A generating family suffices here, whereas detection of isomorphisms requires strong generation.

The finite-set calculation left to the reader in 4.14.13 can be made without any representation theorem. Replace an injection by an actual inclusion $A\subset B$. Consider the poset of pairs of finite subsets $A_0\subset A$ and $B_0\subset B$ with $A_0\subset B_0$. It is directed under union, and the colimit of its arrows $A_0\hookrightarrow B_0$ is $A\hookrightarrow B$. Each finite arrow is finitely presentable in the arrow category: its morphisms to a filtered colimit factor at one stage, and the one commutative-square equation holds at a later stage. Filtered colimits of injections in sets remain injective. Hence these finite injections form a small family of finitely presentable generators of the injection category. For a bijection, take a finite subset of its source and its image in the target; the same argument gives the isomorphism assertion.

Finally, full faithfulness matters in 4.14.14. The formula
\[
 \Hom_{\operatorname{Ind}(\mathcal C)}
    (\varinjlim_iC_i,\varinjlim_jD_j)
 \simeq\varprojlim_i\varinjlim_j\Hom_{\mathcal C}(C_i,D_j)
\]
shows that a full inclusion $\mathcal D\subset\mathcal C$ induces a fully faithful functor on Ind-completions. Together with preservation of filtered colimits, this identifies its essential image with an accessible embedded full subcategory. Mere accessibility of an arbitrary functor would not imply the corresponding assertion about its full image.
\end{dcomment}


\msp{page 381}
\textsc{Démonstration de} 4.14.1. On sait que $\mathcal{M} = \uHom(\mathcal{C},\mathcal{N})$ est accessible (4.9.5, p. 220), et $\uHom^{\Lambda_2}_{\Lambda_1}(\mathcal{C},\mathcal{N})$ est une sous-catégorie strictement pleine de cette catégorie. Il suffit de prouver que c'est une sous-catégorie accessible dans $\mathcal{M}$ (4.13.1, p. 281). Pour tout $\xi = (I,\varphi) \in \Lambda_1$, on a un foncteur
\[
L_\xi \ : \ \mathcal{M} \ = \ \uHom(\mathcal{C},\mathcal{N}) \longrightarrow \uFl(\mathcal{N}) \ ,
\]
donné par
\[
L_\xi(f) \ = \ \text{flèche canonique } \varinjlim_{i\in I} f\varphi(i) \longrightarrow f\varphi(e_I) \ .
\]
Je dis que le foncteur $L_\xi$ est \textit{accessible.} En effet, soit $\pi_0$ un cardinal tel que $\mathcal{N}$ satisfasse à $\mathrm{L}_{\pi_0}$, il en est donc de même de $\mathcal{M} = \uHom(\mathcal{C},\mathcal{N})$ et de $\uFl(\mathcal{N}) \defeq \uHom(\boldsymbol{\Delta}_1,\mathcal{N})$, les limites en question se calculent argument par argument. Je dis que $L_\xi$ est $\pi_0$-accessible, i.e. commute aux $\varinjlim$ grandes devant $\pi_0$. Il suffit de voir qu'il en est ainsi des foncteurs composés avec les foncteurs
\[
\begin{tikzcd}
\uFl(\mathcal{N}) \arrow[r, shift left, "\sigma"] \arrow[r, shift right, "\rho"'] & \mathcal{N}
\end{tikzcd}
\]
``source'' et ``but'', i.e. des foncteurs
\[
f \longmapsto \varinjlim_{i\in I} f\varphi(i) \qquad \text{et} \qquad f \longmapsto f\varphi(e_I) \ .
\]

\msp{page 382}
Or le premier foncteur, $f \longmapsto \varinjlim\limits_{i\in I} f\varphi(i)$, est le composé de
\[
f \longmapsto f\circ\varphi_I \ : \ \uHom(\mathcal{C},\mathcal{N}) \longrightarrow \uHom(I,\mathcal{N})
\]
(lequel commute aux $\varinjlim$ qui sont représentables dans $\mathcal{N}$, et en particulier aux $\varinjlim$ grandes devant $\pi_0$) et du foncteur
\[
\varinjlim_{I} \ : \ \uHom(I,N) \longrightarrow \mathcal{N} \ , % [sic] \uHom(I,N) -> \uHom(I,\mathcal{N}), italic N printed instead of the script N used everywhere else
\]
qui lui aussi commute aux limites inductives représentables dans $\mathcal{N}$, et en particulier à celles grandes devant $\pi_0$. Quant au foncteur $f \longmapsto (f\circ\varphi)(e_I) = f(\varphi(e_I))$, il est clair qu'il commute également aux $\varinjlim$ représentables dans $\mathcal{N}$, donc il est $\pi_0$-accessible.

\printedpage{145}
% head: A. Grothendieck : Dérivateurs, Ch. XVIII, Édition des Universités de Montpellier II et Paris VII
% foot: Version 17 décembre 2025

Ceci posé, notons que
\[
\uHom_{\Lambda_1}(\mathcal{C},\mathcal{N}) \ = \ \bigcap_{\xi\in\Lambda_1} \uHom_{\{\xi\}}(\mathcal{C},\mathcal{N}) \ ,
\]
et que $\uHom_\xi(\mathcal{C},\mathcal{N})$ est l'image inverse, par le foncteur accessible $L_\xi$ précédent, de la sous-catégorie strictement pleine $\uIs(\mathcal{N})$ de $\uFl(\mathcal{N})$ formée des \textit{isomorphismes} dans $\mathcal{N}$. Si on prouve que cette dernière est accessible dans $\uFl(\mathcal{N})$, il s'ensuit que
\msp{page 383}
$\uHom_\xi(\mathcal{C},\mathcal{N})$ est accessible dans $\mathcal{M}$ (4.3.31, page 350). Donc, $\Phi$ étant petite, il s'ensuit que $\uHom_{\Lambda_1}(\mathcal{C},\mathcal{N})$ est accessible dans $\mathcal{M}$ (4.13.7, page 291).

On prouve de la même façon que $\uHom^{\Lambda_2}(\mathcal{C},\mathcal{N})$ est accessible dans $\mathcal{M} = \mathrm{Hom}(\mathcal{C},\mathcal{N})$. Ici, on doit considérer les foncteurs % [sic] \mathrm{Hom}(\mathcal{C},\mathcal{N}) -> \uHom(\mathcal{C},\mathcal{N}), printed without underline, elsewhere underlined
\[
L'_\xi \ : \ \mathcal{M} = \uHom(\mathcal{C},\mathcal{N}) \longrightarrow \uFl(\mathcal{N})
\]
définis par
\[
L'_\xi \ = \ \text{flèche canonique } f\varphi(e'_I) \longrightarrow \varprojlim_{i\in I} f\varphi(i) \ .
\]
Pour prouver que ce foncteur est accessible, on procède comme tantôt, en prenant cependant $\pi_0$ assez grand pour qu'on ait
\[
\left\{
\begin{array}{l}
I \in \Psi \ \Longrightarrow \ \cardred I \leqslant \pi_0. \\[1ex]
\mathcal{N}_{\pi_0} \text{ est génératrice dans } \mathcal{N}.
\end{array}
\right.
\]
Je dis qu'alors $L'_\xi$ est $\pi_0$-accessible : en effet, en procédant comme tantôt pour $L_\xi$, on est ramené à voir que
\[
\varinjlim_{I} \uHom(I,\mathcal{N}) \longrightarrow \mathcal{N}
\]
est $\pi$-accessible. Or cela est contenu dans SGA 4 I 9.8.

\msp{page 384}
Ainsi, on a prouvé que les deux sous-catégories strictement pleines
\[
\uHom_{\Lambda_1}(\mathcal{C},\mathcal{N}) \ , \qquad \uHom^{\Lambda_2}(\mathcal{C},\mathcal{N})
\]
de $\mathcal{M} \defeq \uHom(\mathcal{C},\mathcal{N})$ sont accessibles dans $\mathcal{M}$. Il en est donc de même de leur intersection (4.13.7)
\[
\uHom^{\Lambda_2}_{\Lambda_1}(\mathcal{C},\mathcal{N}) \ = \ \uHom_{\Lambda_1}(\mathcal{C},\mathcal{N}) \cap \uHom^{\Lambda_2}(\mathcal{C},\mathcal{N}) \ ,
\]
q.e.d.

Il reste à prouver le

\textbf{Lemme 4.14.12.} \textit{Soit $\mathcal{N}$ une catégorie accessible. Alors la catégorie $\uFl(\mathcal{N}) \defeq \uHom(\boldsymbol{\Delta}_1,\mathcal{N})$ est accessible (4.9.5), et la sous-catégorie strictement pleine $\uIs(\mathcal{N})$ formée des flèches inversibles dans $\mathcal{N}$ est accessible dans $\uFl(\mathcal{N})$. (Il en est de même de la sous-catégorie strictement pleine $\uMon(\mathcal{N})$ formée des monomorphismes dans $\mathcal{N}$.)}

\textsc{Démonstration}. Soit $\mathcal{C}$ une petite sous-catégorie strictement génératrice de $\mathcal{N}$, et considérons la petite famille des foncteurs
\[
h_x \ = \ (y \longmapsto \mathrm{Hom}(x,y)) \ : \ \mathcal{N} \longrightarrow \Ens \ ,
\]
\printedpage{146}
% head: A. Grothendieck : Dérivateurs, Ch. XVIII, Édition des Universités de Montpellier II et Paris VII
% foot: Version 17 décembre 2025
\msp{page 385}
d'où des foncteurs
\[
\uFl(h_\xi) \ : \ \uFl(\mathcal{N}) \longrightarrow \uFl(\Ens) \ . % [sic] h_\xi -> h_x, the functors defined just before are h_x
\]
Pour qu'une flèche $u : y \longrightarrow z$ dans $\mathcal{N}$ soit isomorphe (resp. monomorphe), il faut et il suffit que pour tout $x \in \Ob\mathcal{C}$, $h_x(u)$ le soit. Ainsi, $\uIs(\mathcal{N})$ (resp. $\uMon(\mathcal{N})$) apparaît comme l'intersection des images inverses, par les foncteurs $\uFl(h_\xi)$ ci-dessus, de la sous-catégorie $\uIs(\Ens)$ (resp. $\uMon(\Ens)$) de $\uFl(\Ens)$. Cela nous ramène au cas où $\mathcal{N} = (\Ens)$. Mais si $\mathcal{C}$ désigne la sous-catégorie strictement pleine de Ens formée des ensembles \textit{finis}, ceux-ci forment une famille strictement génératrice formée d'objets de \textit{présentation finie}, et comme Ens est stable par petites $\varinjlim$ , donc par $\varinjlim$ filtrantes, il s'ensuit qu'on a % [sic] h_\xi -> h_x, the functors defined just before are h_x
\[
\Ind(\mathcal{C}) \longrightarrow (\Ens) \ , \tag*{$(*)$}
\]
d'où aussitôt (SGA 4 I 8.8.5, p. 104)

\msp{page 386}
\[
\Ind(\uFl(\mathcal{C})) \xrightarrow{\ \sim\ } \uFl(\Ens) \ . \tag*{$(*)$}
\]
Donc le résultat d'accessibilité voulu résulte des deux lemmes suivants :

\textbf{Lemme 4.14.13.} \textit{L'équivalence $(*)$ précédente induit des équivalences}
\[
\begin{aligned}
\Ind(\uIs(\mathcal{C})) \quad &\xrightarrow{\ \sim\ } \quad \uIs(\Ens) \\
\Ind(\uMon(\mathcal{C})) \quad &\xrightarrow{\ \sim\ } \quad \uMon(\Ens)
\end{aligned}
\]
\textit{(où $\mathcal{C} = (\mathrm{Ensf})$ est la sous-catégorie pleine de \textup{(Ens)} formée des ensembles \textup{finis}.)}

\textbf{Lemme 4.14.14} (\footnote[166]{\textbf{N.B.} À équivalence près, on trouve les catégories accessibles $\mathcal{M}$ stables par petites $\varinjlim$ filtrantes, et telles que $\mathcal{M}_{\mathrm{pf}}$ (sous-catégorie des objets de présentation \textit{finie}) soit strictement génératrice.}). \textit{Soit $\mathcal{C}$ une catégorie équivalente à une petite catégorie. Alors :}
\begin{itemize}[nosep]
\item[\textit{a)}] \textit{$\mathcal{M} \defeq \Ind(\mathcal{C})$ est une catégorie accessible.}
\item[\textit{b)}] \textit{Pour toute sous-catégorie pleine $\mathcal{D}$ dans $\mathcal{C}$, l'image essentielle de $\Ind\mathcal{D}$ dans $\mathcal{M} = \Ind(\mathcal{C})$ est une sous-catégorie L-accessible dans $\mathcal{M}$.}
\end{itemize}
\textsc{Démonstration} des deux lemmes immédiate et ``laissée au lecteur''. Le b) dans 4.14.14 provient du fait que le foncteur $\Ind(\mathcal{D}) \longrightarrow \Ind(\mathcal{C})$ provenant d'un foncteur $\mathcal{D} \longrightarrow \mathcal{C}$ (ici un
\msp{page 387}
foncteur d'inclusion) est un L-foncteur (i.e. commute aux petites $\varinjlim$ filtrantes), et a fortiori est accessible (et même $\aleph_0$-accessible).

\textbf{4.14.15.} On ne peut s'empêcher de se demander si pour $\mathcal{N}$ accessible, la sous-catégorie strictement pleine $\uEp(\mathcal{N})$ dans $\uFl(\mathcal{N})$, formée des épimorphismes dans $\mathcal{N}$, est accessible dans $\mathcal{N}$. Il est immédiat que si $\mathcal{N}$ satisfait à $\mathrm{L}_{\pi_0}$, alors $\uEp(\mathcal{N})$ est stable dans $\uFl(\mathcal{N})$ par les $\varinjlim$ de type $\mathrm{L}_{\pi_0}$, a fortiori par facteurs directs. Elle aurait donc dû mal à ne pas être accessible dans $\uFl(\mathcal{N})$. Notons quand même : % [sic] dû mal -> du mal, "avoir du mal à"

\textbf{Corollaire 4.14.15.1} (de 4.14.13). \textit{Soit $\mathcal{N}$ une catégorie accessible stable par sommes amalgamées $y \amalg_x z$} (\footnote[167]{Je suis presque sûr que cette hypothèse est superflue. Comparer page 390.}). \textit{Alors $\uEp(\mathcal{N})$ est une sous-catégorie accessible dans $\uFl(\mathcal{N})$.}

\printedpage{147}
% head: A. Grothendieck : Dérivateurs, Ch. XVIII, Édition des Universités de Montpellier II et Paris VII
% foot: Version 17 décembre 2025

\msp{page 388}
Considérons en effet le foncteur
\[
\begin{aligned}
\gamma \ : \ \uFl(\mathcal{N}) \quad &\longrightarrow \quad \uFl(\mathcal{N}) \\
(x \longrightarrow y) \quad &\longmapsto \quad (y \amalg_x y \xrightarrow{\ \ \text{codiag}\ \ } y) \ .
\end{aligned}
\]
Ce foncteur commute aux types de $\varinjlim$ qui sont représentables dans $\mathcal{N}$, en particulier, il est accessible. D'autre part, $\uEp(\mathcal{N})$ est l'image inverse par $\gamma$ de la sous-catégorie strictement pleine $\uIs(\mathcal{N})$ de $\uFl(\mathcal{N})$. Donc le corollaire résulte du fait que par 4.14.13 cette dernière est accessible dans $\uFl(\mathcal{N})$.

\dgrblock{dgrb1098}{Epimorphisms as degenerate pushout squares}
\begin{dcomment}{examples}{p388:Donc le corollaire résulte du fait que}{Epimorphisms as degenerate pushout squares}
Grothendieck turns to epimorphisms because they expose the asymmetry between testing arrows into an object and testing arrows out of it. A small generating family detects monomorphisms. There need not be a corresponding small cogenerating family to detect epimorphisms. The self-pushout replaces all those outgoing tests by one universal comparison.

Let $u:X\to Y$, and suppose its self-pushout $P=Y\amalg_XY$ exists, with maps $i_1,i_2:Y\to P$ and codiagonal $d:P\to Y$. If $u$ is epic, the equality $i_1u=i_2u$ implies $i_1=i_2$. Since $di_1=di_2=\id_Y$, one has $i_1d=\id_P$ by the pushout universal property; hence $d$ is invertible. Conversely, if $d$ is invertible then $i_1=i_2=d^{-1}$. Given $a,b:Y\to Z$ with $au=bu$, the induced map $P\to Z$ shows that $a=b$. Thus $u$ is epic exactly when $d$ is invertible. This proves the inverse-image description in 4.14.15.1, and functoriality of pushouts makes its comparison functor accessible.

For sets, the self-pushout of $u$ consists of two copies of $Y$, with corresponding points identified exactly where they belong to the image of $u$. The codiagonal is bijective precisely when no point remains doubled, that is, when $u$ is surjective. A surjection $A\twoheadrightarrow B$ is a filtered colimit of finite surjections: index by finite subsets $A_0\subset A$ and take $A_0\twoheadrightarrow u(A_0)$. The finite subsets form a directed poset, their source union is $A$, and their image union is $B$. A square from a finite arrow factors through one stage, and its finitely many equations become valid at a later stage. Conversely, a finitely presentable surjection is a retract of one stage of this presentation, and therefore has finite source and target. Consequently the category of surjections of sets is finitely accessible with exactly the finite surjections as its finitely presentable objects.

The categorical formulation remains useful when epimorphism does not mean surjection of underlying sets. In commutative unital rings, $\mathbf Z\to\mathbf Q$ is epic: two ring maps out of $\mathbf Q$ agreeing on $\mathbf Z$ agree on every fraction, since the images of denominators have forced inverses. The corresponding codiagonal is the multiplication map
\[
 \mathbf Q\otimes_{\mathbf Z}\mathbf Q\longrightarrow\mathbf Q,
 \qquad r\otimes s\longmapsto rs.
\]
Its inverse sends $q$ to $q\otimes1$. In the tensor product, $1\otimes n=n\otimes1$ for every nonzero integer $n$; inverting this equality gives $1\otimes n^{-1}=n^{-1}\otimes1$, and therefore $r\otimes s=rs\otimes1$. This proves the claimed inverse directly, although $\mathbf Z\to\mathbf Q$ is far from surjective as a map of sets.

There are useful sharper results in algebraic settings. For a locally $\lambda$-presentable abelian category, \citet[Lemma~10.7]{Positselski2024Limits} proves that its epimorphism category is locally $\lambda$-presentable, with precisely the epimorphisms between $\lambda$-presentable objects as its $\lambda$-presentables. This fixed-rank statement uses the abelian hypotheses and does not settle the unrestricted question on the preceding page.
\end{dcomment}


\textbf{4.14.16.} Je présume que 4.14.1 reste valable sans supposer $\mathcal{N}$ stable par les $\varinjlim$ de type $\Phi$ et par les $\varprojlim$ de type $\Psi$. Si on a une essentiellement petite catégorie $I$, la question est de savoir si la sous-catégorie strictement pleine $\uHom_{\square}(I_\bullet,\mathcal{N})$ de $\uHom(I_\bullet,\mathcal{N})$, formée des foncteurs satisfaisant
\msp{page 389}
la condition a) de 4.14.1, est accessible dans $\uHom(I_\bullet,\mathcal{N})$, et de même pour $\uHom^{\square}({}_\bullet I,\mathcal{N})$. Pour ce deuxième cas, ça a l'air OK, car si $\mathcal{C}$ est une petite sous-catégorie strictement génératrice de $\mathcal{N}$, et si pour tout $x \in \Ob$ \texttt{[}$\mathcal{C}$\texttt{]} on considère encore le foncteur
\[
h_x \ : \ \mathcal{N} \longrightarrow \Ens \ ,
\]
d'où
\[
\bar{h}_x \ : \ \uHom({}_\bullet I,\mathcal{N}) \longrightarrow \uHom({}_\bullet I,\Ens) \ ,
\]
on voit que $\uHom^{\square}({}_\bullet I,\mathcal{N})$ est l'intersection des images inverses des sous-catégories $\uHom^{\square}({}_\bullet I,\Ens)$ de $\uHom({}_\bullet I,\Ens)$. Donc ces dernières étant accessibles dans $\uHom({}_\bullet I,\Ens)$, par 4.14.1 ((Ens) étant stable par petites $\varprojlim$ ), il s'ensuit bien que $\uHom^{\square}({}_\bullet I,\mathcal{N})$ l'est dans $\uHom({}_\bullet I,\mathcal{N})$, q.e.d.

Le cas de $\uHom_{\square}(I_\bullet,\mathcal{N})$ dans $\uHom(I_\bullet,\mathcal{N})$ semble plus délicat, dans le cas où
\msp{page 390}
\[
I = 
\begin{tikzcd}[sep=small]
 & {} \\
\bullet \arrow[ur] \arrow[dr] & \\
 & {}
\end{tikzcd}
\quad , \qquad \text{donc} \qquad
I = 
\begin{tikzcd}[sep=small]
 & {} \arrow[dr] & \\
\bullet \arrow[ur] \arrow[dr] & & {} \\
 & {} \arrow[ur] &
\end{tikzcd}
\quad ,
\]
si l'énoncé envisagé est vrai, on en conclura que $\uEp(\mathcal{N})$ est accessible dans $\uFl(\mathcal{N})$, pour toute catégorie accessible, car c'est l'image inverse de $\uHom_{\square}(I_\bullet,\mathcal{N})$ par le foncteur canonique (accessible)
\[
\begin{aligned}
\uFl(\mathcal{N}) \quad &\longrightarrow \quad \uHom(I_\bullet,\mathcal{N}) \\
(x \xrightarrow{\ u\ } y) \quad &\longmapsto \quad
\left(
\begin{tikzcd}[sep=small, ampersand replacement=\&]
 \& y \arrow[dr, "\mathrm{id}_y"] \& \\
x \arrow[ur, "u"] \arrow[dr, "u"'] \& \& y \\
 \& y \arrow[ur, "\mathrm{id}_y"'] \&
\end{tikzcd}
\right) \ .
\end{aligned}
\]

\printedpage{148}
% head: A. Grothendieck : Dérivateurs, Ch. XVIII, Édition des Universités de Montpellier II et Paris VII
% foot: Version 17 décembre 2025

C'est dire que la généralisation envisagée de 4.14.1 n'est pas oiseuse, mais est sans doute utile.

\msp{page 391}
\textbf{Théorème 4.14.17} (\footnote[168]{Cf. généralisation dans \texttt{[le]} théorème page 422, th. 4.14.26.}). \textit{Soit $p : \mathcal{E} \longrightarrow \mathcal{C}$ un foncteur cofibrant ou fibrant, avec $\mathcal{C}$ essentiellement petite (i.e. équivalente à une petite catégorie), et tel que les catégories fibres soient accessibles, et les foncteurs cochangement de base (cas cofibrant) ou de changement de base (cas fibrant) associés aux flèches $u$ de $\mathcal{C}$ soient accessibles. Alors la catégorie $\uGamma(\mathcal{E}/\mathcal{C})$ des sections de $\mathcal{E}$ sur $\mathcal{C}$ est accessible, et pour tout $s \in \Ob\mathcal{C}$, le foncteur ``évaluation en $s$'' $\varepsilon_s : \uGamma(\mathcal{E}/\mathcal{C}) \longrightarrow \mathcal{E}_s$ est accessible.}

\dgrblock{dgrb1099}{Why the epimorphism question reaches set theory}
\begin{dcomment}{development}{p390}{Why the epimorphism question reaches set theory}
Grothendieck's proposed generalization in 4.14.16 asks whether actual universal cones can be recognized accessibly without assuming that all cones of the same shape have universal objects. For limits, his representable test on page 389 proves the assertion. For colimits, the degenerate square on page 390 puts the epimorphism question inside the proposed generalization. The following consequence explains why closure under filtered colimits cannot by itself finish the proof.

\textbf{Proposition (epimorphism accessibility forces co-wellpoweredness).}\ours\ 
Let $\mathcal K$ be accessible. If the full subcategory $\operatorname{Epi}(\mathcal K)\subset\mathcal K^{[1]}$ is accessible and accessibly embedded, then $\mathcal K$ is co-wellpowered: for each object $X$, its epimorphic quotients form a set up to isomorphism over $X$.

\emph{Proof.} The domain functor $d:\operatorname{Epi}(\mathcal K)\to\mathcal K$ is accessible. Form its iso-comma category with the functor $\mathbf1\to\mathcal K$ selecting $X$. The latter functor is accessible, because a constant diagram indexed by a nonempty filtered category has its chosen value as colimit. By the iso-comma theorem 4.14.21, this category is accessible. Its objects are epimorphisms $e:A\to B$ equipped with an isomorphism $X\simeq A$, and it is equivalent to the full subcategory $\mathcal Q_X$ of $X/\mathcal K$ on the epimorphisms out of $X$.

Between any two objects of $\mathcal Q_X$ there is at most one morphism: if $h e=k e$ and $e$ is epic, then $h=k$. Thus $\mathcal Q_X$ is a preorder. We check directly that every accessible preorder is essentially small. Choose a set $\mathcal G$ of presentable generators. An object $q$ determines the subset
\[
 D_q=\{g\in\mathcal G:g\longrightarrow q\}.
\]
If $D_q\subset D_r$, express $q$ as a filtered colimit of objects of $\mathcal G$. Every object in that presentation has a map to $r$, and these maps are automatically compatible because there is at most one. The colimit supplies $q\to r$. Equality of the two subsets therefore implies arrows both ways, hence an isomorphism. Isomorphism classes inject into the power set of $\mathcal G$, which is a set. Applying this to $\mathcal Q_X$ proves the proposition.\hfill$\square$

The proof is supplied here as a direct consequence of the chapter's own machinery; the marker records our presentation and does not assert priority. We checked the literature on accessible epimorphism categories and on co-wellpoweredness. The argument is closely related to the use of the quotient preorder in \citet[Corollary~7.6]{BagariaCasacubertaMathiasRosicky2015}; we did not locate this exact iso-comma formulation as a separate named result.

There is a definite foundational consequence. The discussion following that corollary, on published pages 574--575, records that co-wellpoweredness of every accessible category implies a proper class of measurable cardinals, and that a proper class of strongly compact cardinals suffices. Our proposition therefore gives the same necessary lower bound for the universal assertion that all accessible categories have accessible epimorphism categories. In particular, that universal assertion cannot be treated as an unconditional theorem of ZFC. This is an obstruction, not a determination of its precise consistency strength. The cited discussion also connects the strong-compact hypothesis to detectability of colimits through Makkai--Par\'e's Theorem~6.3.8(i).

The limit-cone argument has no such reversal of tests. For a small diagram $D:I\to\mathcal K$ with a proposed cone $x\to D$, choose a small dense subcategory $\mathcal A$ of presentables. The cone is limiting exactly when
\[
 \Hom(A,x)\longrightarrow\varprojlim_{i\in I}\Hom(A,D(i))
\]
is bijective for every $A\in\mathcal A$. Density makes this criterion sufficient, and each test is accessible at a sufficiently large common rank. Hence the category of existing limit cones is an accessible intersection of inverse images of the corresponding set-valued cone category. The source thus resolves its limit question while leaving a substantially different colimit question.
\end{dcomment}


\textbf{Corollaire 4.14.18.} \textit{Soit $\pi_0$ un cardinal infini tel que les catégories fibres satisfassent la condition $L_{\pi_0}$, et que les foncteurs cochangement de base (resp. changement de base) entre icelles satisfassent également $L_{\pi_0}$. (Un tel $\pi_0$ existe.)}
\begin{itemize}[nosep]
\item[\circled{a}] \textit{Alors $\mathcal{M} = \uGamma(\mathcal{E}/\mathcal{C})$ satisfait la condition $\underline{\mathrm{L}}_{\pi_0}$, et les foncteurs $\varepsilon_s$ aussi, i.e. les $\varinjlim$ grandes devant $\pi_0$ dans $\uGamma(\mathcal{E}/\mathcal{C})$ se calculent argument par argument.}
\item[\circled{b}] \textit{Soit un cardinal $\pi \geqslant \sup(\pi_0,\cardred\mathcal{C})$. Pour qu'un objet $f$ de $\mathcal{M}$ appartienne à $\mathcal{M}_\pi$, il \textup{suffit} que $\forall\ s \in \Ob C$, $f(s) \defeq \varepsilon_s(x)$ \textup{\texttt{[plutôt} $\varepsilon_s(f)$\texttt{]}} soit dans $\mathcal{E}_{s\,\pi}$. (Comparer 4.5.8, page 119.)} % [sic] \Ob C -> \Ob\mathcal{C}, italic C printed instead of script C
\item[\circled{c}] \textit{Supposons $pi_0$ choisi de façon} % [sic] pi_0 -> \pi_0, the letters "pi" printed instead of the Greek letter
\msp{page 392}

\textit{à satisfaire les conditions supplémentaires suivantes : $\pi_0$ est utile pour $\mathcal{M} = \uGamma(\mathcal{E}/\mathcal{C})$, et pour chacune des catégories fibres $\mathcal{E}_s$ ($s \in \Ob\mathcal{C}$) (un tel $\pi_0$ existe). Soit}
\[
\pi_1 \ = \ \sup(\pi_0, \cardred\mathcal{C}, \cardred\mathcal{M}_{\pi_0}) \ .
\]
\textit{Alors pour tout cardinal $\pi \geqslant \pi_1$ et tel que $\pi^{\pi_0} = \pi$, on a}
\[
\mathcal{M}_\pi \ = \ \left\{
\begin{array}{l}
\textit{sous-catégorie strictement pleine de $\mathcal{M}$ formée des $f \in$} \\
\textit{$\Ob\mathcal{M}$ telles que $f(s) \in \Ob\mathcal{E}_{s\,\pi}$ $\forall\ s \in S$ \textup{\texttt{[plutôt}} $s \in$} \\
\textit{$\Ob C$\textup{\texttt{]}}.} % [sic] \Ob C -> \Ob\mathcal{C}, italic C printed instead of script C
\end{array}
\right.
\]
\textit{Si de plus $\pi$ est choisi tel qu'on ait}
\[
\pi \ \geqslant \ \sup_{s\in\Ob\mathcal{C}} \cardred\mathcal{E}_{s\,\pi_0} \ ,
\]
\textit{alors la sous-catégorie strictement pleine $\mathcal{E}(\pi)$ de $\mathcal{E}$ telle que $\Ob\mathcal{E}(\pi) = \bigcup_{s\in\Ob\mathcal{C}} \Ob\mathcal{E}_{s\,\pi}$ est une sous-catégorie $\mathcal{C}$-cofibrée (resp. $\mathcal{C}$-fibrée) de $\mathcal{E}$, et on a}
\[
\uGamma(\mathcal{E}/\mathcal{C})_\pi \xleftarrow{\ \sim\ } \uGamma(\mathcal{E}(\pi)/C), % [sic] C -> \mathcal{C}, italic C printed instead of script C
\]
\textit{donc ($\pi$ étant utile pour $\mathcal{E}$) on a une équivalence de catégories}
\[
\uGamma(\mathcal{E}/\mathcal{C}) \xleftarrow{\ \sim\ } \Ind_\pi(\uGamma(\mathcal{E}(\pi)/\mathcal{C})) \ .
\]
\end{itemize}

% ---- from der-ch18_p151-155.tex
\printedpage{149}
% head: A. Grothendieck : Dérivateurs, Ch. XVIII, Édition des Universités de Montpellier II et Paris VII

\msp{page 393}
Toutes ces assertions se prouvent facilement dans l’ordre où elles sont données, en utilisant 4.14.17 et les résultats généraux déjà prouvés. Bien sûr, on peut commencer par remplacer $\mathcal{C}$ par une sous-catégorie réduite $\mathcal{C}_0$ équivalente et $\mathcal{E}$ par $\mathcal{E}_0 = \mathcal{E}|\mathcal{C}$ \texttt{[plutôt} $\mathcal{E}|\mathcal{C}_0$\texttt{]}, de façon à se réduire (dans b) et c)) au cas où $\operatorname{card}\mathcal{C} \leqslant \pi$.

\dgrblock{dgrb1100}{Sections and the two directions of transport}
\begin{dcomment}{development}{p392}{Sections and the two directions of transport}
Grothendieck's section theorem assembles accessible categories over a small base. The objects are families in the fibres together with compatible transition arrows, so accessibility should follow by adjoining structure and imposing equations. The main theorem does hold. The presentability estimate in 4.14.18(b), however, needs a distinction between the two directions of transport that is absent from its printed statement.

First consider a fibration, and fix a cleavage. A section $x$ consists of objects $x_s$ and arrows $\alpha_a:x_s\to a^*x_t$ for $a:s\to t$, satisfying the unit and composition equations with the coherence isomorphisms of the cleavage. For sections $x,y$, their morphisms are the equalizer
\[
 \Hom_{\Gamma}(x,y)\longrightarrow
 \prod_{s\in\Ob\mathcal C}\Hom_{\mathcal E_s}(x_s,y_s)
 \rightrightarrows
 \prod_{a:s\to t}\Hom_{\mathcal E_s}(x_s,a^*y_t).
\]
The two images of a family $(h_s)$ at $a$ are $a^*(h_t)\alpha_a$ and $\alpha_a^y h_s$. Take a regular $\lambda$ exceeding the number of arrows of the base, suppose the transports preserve $\lambda$-filtered colimits, and suppose every $x_s$ is $\lambda$-presentable. All terms on the right preserve those filtered colimits in the variable $y$, and the products and equalizer have fewer than $\lambda$ entries. Interchange in sets proves that $x$ is $\lambda$-presentable. This establishes the fibration half of 4.14.18(b).

For an opfibration, the second product instead has factors
\[
 \Hom_{\mathcal E_t}(a_*x_s,y_t).
\]
Presentability of $x_s$ does not guarantee presentability of $a_*x_s$ merely because $a_*$ is accessible. That missing implication cannot be repaired by the cardinality of the base.

\textbf{Counterexample (pointwise presentability need not suffice).}\ours\ 
For every infinite regular $\lambda$, take the base $[1]$, both fibres equal to $\mathbf{Set}$, and the transition functor constant at a set $B$ with $|B|=\lambda$. This functor preserves every filtered colimit: filtered categories are nonempty and connected. Its opfibration has section category
\[
 \mathbf{Set}\times(B/\mathbf{Set}),
\]
where a section is $(X,B\to Y)$. The section $x=(\varnothing,B\to1)$ has finite, hence $\lambda$-presentable, values in both fibres.

To test $x$ itself, index by the poset of subsets $S\subset B$ of cardinality less than $\lambda$. This poset is $\lambda$-filtered, since the union of fewer than $\lambda$ such subsets still has cardinality less than $\lambda$. Let $B/S$ denote the quotient identifying all elements of $S$ and leaving the other elements separate. The sections
\[
 x_S=(\varnothing,B\longrightarrow B/S)
\]
form a $\lambda$-filtered diagram. Its colimit is $x$, because every pair of elements of $B$ is identified at some stage. No map $x\to x_S$ exists over $B$: such a map would make $B\to B/S$ constant, although $S$ has size less than $|B|$. Thus
\[
 \varinjlim_S\Hom(x,x_S)=\varnothing,
 \qquad
 \Hom(x,\varinjlim_Sx_S)=\Hom(x,x)\ne\varnothing.
\]
This disproves the opfibration half of 4.14.18(b) as written. For Grothendieck's indexing choose $\lambda=\pi^+$; the base is finite and the constant transition satisfies every condition $L_{\pi_0}$ in the statement.

The corrected sufficient condition is explicit: require each transported object $a_*x_s$, as well as each $x_s$, to be $\lambda$-presentable. Choosing a common strong accessibility rank for the small family of transports ensures this condition for all pointwise presentable sections. After also enlarging the rank so that every evaluation preserves presentables, one obtains an eventual pointwise characterization. The printed numerical threshold in 4.14.18(c), however, does not itself ensure this enlargement.

\textbf{Counterexample (the section skeleton does not bound evaluation).}\ours\ 
Let $\mathcal P(S)$ be the poset of subsets of a set $S$, ordered by inclusion. It is locally finitely presentable: unions compute colimits, finite subsets are finitely presentable, and every subset is the directed union of its finite subsets. Consider the $[1]$-opfibration given by the constant functor $\mathbf1\to\mathcal P(S)$ with value $S$. Its sections are pairs consisting of the unique object of $\mathbf1$ and an inclusion $S\subset Y\subset S$. There is therefore exactly one section, and its structure arrow is an isomorphism. Both the section category and its cocartesian subcategory are $\mathbf1$.

Choose $\pi_0=\aleph_0$, $\pi=2^{\aleph_0}$, and $|S|>\pi$. The fibres and section category admit presentations at $\pi_0$, and transport preserves every filtered colimit. The printed bound is $\pi_1=\aleph_0$, since the section skeleton has one object and one arrow, and $\pi^{\pi_0}=\pi$. The unique section is presentable at every rank. Its value $S$ is not $\pi^+$-presentable in $\mathcal P(S)$: the subsets of $S$ of cardinality at most $\pi$ form a $\pi^+$-filtered poset, their union is $S$, and no stage admits an inclusion from $S$. This refutes the first pointwise equality in 4.14.18(c), and also the pointwise equality in 4.14.20(c). The later additional bound involving the fibre skeletons excludes this particular example: the skeleton of $(\mathcal P(S))_{\pi_0}$ already contains all singletons.

These two examples were obtained by following the two missing size implications separately. We checked the fixed-rank results and hypotheses in \citet{Positselski2024Limits}; we did not locate these exact counterexamples there or in the other sources consulted. The markers record the present constructions, without a claim of priority.

Cocartesian sections behave better even at a fixed rank. Their structure arrows identify $a_*x_s$ with $x_t$, so pointwise presentability already gives all the source objects needed by the equalizer. Thus the pointwise sufficient condition in 4.14.20(b) remains true for cocartesian sections, independently of the false unrestricted opfibration assertion. Cartesian sections likewise satisfy the fibration argument. Their filtered colimits are pointwise whenever transport preserves those colimits, because the induced structure comparisons remain isomorphisms. This accounts for the accessibility and reconstruction claims without conflating the two variance conventions.
\end{dcomment}


\textbf{Corollaire 4.14.19.} \textit{Sous les conditions de 4.14.17, la sous-catégorie $\uGamma_{\cocart}(\mathcal{E}/\mathcal{C})$ de $\mathcal{M} = \uGamma(\mathcal{E}/\mathcal{C})$ (cas $p$ cofibrant), resp. la sous-catégorie $\uGamma_{\cart}(\mathcal{E}/\mathcal{C})$ (cas $p$ fibrant) est accessible dans $\mathcal{M}$, donc accessible et le foncteur d’inclusion} \texttt{[est]} \textit{accessible.}

(\textbf{N.B.} Ces deux catégories sont strictement pleines dans $\uGamma(\mathcal{E}/\mathcal{C})$, donc l’assertion qu’elles soient accessibles dans celle-ci \textit{équivaut} à la dernière assertion.)

\textsc{Démonstration de} 4.14.19. On peut supposer $\mathcal{C}$ réduite, donc petite. Considérons le cas où $p$ est cofibrant. Pour toute flèche $\lambda : s \longrightarrow t$ dans $\mathcal{C}$, considérons le foncteur
\[
\varphi_\lambda \ : \ \mathcal{M} \longrightarrow \uFl(\mathcal{E}_t)
\]
défini par
\msp{page 394}
\[
\varphi_\lambda(f) \ = \ (\lambda_*(f(s)) \xrightarrow{f(\lambda)^\natural} f(t)) \ .
\]
C’est un foncteur accessible, car son composé avec les foncteurs source et but sont les foncteurs $\lambda_*\varepsilon_s$ et $\varepsilon_t$, donc accessibles, $\varepsilon_s$, $\varepsilon_t$ et $\lambda_*$ étant accessibles. Considérons la sous-catégorie strictement pleine
\[
\uIs(\mathcal{E}_t) \ \subset \ \uFl(\mathcal{E}_t) \ ,
\]
on sait qu’elle est accessible dans $\uFl(\mathcal{E}_t)$ (4.14.12). Donc son image inverse par $\varphi_\lambda$ est accessible dans $\mathcal{M}$, et par suite aussi
\[
\uGamma_{\cart}(\mathcal{E}/\mathcal{C}) \ \overset{\text{déf}}{=} \ \bigcap_{\lambda\in\Ob\mathcal{C}} \varphi_\lambda^{-1}(\uIs(\mathcal{E}_t))
\]
% [sic] \uGamma_{\cart} -> \uGamma_{\cocart}, the case treated is p cofibrant
% [sic] \lambda\in\Ob\mathcal{C} -> \lambda\in\Fl\mathcal{C}, \lambda runs over arrows of \mathcal{C}
comme intersection d’une \textit{petite} famille de sous-catégories accessibles dans $\mathcal{M}$, q.e.d. Le cas où $p$ est fibrant se traite de façon similaire, en introduisant les foncteurs
\[
\psi_\lambda \ : \ \mathcal{M} = \uGamma(\mathcal{E}/\mathcal{C}) \longrightarrow \uFl(\mathcal{E}_s)
\]
donnés par
\[
\psi_\lambda(f) \ = \ (f(s) \xrightarrow{f(\lambda)^\flat} \lambda^* f(t)) \ .
\]

\msp{page 395}
\textbf{Corollaire 4.14.20.} \textit{Soit $\pi_0$ un cardinal infini choisi comme dans 4.14.18 a).}
\begin{itemize}[leftmargin=2.5em]
\item[\circled{a}] \textit{Alors $\mathcal{M}_{\cc} = \uGamma_{\cocart}(\mathcal{E}/\mathcal{C})$ (resp. $\mathcal{M}_{\cc} = \uGamma_{\cart}$) est $L_{\pi_0}$-stable dans $\mathcal{M} = \uGamma(\mathcal{E}/\mathcal{C})$, ou ce qui revient au même (comme elle est strictement pleine), elle satisfait $L_{\pi_0}$, et le foncteur d’inclusion satisfait à $L_{\pi_0}$. Les $\varinjlim$ grandes devant $\pi_0$ dans cette catégorie se calculent argument par argument.}
\item[\circled{b}] \textit{Pour tout cardinal $\pi \geqslant \pi_0$, tout objet de $\mathcal{M}_{\cc}$ qui est $\pi$-accessible dans $\mathcal{M}$ l’est (tautologiquement) dans $\mathcal{M}_{\cc}$. Si $\pi \geqslant \sup(\pi_0, \cardred\mathcal{C})$, alors (par 4.14.18 b) tout objet $f$ de $\mathcal{M}_{\cc}$ tel que $f(s) \in \Ob\mathcal{E}_{s\,\pi}$ $\forall\ s \in S$} \texttt{[plutôt} $s \in \Ob C$\texttt{]} \textit{est dans $\mathcal{M}_\pi$} \textup{et} \textit{dans $\mathcal{M}_{\cc\,\pi}$.}
% foot: Version 17 décembre 2025

\printedpage{150}
% head: A. Grothendieck : Dérivateurs, Ch. XVIII, Édition des Universités de Montpellier II et Paris VII

\item[\circled{c}] \textit{Supposons $\pi_0$ choisi comme dans 4.14.18 c, et que de plus $\mathcal{M}_{\cc}$ soit $\pi_0$-accessible dans $\mathcal{M}$. (Un tel $\pi_0$ existe.) On a donc, pour $\pi \geqslant \pi_1 = \sup(\pi_0, \cardred\mathcal{C}, \cardred\mathcal{M}_\pi)$ tel que $\pi^{\pi_0} = \pi$,}
\[
(\mathcal{M}_{\cc})_\pi \ = \ \mathcal{M}_\pi{\cap}\mathcal{M}_{\cc} \ = \quad
\begin{array}{l}
\textit{Sous-catégorie strictement pleine de $\mathcal{M}_{\cc}$ formée des $f$}\\
\textit{telles que $f(s) \in \Ob\mathcal{E}_{s\,\pi}$ $\forall\ s \in \Ob\mathcal{C}$.}
\end{array}
\]

\msp{page 396}
\textit{Si de plus $\pi$ est choisi tel qu’on ait $\pi \geqslant \sup_{s\in\Ob\mathcal{C}} \cardred\mathcal{E}_{s\,\pi_0}$, alors}
\[
\left\{
\begin{array}{rrcl}
 & \mathcal{M}_{\cc\,\pi} & \xleftarrow{\ \sim\ } & \uGamma_{\cocart}(\mathcal{E}(\pi)/\mathcal{C})\\[1ex]
\textit{resp.} & \mathcal{M}_{\cc\,\pi} & \xleftarrow{\ \sim\ } & \uGamma_{\cart}(\mathcal{E}(\pi)/\mathcal{C}) \ ,
\end{array}
\right.
\]
\textit{donc ($\pi$ étant utile pour $\mathcal{M}_{\cc}$) une équivalence de catégories}
\[
\begin{array}{rrcl}
 & \uGamma_{\cocart}(\mathcal{E}/\mathcal{C}) & \xleftarrow{\ \sim\ } & \mathrm{Ind}_\pi(\uGamma_{\cocart}(\mathcal{E}(\pi)/\mathcal{C}))\\
\textit{resp.}\quad & \uGamma_{\cart}(\mathcal{E}/\mathcal{C}) & \xleftarrow{\ \sim\ } & \mathrm{Ind}_\pi(\uGamma_{\cart}(\mathcal{E}(\pi)/\mathcal{C})).
\end{array}
\]
\end{itemize}

\textsc{Démonstration}. Tout est clair dans l’ordre où c’est donné en utilisant les deux corollaires précédents.

\textbf{Corollaire 4.14.21.} \textit{Considérons un diagramme de foncteurs}
\[
\begin{tikzcd}[column sep=small, row sep=small]
\mathcal{E}_1 \arrow[dr, "p_1"] & & \mathcal{E}_2 \arrow[dl, "p_2"'] \\
& \mathcal{E}_0 \ , &
\end{tikzcd}
\tag*{$(*)$}
\]
\textit{les catégories et les foncteurs étant accessibles.}

\msp{page 397}
\begin{itemize}[leftmargin=2.5em]
\item[\textit{a)}] \textit{Le produit 2-fibré}
\[
\mathcal{M}_0 \ = \ \mathcal{E}_1 \overset{2}{\times}_{\mathcal{E}_0} \mathcal{E}_2
\]
\textit{est accessible, et les foncteurs canoniques $\mathcal{M} \longrightarrow \mathcal{E}_1$, $\mathcal{M} \longrightarrow \mathcal{E}_2$ sont accessibles.} % [sic] \mathcal{M} \longrightarrow \mathcal{E}_1, \mathcal{M} \longrightarrow \mathcal{E}_2 -> \mathcal{M}_0 \longrightarrow \mathcal{E}_1, \mathcal{M}_0 \longrightarrow \mathcal{E}_2, assertion a) concerns \mathcal{M}_0
\item[\textit{b)}] (\footnote[169]{Ceci est contenu dans 4.14.22.1 ci-dessus (p. 406), appliqué à $\Pi = \mathcal{E}_1{\times}\mathcal{E}_2$, $\Pi' = \mathcal{E}_0$, et $F, G : \Pi \rightrightarrows \Pi'$ donnés par $F = p_1\circ\pr_1$ et $G = p_2\circ\pr_2$. D’autre part, a) et c) résultent facilement de b).}). \textit{La catégorie $\mathcal{M}$ des triples $(x_1,x_2,u)$ avec $x_1 \in \Ob\mathcal{E}_1$, $x_2 \in \Ob\mathcal{E}_2$, et}
\[
u \ : \ p_1(x_1) \ \longrightarrow \ p_2(x_2) \ ,
\]
\textit{est accessible, et les foncteurs canoniques $\mathcal{M} \longrightarrow \mathcal{E}_1$, $\mathcal{M} \longrightarrow \mathcal{E}_2$ sont accessibles.}
\item[\textit{c)}] \textit{La sous-catégorie strictement pleine $\mathcal{M}_0$ de $\mathcal{M}$ est accessible dans $\mathcal{M}$.}
\end{itemize}

\noindent\textsc{Démonstration}. Soit $\mathcal{C}$ la catégorie
\[
\begin{tikzcd}[column sep=small]
1 \arrow[dr] & & 2 \arrow[dl] \\
& 0 \ , &
\end{tikzcd}
\]
et $\mathcal{E}$ la catégorie cofibrée sur $\mathcal{C}$ définie par le diagramme $(*)$. On est sous les conditions de 4.14.17. La catégorie $\Gamma = \uGamma(\mathcal{E}/\mathcal{C})$ s’interprète
% foot: Version 17 décembre 2025
\printedpage{151}
% head: A. Grothendieck : Dérivateurs, Ch. XVIII, Édition des Universités de Montpellier II et Paris VII
\msp{page 398}
comme la catégorie des quintuplets $(x_0,x_1,x_2,u_1,u_2)$, avec $x_i \in \Ob\mathcal{E}$ pour $i \in \{0,1,2\}$, et $u_1, u_2 \in \Fl\mathcal{E}_0$,
\[
\begin{tikzcd}[column sep=small, row sep=small]
p_1(x_1) \arrow[dr, "u_1"] & & p_2(x_2) \arrow[dl, "u_2"'] \\
& x_0 \ . &
\end{tikzcd}
\tag*{$(**)$}
\]
En vertu de 4.14.18, cette catégorie $\uGamma$ est accessible. La catégorie $\mathcal{M}$ est équivalente à la sous-catégorie strictement pleine $\Gamma_0$ de $\Gamma$, formée des quintuplets pour lesquels on a $u_2 \in \uIs\mathcal{E}_0$. On a en effet des foncteurs quasi-inverses l’un de l’autre % [sic] \Gamma_0 -> \Gamma_1, the subcategory is called \Gamma_1 in what follows
\[
\begin{tikzcd}
\mathcal{M} \arrow[r, shift left, "\alpha"] & \Gamma_1 \arrow[l, shift left, "\beta"]
\end{tikzcd}
\]
donnés par
\[
\begin{aligned}
\alpha(x_1,x_2,u) \ &= \ (p_2(x_2),x_1,x_2,u,\id_{p_2(x_2)})\\
\beta(x_0,x_1,x_2,u_1,u_2) \ &= \ (x_1,x_2,u_2^{-1}u_1) \ ,
\end{aligned}
\]
on a alors
\[
\begin{aligned}
\beta\alpha \quad &= \quad \id_{\mathcal{M}}\\
\alpha\beta \quad &\xrightarrow{\ \sim\ } \quad \id_{\Gamma_1} \ .
\end{aligned}
\]
D’autre part, $\Gamma_1$ est évidemment accessible dans $\mathcal{M}$, comme image inverse de $\uIs(\mathcal{E}_0)$ (qui % [sic] accessible dans \mathcal{M} -> accessible dans \Gamma, \Gamma_1 is a subcategory of \Gamma
\msp{page 399}
est accessible dans $\uFl(\mathcal{E}_0)$ par le foncteur accessible
\[
\uGamma \ \longrightarrow \ \uFl(\mathcal{E}_0)
\]
donné par $(x_1, x_1, x_2, u_1, u_2) \longmapsto u_2$. Cela prouve que $\mathcal{M}$, qui est équivalente à $\Gamma_1$, est également accessible, et l’accessibilité du foncteur canonique $\mathcal{M} \longrightarrow \mathcal{E}_i$ ($i \in \{1,2\}$) résulte de celle de $\Gamma \longrightarrow \mathcal{E}_i$, impliquant celle du composé $\Gamma_1 \hookrightarrow \Gamma \longrightarrow \mathcal{E}_i$. Cela prouve donc b). Pour prouver a), il suffit de voir que $\mathcal{M}_0$ est accessible dans $\mathcal{M}$, or c’est encore l’image inverse de $\uIs(\mathcal{E}_0) \subset \uFl(\mathcal{E}_0)$ par le foncteur accessible canonique % [sic] (x_1, x_1, x_2, u_1, u_2) -> (x_0, x_1, x_2, u_1, u_2), the quintuplets begin with x_0
\[
\mathcal{M} \ \longrightarrow \ \uFl(\mathcal{E}_0) \ ,
\]
donc $\mathcal{M}$ est bien accessible dans $\mathcal{M}$, comme image inverse d’une sous-catégorie accessible dans $\uFl(\mathcal{E})$. On peut d’ailleurs aussi interpréter $\mathcal{M}$, à équivalence de catégories près, comme % [sic] donc \mathcal{M} est bien accessible -> donc \mathcal{M}_0 est bien accessible, \mathcal{M}_0 is the subcategory in question
% [sic] \uFl(\mathcal{E}) -> \uFl(\mathcal{E}_0), as in the preceding display
% [sic] interpréter \mathcal{M} -> interpréter \mathcal{M}_0, as in the following display
\msp{page 400}
\[
\mathcal{M}_0 \ \simeq \ \uGamma_{\cocart}(\mathcal{E}/\mathcal{C}) \ ,
\]
les foncteurs canoniques $\mathcal{M}_0 \longrightarrow \mathcal{E}_i$ ($i \in \{0,1\}$) s’interprétant comme des foncteurs évaluation, ce qui donne a) directement, comme cas particulier de 4.14.19. % [sic] i \in \{0,1\} -> i \in \{1,2\}, the canonical functors of a) go to \mathcal{E}_1, \mathcal{E}_2
% foot: Version 17 décembre 2025

\dgrblock{dgrb1101}{A two-fibre product remembers the comparison isomorphism}
\begin{dcomment}{exposition}{p400:comme cas particulier de 4.14.19.}{A two-fibre product remembers the comparison isomorphism}
Grothendieck closes this sequence of corollaries by giving a useful model for gluing accessible categories. Given $p_1:\mathcal E_1\to\mathcal E_0$ and $p_2:\mathcal E_2\to\mathcal E_0$, the ordinary comma category remembers an arrow $p_1x_1\to p_2x_2$; its iso-comma subcategory requires that arrow to be invertible. The superscript $2$ is substantive: equality of two objects in the base is replaced by a specified isomorphism.

The cospan section category makes the proof concrete. Its objects consist of $x_i\in\mathcal E_i$ and arrows $u_i:p_ix_i\to x_0$. Require $u_2$ to be invertible. One can then discard $x_0$ and replace $u_1$ by $u_2^{-1}u_1$. The inverse operation sends $(x_1,x_2,u)$ to the quintuple with $x_0=p_2x_2$ and $u_2=\id$. The natural isomorphism between the composite and the original quintuple has component $u_2$ at the vertex $0$ and identity components at $1,2$. This checks the equivalence, including its action on morphisms, rather than merely comparing its object descriptions.

The condition on $u_2$ is an accessible inverse image of the isomorphism category of $\mathcal E_0$. The comma category is consequently accessible. Requiring the remaining comparison $u$ to be an isomorphism is another such inverse image, yielding the iso-comma theorem. This proof uses only the opfibration section theorem already established earlier, so its use in the subsequent proof for fibrations creates no circularity.

As an example, take two copies of $\mathbf1$ mapping into $\mathbf{Set}$ by selecting the same set $A$. Their strict pullback is $\mathbf1$. Their iso-comma category is the discrete category of permutations of $A$: an object is a bijection $u:A\to A$, and a morphism between two such objects must have identity components in both copies of $\mathbf1$, so it exists exactly when the two bijections are equal. The set of comparison isomorphisms is therefore part of the categorical answer. Discarding it would change the result even in this elementary diagram.

For the cospan opfibration, a cocartesian section requires both $u_1$ and $u_2$ to be invertible. Its section category is thus equivalent to the same iso-comma category, exactly as in the final display of the source. The missing subscripts and exchanged names in the surrounding paragraphs are edition slips; none affects this argument once the three categories, the comma category, its iso-comma subcategory, and the cospan section category, are kept distinct.
\end{dcomment}


\printedpage{152}
% head: A. Grothendieck : Dérivateurs, Ch. XVIII, Édition des Universités de Montpellier II et Paris VII

\textbf{4.14.22.} \textsc{Démonstration de} 4.14.17. Le cas où $p$ est cofibrant a déjà été démontré (th. 4.10.8, p. 263). Nous allons voir que le cas où $p$ est fibrant se ramène au théorème précédent (dans un cas très particulier), en utilisant à fond les résultats généraux sur les parties accessibles d’une catégorie accessible. Notons que les assertions a) et b) du corollaire 4.14.18 se prouvent de façon AQT \texttt{[âne qui trotte]} sans avoir à faire appel à 4.14.17. Le a) implique déjà que les foncteurs évaluation $\varepsilon_s : \mathcal{M} \longrightarrow \mathcal{E}_s$ sont accessibles. Il reste donc à prouver
\msp{page 401}
que $\mathcal{M} = \uGamma(\mathcal{E}/\mathcal{C})$ est accessible, dans le cas où $p$ est \textit{fibrant}. On supposera de plus que $\mathcal{C}$ est petite (quitte à remplacer $\mathcal{C}$ par une petite sous-catégorie pleine équivalente).

Soient
\[
\begin{aligned}
\Pi \ &= \ \prod\nolimits_{s\in\Ob\mathcal{C}} \mathcal{E}_s\\
\Pi' \ &= \ \prod\nolimits_{\lambda\in\Fl\mathcal{C}} \mathcal{E}_{s(\lambda)} \ ,
\end{aligned}
\]
où $s(\lambda)$ désigne la \textit{source} $s$ de la flèche $\lambda : s \longrightarrow t$ dans $\mathcal{C}$. $\Pi$, $\Pi'$ sont des catégories accessibles, comme (petits) produits de catégories accessibles. (Cas particulier trivial de 4.10.8, pour une catégorie d’indices discrète.) On va définir deux foncteurs
\[
\Pi \ \overset{F}{\underset{G}{\rightrightarrows}} \ \Pi'
\]
en définissant pour tout $\lambda : s_0 \longrightarrow t_0$ dans $\Fl\mathcal{C}$ deux foncteurs
\msp{page 402}
\[
\left\{
\begin{array}{lcll}
\Pi & \overset{F_\lambda}{\underset{G_\lambda}{\rightrightarrows}} & \mathcal{E}_{s(\lambda)=s_0} & \\
F_\lambda & = & \pr_{s_0} & \text{projection de } \Pi = \prod\nolimits_s \mathcal{E}_s \text{ sur le facteur } \mathcal{E}_{s_0}\\
G_\lambda & = & \lambda^*\pr_{s_1} & \text{composé } \Pi \xrightarrow{\pr_{s_1}} \mathcal{E}_{s_1} \xrightarrow{\lambda^*} \mathcal{E}_{s_0}\\
 & & & \text{\texttt{[plutôt} } \lambda^*\pr_{t_0} \text{ \texttt{resp.} } \Pi \xrightarrow{\pr_{t_0}} \mathcal{E}_{t_0} \xrightarrow{\lambda^*} \mathcal{E}_{s_0}\text{\texttt{]}}
\end{array}
\right.
\qquad \bigl(\mbox{\footnotemark[170]}, \mbox{\footnotemark[171]}\bigr) \ .
\]
\footnotetext[170]{Il est immédiat que les $F_\lambda$, $G_\lambda$, donc $F$ et $G$, sont des foncteurs \textit{accessibles}.}
\footnotetext[171]{\textbf{N.B.} On suppose choisi pour tout $\lambda \in \Fl\mathcal{C}$ un foncteur image inverse $\lambda^* : \mathcal{E}_t \longrightarrow \mathcal{E}_s$. On suppose que si $\lambda$ est une flèche identique, $\id_s$, \texttt{[on]} choisit $\lambda^* = \id_{\mathcal{E}_s}$.}

\dgrblock{dgrb1102}{First choose the arrows, then impose the laws}
\begin{dcomment}{exposition}{p401}{First choose the arrows, then impose the laws}
Grothendieck is now extracting a reusable construction from the section theorem. A section contains objects in the fibres, arrows above the arrows of the base, and equations expressing functoriality. The first two sorts of data can be assembled before the equations are imposed. This separation explains why accessibility survives: the preliminary construction is an inserter, and the final constraints are a small family of equifiers.

Choose a normalized cleavage of the fibration. A prospective section is a family $X=(X_s)$ together with arrows
\[
 a_\lambda:X_s\longrightarrow\lambda^*X_t
 \qquad(\lambda:s\longrightarrow t).
\]
For the moment these arrows need satisfy no laws. A morphism from $(X,a)$ to $(Y,b)$ is a family $h_s:X_s\to Y_s$ satisfying
\[
 \lambda^*(h_t)a_\lambda=b_\lambda h_s.
\]
This is precisely the inserter of the two functors $F,G:\Pi\rightrightarrows\Pi'$ on the page. The word ``diagram'' here means the underlying directed graph: it does not already incorporate identities or composition. Consequently the preliminary category is larger than the section category even when the underlying objects have been fixed.

The cleavage need not be split. If $\lambda:s\to t$ and $\mu:t\to u$, cartesian factorisation supplies a canonical isomorphism
\[
 c_{\lambda,\mu}:\lambda^*\mu^*\xrightarrow{\sim}(\mu\lambda)^*.
\]
The composition equation in fibre coordinates is
\[
 a_{\mu\lambda}
 =c_{\lambda,\mu,X_u}\,\lambda^*(a_\mu)\,a_\lambda.
\]
The unit equation is $a_{\id_s}=\id_{X_s}$ with the chosen normalization. Associativity of these comparisons is inherited from composition in $\mathcal E$. There is therefore no extra assumption that all transition functors compose as literal equalities. Grothendieck's use of the actual arrow $f_1(\mu)f_1(\lambda)$ keeps this coherence implicit and correct.

To see what the cardinal estimates must control, choose a regular cardinal $\kappa$ at which the fibres and the small family of transition functors admit and preserve filtered colimits. For a filtered diagram of prospective sections, put $X_s=\varinjlim_iX_{i,s}$. The structure arrow $a_\lambda$ is induced by the compatible arrows $X_{i,s}\to\lambda^*X_{i,t}$ and the comparison $\varinjlim_i\lambda^*X_{i,t}\simeq\lambda^*X_t$. The equations for a genuine section pass to these colimits. This proves colimit preservation by the relevant comparison functors; accessibility of the equation loci additionally uses the preceding small-generation results. Stability under filtered colimits alone would not supply that conclusion.

There is also a useful formula for the morphisms between two sections. They form the equalizer
\[
 \Hom_{\Gamma}(X,Y)\longrightarrow
 \prod_s\Hom_{\mathcal E_s}(X_s,Y_s)
 \rightrightarrows
 \prod_{\lambda:s\to t}
 \Hom_{\mathcal E_s}(X_s,\lambda^*Y_t).
\]
The two arrows send $(h_s)$ respectively to $(\lambda^*h_t)a_\lambda$ and $b_\lambda h_s$. If $\kappa$ exceeds the number of arrows of the base and every $X_s$ is $\kappa$-presentable, the right-hand expression preserves $\kappa$-filtered colimits in $Y$, provided the transition functors do. This gives the pointwise sufficient condition for presentability directly. The converse requires a larger common bound; forgetting structure does not preserve presentability at every initial rank.
\end{dcomment}


Soit $f \in \mathcal{M} = \uGamma(\mathcal{E}/\mathcal{C})$, alors pour tout $\lambda \in \Fl\mathcal{C}$, $f$ définit une flèche
\[
\alpha_\lambda(f) \ : \quad \underbrace{f(s)}_{= F_\lambda(\varepsilon(f))} \quad \xrightarrow{f(\lambda)^\flat} \quad \underbrace{\lambda^*(f(t))}_{= G_\lambda(\varepsilon(f))} \ ,
\]
\textit{fonctorielle} en $f$. On a un foncteur
\[
\varepsilon \ = \ (\varepsilon_s)_{s\in\Ob\mathcal{C}} \ : \ \mathcal{M} \ \longrightarrow \ \Pi \ ,
\]
et pour tout $\lambda \in \Fl\mathcal{C}$ un homomorphisme fonctoriel
\[
\alpha_\lambda \ : \ F_\lambda \circ \varepsilon \ \longrightarrow \ G_\lambda \circ \varepsilon \ ,
\]
d’où un homomorphisme fonctoriel
\[
\alpha \ = \ (\alpha_\lambda)_{\lambda\in\Fl\mathcal{C}} \ : \ F \circ \varepsilon \ \longrightarrow \ G \circ \varepsilon \ .
\]
Soit alors $\mathcal{M}'$ la catégorie suivante :
% foot: Version 17 décembre 2025
\printedpage{153}
% head: A. Grothendieck : Dérivateurs, Ch. XVIII, Édition des Universités de Montpellier II et Paris VII
\msp{page 403}
\[
\mathcal{M}' \quad = \quad \parbox[t]{0.5\textwidth}{catégorie des couples $(X,u)$, avec $X$ dans $\Ob\Pi$, et $u : F(X) \longrightarrow G(X)$ une flèche dans $\Pi'$.}
\]
L’homomorphisme fonctoriel $\alpha : F \circ \varepsilon \longrightarrow G \circ \varepsilon$ définit un foncteur
\[
\mathcal{M} \ \longrightarrow \ \mathcal{M}' \ ,
\]
et il est immédiat que ce foncteur est pleinement fidèle. En fait, la catégorie $\mathcal{M}'$ est isomorphe à la catégorie formée des couples $(f_0,f_1)$ de deux applications
\[
\begin{array}{lcccccc}
f_0 : & \Ob\mathcal{C} & \longrightarrow & \Ob\mathcal{E} & s & \longmapsto & f_0(s)\\[1ex]
f_1 : & \Fl\mathcal{C} & \longrightarrow & \Fl\mathcal{E} & \lambda & \longmapsto & f_1(\lambda)
\end{array}
\]
de telle façon que $f_1$ et $f_0$ soient compatibles avec les applications source et but, i.e. qu’ils définissent un homomorphisme de \textit{diagrammes} $f : \diag\mathcal{C} \longrightarrow \diag\mathcal{E}$, et que cet homomorphisme soit inverse à droite de $\mathrm{diag}(p) : \diag\mathcal{E} \longrightarrow \diag\mathcal{C}$.

\msp{page 404}
Le foncteur précédent $\mathcal{M} \hookrightarrow \mathcal{M}'$ s’interprète comme le foncteur qui associe à toute sections $f$ de $\mathcal{E}$ sur $\mathcal{C}$, le morphisme de diagrammes correspondant. C’est donc l’inclusion d’une sous-catégorie \textit{strictement pleine} de $\mathcal{M}'$. Donc pour prouver que $\mathcal{M}$ est accessible, il suffit de prouver : % [sic] toute sections -> toute section, agreement of toute
\begin{itemize}[leftmargin=2.5em]
\item[a)] $\mathcal{M}'$ est accessible.
\item[b)] $\mathcal{M}$ est accessible dans $\mathcal{M}'$.
\end{itemize}
Prouvons d’abord b) (en admettant a). La sous-catégorie $\mathcal{M}$ de $\mathcal{M}'$ est définie par deux types des conditions, sur un objet $f = (f_0,f_1)$ de $\mathcal{M}'$ :
\begin{itemize}[leftmargin=2.5em]
\item[1°)] $\forall\ s \in \Ob\mathcal{C}$, $f_1(\id_s) = \id_{f_0(s)}$.
\item[2°)] Si $s \xrightarrow{\lambda} t \xrightarrow{\mu} u$ sont deux flèches composables dans $\mathcal{C}$, alors
\[
f_1(\mu\lambda) \ = \ f_1(\mu)f_1(\lambda) \ .
\]
\end{itemize}
Ces conditions s’explicitent l’une et \texttt{[l’autre par]} l’égalité d’une double-flèche dans $\mathcal{E}_s$,
\msp{page 405}
savoir
\[
f_0(s) \overset{f_1(\id_s)}{\underset{\id_{f_0(s)}}{\rightrightarrows}} f_0(s)
\]
dans le cas 1°), et
\[
f_0(s) \overset{f_1(\mu\lambda)^\flat}{\underset{(f_1(\mu)f_1(\lambda))^\flat}{\rightrightarrows}} (\mu\lambda)^*(f_0(u))
\]
dans le cas 2°). Dans l’un et l’autre cas, on a un foncteur
\[
\mathcal{M}' \ \xrightarrow{\gamma_s \text{ resp. } \gamma_{\lambda,\mu}} \ \uHom(I,\mathcal{E}_s)
\]
% foot: Version 17 décembre 2025

% ---- from der-ch18_p156-160.tex
\printedpage{154}
% head: A. Grothendieck : Dérivateurs, Ch. XVIII, Édition des Universités de Montpellier II et Paris VII
dans la catégorie des double-flèches de $\mathcal{E}_s$ ($I$ désignant la catégorie double-flèche type $0\rightrightarrows 1$), foncteur qui est (trivialement) accessible. Soit alors $I\xrightarrow{\;i\;}\bDelta_1$ le foncteur canonique, et
\[
\underbrace{\uHom(\bDelta_1,\mathcal{E}_s)}_{\simeq\,\uFl(\mathcal{E}_s)}\ \hookrightarrow\ \uHom(I,\mathcal{E}_s)
\]
le foncteur $u\longmapsto u\circ i$, qui est un foncteur strictement pleine fidèle, ayant comme image la sous-catégorie des double-flèches $(u,v)$ telles que $u=v$. Comme le foncteur d'inclusion est accessible et $\uHom(\bDelta_1,\mathcal{E}_s)$ et $\uHom(I,\mathcal{E}_s)$ accessibles
\msp{page 406}
(cas particulier de 4.14.1, ou aussi de 4.10.8!), on voit que la sous-catégorie strictement pleine de $\uHom(I,\mathcal{E}_s)$ des double-flèches égalisées dans $\mathcal{E}_s$, est accessible dans $\uHom(I,\mathcal{E}_s)$. Donc, son image inverse dans $\mathcal{M}'$ est accessible. Donc l'intersection de ces images inverses pour $s\in\Ob\mathcal{C}$ et $(\lambda,\mu)\in\Fl_2(\mathcal{C})$ variables, l'est également ($\mathcal{C}$ étant petite).

Donc tout revient à prouver que $\mathcal{M}'$ est accessible, ce qui est contenu dans le

% vertical bar printed in the left margin alongside Lemme 4.14.22.1
\textbf{Lemme 4.14.22.1.} \textit{Soient $\Pi$, $\Pi'$ deux catégories accessibles, $F,G:\Pi\rightrightarrows\Pi'$ deux foncteurs accessibles, et $\mathcal{M}'$ la catégorie des couples $(X,u)$, avec $X\in\Ob\Pi$ et $u:F(x)\longrightarrow G(x)$ une flèche de $\Pi'$. Alors $\mathcal{M}'$ est accessible, et le foncteur canonique $\mathcal{M}'\longrightarrow\Pi$ est accessible.} % [sic] (X,u), X \in Ob Pi ... F(x) -> the same letter (X or x) throughout, capital X and lowercase x used for the same object

\textsc{Démonstration du Lemme.} La donnée de $\Pi$, $\Pi'$, $F$, $G$ définit une catégorie cofibrée $\mathcal{F}$ sur la catégorie $I=(0\rightrightarrows 1)$
\msp{page 407}
précédente. La catégorie $\uGamma(\mathcal{F}/I)$ s'interprète, à isomorphisme près, comme la catégorie des triples $(x,u,v)$ avec $x\in\Ob\Pi$, $x'\in\Ob\Pi'$, et $u$ et $v$ des flèches
\[
\begin{tikzcd}[column sep=small]
F(x) \arrow[dr, "u"'] & & G(x) \arrow[dl, "v"] \\
& x' &
\end{tikzcd}
\]
dans $\Pi'$. La catégorie $\mathcal{M}'$ est équivalente à la sous-catégorie strictement pleine de $\uGamma(\mathcal{F}/I)$, formée des triples $(x,u,v)$ tels que $v$ soit un isomorphisme, i.e. à la sous-catégorie de $\uGamma(\mathcal{F}/I)$, image inverse de $\uIs(\Pi')$ par le foncteur $(x,u,v)\longmapsto v$,
\[
\uGamma(\mathcal{F}/I)\ \longrightarrow\ \uFl(\Pi')\ .
\]
Or par 4.10.8, $\uGamma(\mathcal{F}/I)$ est accessible, et le foncteur précédent est bien sûr accessible. Enfin, par 4.14.12,
\msp{page 408}
$\uIs(\Pi')$ est accessible dans la catégorie accessible $\uFl(\Pi')$. Donc son image inverse $\mathcal{M}'$ dans $\uGamma(\mathcal{F}/I)$ est accessible dans $\uGamma(\mathcal{F}/I)$, donc accessible. Donc il en est de même de la catégorie équivalente, q.e.d.

\dgrblock{dgrb1103}{The inserter and the later limit theorems}
\begin{dcomment}{later}{p407}{The inserter and the later limit theorems}
The auxiliary lemma is the familiar inserter theorem. Its proof in the manuscript is especially useful because it avoids treating an inserter as an unexplained closure operation. A pair of arrows $FX\xrightarrow{u}Y\xleftarrow{v}GX$ becomes an inserter object when $v$ is invertible, by replacing the pair with $v^{-1}u:FX\to GX$. Conversely, an inserter object gives such a pair with $Y=GX$ and $v=\id$. On morphisms the two constructions are compatible, and the resulting comparison is an equivalence. The conclusion therefore follows from the already established accessibility of opfibration sections and of the isomorphism locus.

The surrounding theory had already acquired a systematic form in \citet[\S5.1]{MakkaiPare1989}. \citet[Theorems 3.1, 4.1, and 8.1--8.3]{Positselski2024Limits} gives explicit rank control for equifiers, inserters, and lax limits under a smaller-chain-colimit hypothesis. These hypotheses explain why existence of some accessibility rank is easier than retaining a prescribed rank. This literature confirms the construction here; it does not turn an arbitrary strict equalizer of accessible functors into an accessible category.

A small example makes the distinction visible. Let both $F$ and $G$ be the identity of $\mathrm{Set}$. Their inserter is the category of sets equipped with one endomorphism, equivalently sets with an action of the free monoid $\mathbb N$. The object $(\mathbb N,n\mapsto n+1)$ is free on one generator. A morphism from it to $(X,u)$ is uniquely specified by the image $x$ of $0$, because the image of $n$ must be $u^n(x)$. Hence
\[
 \Hom\bigl((\mathbb N,+1),(X,u)\bigr)\simeq X.
\]
Filtered colimits of these actions are computed on the underlying sets, so the free action is finitely presentable. Its underlying set is infinite and is not finitely presentable in $\mathrm{Set}$. To verify the latter assertion, express $\mathbb N$ as the directed union of its finite subsets: its identity map does not factor through a finite subset. Thus the inserter projection is accessible while failing to preserve finitely presentable objects. There is no conflict with the manuscript's eventual comparisons after enlarging the cardinal.

For higher categories the corresponding section statement is explicit in \citet[Proposition 5.4.7.11 and Remarks 5.4.7.13--14]{Lurie2009HTT}: accessible fibres and accessible transport over a small simplicial base give accessible categories of sections, with specified cartesian or cocartesian conditions allowed. The higher theorem includes coherence as part of the section space. The ordinary equations written here give its one-categorical counterpart.
\end{dcomment}


Ainsi, grâce aux résultats généraux sur les parties accessibles des catégories accessibles, le théorème 4.14.17, dans le cas où $p$ est fibrant, est ramené aux cas cofibrant, dans le cas très particulier où la catégorie d'indices est $I=(0\rightrightarrows 1)$. On pourrait d'ailleurs faire une réduction similaire dans le cas où $p$ est cofibrant, par une démonstration pratiquement % [sic] aux cas cofibrant -> au cas cofibrant, singular noun
% foot: Version 17 décembre 2025
\printedpage{155}
% head: A. Grothendieck : Dérivateurs, Ch. XVIII, Édition des Universités de Montpellier II et Paris VII
identique. En d'autres termes, le théorème 4.10.8 se ramène au cas où la catégorie d'indices est $I=(0\rightrightarrows 1)$. Mais je suspecte que la démonstration du théorème dans ce cas particulier, n'est pas essentiellement plus simple que dans le cas général.

\msp{page 409}

\textbf{4.14.23.} On peut se demander de trouver un énoncé assurant l'accessibilité d'une catégorie $\uGamma(\mathcal{E}/\mathcal{C})$, qui coiffe les deux cas envisagés dans 4.14.17. Voici ce que je peux proposer. Je vais faire l'hypothèse suivante sur les foncteur $p:\mathcal{E}\longrightarrow\mathcal{C}$ (en plus de celle que $\mathcal{C}$ soit essentiellement petite) : % [sic] les foncteur -> le foncteur, singular functor p

\begin{itemize}[leftmargin=*]
\item[H)] Pour tout $\lambda\in\Fl\mathcal{C}$, la catégorie $\mathcal{E}_\lambda=\mathcal{E}\times_{\mathcal{C}}(\bDelta_1,i_\lambda)$, où $i_\lambda:\bDelta_1\longrightarrow\mathcal{C}$ est le foncteur défini par $\lambda$, est accessible, et le foncteur canonique $\mathcal{E}_\lambda\xrightarrow{\;p_\lambda\;}\bDelta_1$ est accessible.
\end{itemize}

Notons que si $\lambda:s\longrightarrow t$ dans $\mathcal{C}$ est tel que $\mathcal{E}_\lambda$ et $p_\lambda$ soient accessibles, alors $\mathcal{E}_s$ et $\mathcal{E}_t$ sont accessible, car isomorphes aux catégories image inverse par $p_\lambda:\mathcal{E}_\lambda\longrightarrow\bDelta_1$ des sous-catégories $\{0\}$ resp. $\{1\}$ de $\Delta_1$, lesquelles sont accessibles dans la catégorie accessible $\bDelta_1$, donc % [sic] sont accessible -> sont accessibles, plural agreement
\msp{page 410}
$\mathcal{E}_{\lambda_0}$ et $\mathcal{E}_{\lambda_1}$ sont accessibles dans $\mathcal{E}_\lambda$, donc sont accessibles, donc aussi $\mathcal{E}_s$ et $\mathcal{E}_t$. Ainsi, l'hypothèse H) implique que les catégories fibres sont accessibles. Notons de plus :

\textbf{Lemme 4.14.23.1.} \textit{Supposons que $p:\mathcal{E}\longrightarrow\mathcal{C}$ soit cofibrant ou fibrant, et soit $(\lambda:s\longrightarrow t)\in\Fl(\mathcal{C})$. Pour que $\mathcal{E}_\lambda$ soit accessible et $\mathcal{E}_\lambda\xrightarrow{\;p_\lambda\;}\bDelta_1$ accessible, il faut et il suffit que $\mathcal{E}_s$, $\mathcal{E}_t$ soient accessibles et $\lambda_*:\mathcal{E}_s\longrightarrow\mathcal{E}_t$ resp. $\lambda^*:\mathcal{E}_t\longrightarrow\mathcal{E}_s$ soit un foncteur accessible.}

\textsc{Démonstration.} On peut supposer que $\mathcal{C}=\bDelta_1$, et que $\mathcal{E}$ est la catégorie cofibrée (resp. fibrée) sur $\bDelta_1$ définie par un foncteur $\varphi:\mathcal{E}_0\longrightarrow\mathcal{E}_1$ (resp. $\psi:\mathcal{E}_1\longrightarrow\mathcal{E}_0$), et à montrer que $\mathcal{E}$ et $p:\mathcal{E}\longrightarrow\bDelta_1$ sont accessibles si et seulement si $\mathcal{E}_0$ et $\mathcal{E}_1$ et $\varphi$ (resp. $\psi$) le sont.

\msp{page 411}

Prouvons le “il faut”. On a déjà vu que si $\mathcal{E}$, $p$ sont accessibles, $\mathcal{E}_0$ et $\mathcal{E}_1$ sont accessibles dans $\mathcal{E}_s$, donc accessibles, et les foncteurs d'inclusion % [sic] dans E_s -> dans E, here C = Delta_1 and the subcategories lie in E
\[
\mathcal{E}_0\ \hookrightarrow\ \mathcal{E}\ \hookleftarrow\ \mathcal{E}_1
\]
accessibles. Dire que $p$ est cofibrant (resp. fibrant) signifie que le foncteur d'inclusion $\mathcal{E}_1\hookrightarrow\mathcal{E}$ admet un adjoint à droite $\bar{\varphi}$ (resp. que le foncteur $\mathcal{E}_0\hookrightarrow\mathcal{E}$ admet un adjoint à gauche $\tilde{\psi}$). Alors $\varphi$ (resp. $\psi$) est le composé $\mathcal{E}_0\longrightarrow\mathcal{E}\xrightarrow{\;\bar{\varphi}\;}\mathcal{E}_1$ (resp. le composé $\mathcal{E}_1\hookrightarrow\mathcal{E}\xrightarrow{\;\tilde{\psi}\;}\mathcal{E}_1$). Donc pour prouver que $\varphi$ (resp. $\psi$) est accessible, il suffit de prouver que $\bar{\varphi}$ (resp. $\tilde{\psi}$) l'est. Cela est un cas particulier du % [sic] E -> E_1 (target of psi-tilde composite) -> E_0, psi goes from E_1 to E_0

\textbf{Lemme 4.14.23.2.} \textit{Soient $\mathcal{E}$, $\mathcal{E}'$ deux catégories accessibles} (\footnote[172]{\textbf{N.B.} Il suffit que $\mathcal{E}$ admette une petite sous-catégorie génératrice, et que les objets de $\mathcal{E}'$ soient accessibles.}), \textit{et $\mathcal{E}\underset{G}{\overset{F}{\rightleftarrows}}\mathcal{E}'$ un couple de foncteurs adjoints entre $\mathcal{E}$ et $\mathcal{E}'$. Alors $F$ et $G$ sont accessibles.}

\msp{page 412}

En effet, $F$ est accessible, car il commute aux petites $\varinjlim$ quelconques, a fortiori il satisfait $\mathrm{L}_{\pi_0}$ (où $\pi_0$ est tel que $\mathcal{E}$, $\mathcal{E}'$ satisfassent $\mathrm{L}_{\pi_0}$). Prouvons que $G$ est accessible. Si $S$ est un
% foot: Version 17 décembre 2025
\printedpage{156}
% head: A. Grothendieck : Dérivateurs, Ch. XVIII, Édition des Universités de Montpellier II et Paris VII
petit ensemble d'éléments générateur dans $\mathcal{E}$, de sorte que la petite famille des foncteurs
\[
h_\xi\ =\ (x\longmapsto\Hom_{\mathcal{E}}(\xi,x))\ :\ \mathcal{E}\ \longrightarrow\ \Ens
\]
est conservative, pour prouver que $G$ est accessible, il suffit de prouver que les $h_\xi\circ G$ ($\xi\in S\subset\Ob\mathcal{E}$) le sont. Or
\[
L_\xi\circ G(x')\ =\ \Hom(\xi,G(x'))\ \simeq\ \Hom_{\mathcal{E}'}(F(\xi),x')\ , % [sic] L_xi -> h_xi, the functors were named h_xi
\]
et comme $F(\xi)$ est un objet accessible de $\mathcal{E}'$, le foncteur
\[
x'\ \longmapsto\ (L_\xi\circ G)(x')\ \simeq\ \Hom_{\mathcal{E}'}(F(\xi),x') % [sic] L_xi -> h_xi, the functors were named h_xi
\]
est bien accessible, q.e.d.

Cela prouve le “il faut” dans 4.14.23.1. Prouvons le “il suffit”, i.e. supposons $\mathcal{E}_s$, $\mathcal{E}_t$ accessibles, et $\varphi$ (resp. $\psi$) accessible.

\msp{page 413}

\dgrblock{dgrb1104}{Why both adjoints are accessible}
\begin{dcomment}{exposition}{p412}{Why both adjoints are accessible}
Grothendieck needs the adjunction lemma to recover a transition functor from the total category over an interval. The lemma concerns the existence of a sufficiently large filtered-colimit rank. It does not assert that a right adjoint preserves all filtered colimits, or that the same rank works for every adjunction.

Write $F:\mathcal A\rightleftarrows\mathcal B:G$ with $F\dashv G$. Choose a small strong generating family $S$ of presentable objects of $\mathcal A$. Since the family is small and all the objects $F(s)$ are presentable at some ranks, choose one regular $\kappa$ such that both categories admit $\kappa$-filtered colimits and every $s$ and $F(s)$ is $\kappa$-presentable. For a $\kappa$-filtered diagram $(b_i)$, consider
\[
 q:\varinjlim_iG(b_i)\longrightarrow G(\varinjlim_i b_i).
\]
For each $s\in S$, adjunction and the two presentability assertions give
\[
\begin{aligned}
 \Hom_{\mathcal A}(s,\varinjlim_iG b_i)
 &\simeq\varinjlim_i\Hom_{\mathcal A}(s,G b_i)\\
 &\simeq\varinjlim_i\Hom_{\mathcal B}(F s,b_i)\\
 &\simeq\Hom_{\mathcal B}(F s,\varinjlim_i b_i)\\
 &\simeq\Hom_{\mathcal A}(s,G\varinjlim_i b_i).
\end{aligned}
\]
This composite is $\Hom(s,q)$, so the strong generator implies that $q$ is an isomorphism. The left adjoint preserves all colimits that exist, which proves its accessibility at once. The modern reference for the lemma is \citet[Proposition 2.23]{AdamekRosicky1994}; the argument just given makes the otherwise implicit common-cardinal choice explicit.

The word ``generating'' on page 412 must include the stated conservativity, and the calculation also uses presentability of the chosen generators. Both are available for an accessible category. Footnote 172 suggests a weaker hypothesis on $\mathcal A$; the printed argument does not explain how to dispense with preservation by these representables or with the relevant colimits. We use the lemma under its stated accessible-category hypotheses and do not certify that additional weakening.

Two orientation slips on page 411 should be read through this argument. For an opfibration classified by $\varphi:\mathcal E_0\to\mathcal E_1$, extend $\varphi$ to $\bar\varphi:\mathcal E\to\mathcal E_1$ by the identity on $\mathcal E_1$. Then
\[
 \Hom_{\mathcal E}(x,i_1y)
 \simeq\Hom_{\mathcal E_1}(\bar\varphi x,y),
 \qquad\bar\varphi\dashv i_1.
\]
For a fibration classified by $\psi:\mathcal E_1\to\mathcal E_0$, the extension satisfies $i_0\dashv\widetilde\psi$. The edition prints the opposite adjunction directions, as well as the wrong final fibre in the composite defining $\psi$. Correcting the types leaves the proof of Lemma 4.14.23.1 intact because the next lemma applies to either adjoint.
\end{dcomment}


Dans les deux cas, $\mathcal{E}_0$ est un ouvert de $\mathcal{E}$, $\mathcal{E}_1$ est le fermé complémentaire, et $\mathcal{E}$ est la catégorie “recollée” à l'aide du bifoncteur
\[
H\ :\ \mathcal{E}_0^{\circ}\times\mathcal{E}_1\ \longrightarrow\ \Ens
\]
donné par
\[
\left\{
\begin{aligned}
H(x_0,x_1)\ &=\ \Hom(\varphi(x_0),x_1)\\
&\hspace{-2em}\text{resp.}\\
H(x_0,x_1)\ &=\ \Hom(x_0,\psi(x_1))\ ,
\end{aligned}
\right.
\]
suivant qu'on est dans le cas $p$ cofibrant ou $p$ fibrant. On va utiliser la

\textbf{Proposition 4.14.24} (\footnote[173]{cf. énoncés plus généraux et démonstration sympathique, 4.15.1 et 4.15.4 (p. 425, 432).}). \textit{Soit $\mathcal{E}$ une catégorie, $\mathcal{E}_0$ une sous-catégorie ouverte, $\mathcal{E}_1$ la sous-catégorie fermée complémentaire. Conditions équivalentes :}
\begin{itemize}[leftmargin=*]
\item[\textit{a)}] \textit{$\mathcal{E}$ est accessible, et $\mathcal{E}\longrightarrow\bDelta_1$ (de fibres $\mathcal{E}_0$, $\mathcal{E}_1$) est accessible.}
\item[\textit{b)}] \textit{$\mathcal{E}$ est accessible, et $\mathcal{E}_0$, $\mathcal{E}_1$ sont accessibles dans $\mathcal{E}$.}
\item[\textit{c)}] \textit{$\mathcal{E}_0$, $\mathcal{E}_1$ sont accessibles, et pour tout $x_0\in\Ob\mathcal{E}_0$, le foncteur}
\[
x_1\ \longmapsto\ \uHom_{\mathcal{E}}(x_0,x_1)\ :\ \mathcal{E}_1\ \longrightarrow\ \Ens
\]
\msp{page 414}
\textit{est accessible, et il existe $\pi$ tel que pour tout $x_1\in\Ob\mathcal{E}_1$, le foncteur}
\[
x_0\ \longmapsto\ \uHom_{\mathcal{E}}(x_0,x_1)^{\circ}\ :\ \mathcal{E}_0\ \longrightarrow\ (\Ens)^{\circ}
\]
\textit{soit $\pi$-accessible.}
\end{itemize}
% foot: Version 17 décembre 2025

\printedpage{157}
% head: A. Grothendieck : Dérivateurs, Ch. XVIII, Édition des Universités de Montpellier II et Paris VII
Admettons ce lemme \texttt{[plutôt cette proposition]}, achevons la démonstration de 4.14.23.1. On suppose $\mathcal{E}_0$, $\mathcal{E}_1$ et $\varphi$ (resp. $\psi$) accessibles, on veut prouver la condition a) de 4.14.23.3. Pour ceci, il suffit de prouver la condition c). Or on a, dans le cas $p$ cofibrant, % [sic] 4.14.23.3 -> 4.14.24, the conditions a), c) are those of Proposition 4.14.24
\[
\Hom_{\mathcal{E}}(x_0,x_1)\ \simeq\ \Hom_{\mathcal{E}_1}(\varphi(x_0),x_1)\ ,
\]
c'est bien accessible en $x_1$, puisque $\mathcal{E}_1$ est accessible, et $\pi$-accessible en $x_0$ si $\varphi$ l'est, comme composé de $\varphi$ avec $\mathcal{E}_1\longrightarrow\Ens$, $\Hom(-,x_1)$, qui l'est. Dans le cas $p$ fibrant, on a
\[
\Hom_{\mathcal{E}}(x_0,x_1)\ \simeq\ \Hom_{\mathcal{E}_0}(x_0,\psi(x_1))\ ,
\]
c'est accessible en $x_1$, comme le composé du foncteur accessible $\psi:\mathcal{E}_1\longrightarrow\mathcal{E}_0$, et du foncteur également accessible $\Hom(x_0,-):\mathcal{E}_0\longrightarrow\Ens$, qui est $\pi$-accessible en $x_0$ tautologiquement.

Pour achever la démonstration de 4.14.23.1, il reste donc à prouver la proposition 4.14.24. On a déjà noté qu'on a a $\Longrightarrow$ b, et b $\Longrightarrow$ c clair, puisque le foncteur envisagé $\mathcal{E}\longrightarrow\Ens$ est le composé $\mathcal{E}_1\overset{\text{incl.}}{\hookrightarrow}\mathcal{E}\xrightarrow{\;\Hom_{\mathcal{E}}(x_0,-)\;}\Ens$, est accessible comme composé de deux
\msp{page 415}
foncteurs accessibles. Reste donc à prouver c $\Longrightarrow$ a. Soit donc $\pi_0$ un cardinal utile pour $\mathcal{E}_0$, et tel que les foncteurs
\[
x_0\ \longmapsto\ H(x_0,x_1)\ :\ \mathcal{E}_0\longrightarrow\Ens^{\circ}
\]
soient $\pi_0$-accessibles, soit $\pi$ un cardinal $\geqslant\sup(\pi_0,\mathrm{card}\,\text{réd}\,\mathcal{E}_{0\,\pi_0})$ tel que $\pi^{\pi_0}=\pi$ (donc utile pour $\mathcal{E}_0$), utile pour $\mathcal{E}_1$, et tel que les foncteurs
\[
x_1\ \longmapsto\ H(x_0,x_1)\ :\ \mathcal{E}_1\longrightarrow\Ens
\]
soient $\pi$-accessibles, pour $x_0\in\Ob\mathcal{E}_{0\,\pi_0}$. Soit $\mathcal{E}(\pi)$ la sous-catégorie strictement pleine de $\mathcal{E}$, telle que
\[
\Ob\mathcal{E}(\pi)\ =\ \Ob\mathcal{E}_{0\,\pi}\cup\Ob\mathcal{E}_{1\,\pi}\ .
\]
Je dis que
\begin{itemize}[leftmargin=*]
\item[1°)] $\mathcal{E}$ satisfait à $\mathrm{L}_\pi$, et $\mathcal{E}_0$ et $\mathcal{E}_1$ sont stables dans $\mathcal{E}$ sous $\mathrm{L}_\pi$, les foncteurs $\mathcal{E}_0\longrightarrow\mathcal{E}$, $\mathcal{E}_1\longrightarrow\mathcal{E}$, $\mathcal{E}\longrightarrow\bDelta_1$ satisfont $\mathrm{L}_\pi$ (\footnote[174]{il suffit que $\mathcal{E}_0$ et $\mathcal{E}_1$ satisfassent $\mathrm{L}_\pi$, et que
\[
x_0\ \longmapsto\ H(x_0,x_1)\ :\ \mathcal{E}_0\longrightarrow\Ens^{\circ}
\]
satisfasse $\mathrm{L}_\pi$ $\forall\ x\in\Ob\mathcal{E}$.}).
\item[2°)] $\mathcal{E}(\pi)=\mathcal{E}_\pi$ (d'où un foncteur pleinement fidèle
\[
\mathrm{Ind}_\pi(\mathcal{E}(\pi))\ \longrightarrow\ \mathcal{E}\ \ ).
\]
\item[3°)] Ledit foncteur est essentiellement surjectif,
\msp{page 416}
donc une équivalence de catégories.
\end{itemize}
Notons que 2°) et 3°) impliquent que $\mathcal{E}$ est accessible, et 1°) que $\mathcal{E}\longrightarrow\bDelta_1$ est accessible, d'où la condition a) du lemme qu'il s'agit d'établir. Donc tout revient à prouver les 1°, 2°, 3° ci-dessus.
% foot: Version 17 décembre 2025

\dgrblock{dgrb1105}{A collage is controlled by its cross-Homs}
\begin{dcomment}{exposition}{p414}{A collage is controlled by its cross-Homs}
The passage to a gluing bifunctor frees the argument from the existence of chosen lifts. An open subcategory here is a sieve, and its closed complement is a cosieve. Thus there are no arrows from $\mathcal E_1$ to $\mathcal E_0$. The remaining arrows between the two pieces are recorded by
\[
 H:\mathcal E_0^{\mathrm{op}}\times\mathcal E_1\longrightarrow\mathrm{Set}.
\]
The category reconstructed from this information is now called the collage of $H$. Its composition laws involving a cross-arrow are exactly the two functorial actions of $H$.

The two preservation conditions in Proposition 4.14.24 have different jobs. The contravariant condition says that a colimit formed in $\mathcal E_0$ remains universal when tested against objects of $\mathcal E_1$. The covariant condition says that an object of $\mathcal E_0$ eventually remains small when maps from it are tested on filtered diagrams in $\mathcal E_1$. The first condition needs a common rank for all target objects; the second can initially have a different rank for each source object. Accessibility permits us to bound that second family on a small generating subcategory.

Here is the colimit argument without the incidental notation slips. Fix a regular $\kappa$ such that both fibres have $\kappa$-filtered colimits and each $H(-,y)$ converts those in $\mathcal E_0$ to limits. A diagram contained in $\mathcal E_0$ has its fibrewise colimit in the total category because it has the required universal maps to targets in either fibre. A diagram that has reached $\mathcal E_1$ stays there after passing to a cofinal tail. Its colimit in $\mathcal E_1$ is also a colimit in the total category: there are no maps from this nonempty diagram to an object of $\mathcal E_0$. These two cases exhaust a filtered diagram.

The smallness argument is equally concrete. An object $x\in\mathcal E_1$ is $\kappa$-presentable in the collage precisely when it is $\kappa$-presentable in its fibre. For $x\in\mathcal E_0$ one must additionally require that $H(x,-)$ preserve $\kappa$-filtered colimits. The calculation on page 419 obtains this latter requirement from a small presentation of $x$ in $\mathcal E_0$. Contravariance gives
\[
 H(\varinjlim_j a_j,-)\simeq\varprojlim_jH(a_j,-),
\]
and a sufficiently filtered colimit of sets commutes with this small limit. This is why the cardinal controlling the presentation, the variance condition, and the size of its indexing category must be chosen together.

Once these properties hold at a useful common cardinal, the union of the two small generating subcategories generates the collage. This proves accessibility itself, not merely closure under filtered colimits. The proof also explains the two familiar cases. For an opfibration, $H(x,y)=\Hom(\varphi x,y)$, and the relevant size of a lower-fibre object is the size of $\varphi x$. For a fibration, $H(x,y)=\Hom(x,\psi y)$, and preservation by $\psi$ supplies the corresponding comparison.
\end{dcomment}


\printedpage{158}
% head: A. Grothendieck : Dérivateurs, Ch. XVIII, Édition des Universités de Montpellier II et Paris VII
1°) Comme $\mathcal{E}_1$ est fermée dans $\mathcal{E}$, il est clair que $\mathcal{E}$ est stable par toutes $\varinjlim$ représentables dans $\mathcal{E}$, et plus précisément, que pour un foncteur $I\longrightarrow\mathcal{E}_1$, sa limite inductive dans $\mathcal{E}_1$ existe si et seulement si elle existe dans $\mathcal{E}$, et les deux sont isomorphes. Supposons $I$ grande devant $\pi$, et soit $\varphi:I\longrightarrow\mathcal{E}$, $i\longmapsto x_i$. Distinguons deux cas.

a) $\varphi$ prend ses valeurs dans $\mathcal{E}_0$. Soit $x_0$ sa limite inductive \textit{dans} $\mathcal{E}_0$. Je dis que c'est aussi une limite inductive dans $\mathcal{E}$. Il faut prouver que pour $x_1\in\Ob\mathcal{E}_1$, la flèche
\[
\Hom_{\mathcal{E}}(x_0,x_1)\ \longrightarrow\ \varprojlim_{i\in I}\Hom_{\mathcal{E}}(x_i,x_1)
\]
est isomorphe. Or c'est le cas
\msp{page 417}
par hypothèse sur $\pi$,
\[
\xi\ \longmapsto H(\xi,x_1)\ :\ \mathcal{E}_0\ \longrightarrow\ (\Ens)^{\circ}
\]
satisfait à $\mathrm{L}_\pi$ pour tout $x_1$ dans $\Ob\mathcal{E}_1$.

b) $\varphi$ ne se factorise pas par $\mathcal{E}_0$, i.e. $\exists\ i_0\in\Ob I$ avec $\varphi(i_0)\in\Ob\mathcal{E}_1$. Alors on a $\varphi(i)\in\Ob\mathcal{E}_1$ pour tout $i\in\Ob I[i_0]$, où $I[i_0]$ désigne la sous-catégorie strictement pleine de $I$ formée des $i$ tels que $i\geqslant i_0$. Cette catégorie est cofinale dans $I$, donc grande devant $\pi$ comme $I$, et pour que $\varinjlim\limits_{I}x_i$ existe, il faut et il suffit que $\varinjlim\limits_{I[i_0]}x_i$ existe, ce qui est le cas d'après les cas a) déjà traité. % [sic] les cas a) déjà traité -> le cas a) déjà traité, singular agreement

Cela prouve donc en même temps que $\mathcal{E}$ satisfait à $\mathrm{L}_\pi$, et que $\mathcal{E}_0$ et $\mathcal{E}_1$ sont $\mathrm{L}_\pi$-stables. La démonstration montre en même temps que le foncteur $p:\mathcal{E}\longrightarrow\bDelta_1$ satisfait à $\mathrm{L}_\pi$.

2°) De ceci résulte aussitôt que $\mathcal{E}_\pi\subset\mathcal{E}(\pi)$. Prouvons qu'on a aussi
\[
\mathcal{E}(\pi)\ \subset\ \mathcal{E}_\pi\ ,
\]
et distinguons encore deux cas, pour
\msp{page 418}
$x\in\Ob E(\pi)$, suivant que $x$ est dans $\mathcal{E}_{0\,\pi}$ ou $\mathcal{E}_{1\,\pi}$. S'il est dans $\mathcal{E}_{1\,\pi}$, pour prouver que % [sic] E(pi) -> \mathcal{E}(pi), script letter elsewhere
\[
\Hom_{\mathcal{E}}(x,\underbrace{\varinjlim y_i}_{y})\ \xleftarrow{\ \sim\ }\ \varinjlim_{i}\Hom_{\mathcal{E}}(x,y_i)
\tag*{$(*)$}
\]
(pour un système inductif $(y_i)_{i\in I}$ dans $\mathcal{E}$ grand devant $\pi$), si les $y_i$ sont dans $\mathcal{E}_0$, donc $y\in\mathcal{E}_0$, alors les deux membres sont vides, et c'est OK. Sinon, on peut supposer (quitte à remplacer $I$ par un $I[i_0]$) que les $x_i$ sont dans $\mathcal{E}_1$, et alors c'est OK aussi, puisque $y$ est aussi la $\varinjlim$ dans $\mathcal{E}_1$, d'après 1°). % [sic] les x_i -> les y_i, the inductive system is (y_i)

Si $x$ est dans $\mathcal{E}_{0\,\pi}$, alors $(*)$ est clair si les $y_i$ sont dans $\mathcal{E}_0$. Sinon, on peut supposer encore que les $y_i$ sont dans $\mathcal{E}_0$, et $(*)$ s'écrit % [sic] y_i sont dans E_0 (second occurrence) -> E_1, the case "sinon" and the following "les y_i dans E_1"
\[
H(x,\varinjlim y_i)\ \xleftarrow{\ \sim\ }\ \varinjlim H(x,y_i)\ ,
\]
pour $x\in\Ob\mathcal{E}_{0\,\pi}$, les $y_i$ dans $\mathcal{E}_1$, $I$ grande devant $\pi$. En d'autres termes, il faut prouver que le foncteur
\[
x_1\ \longmapsto\ H(x_0,x_1)\ :\ \mathcal{E}_1\ \longrightarrow\ \Ens
\]
% foot: Version 17 décembre 2025

% ---- from der-ch18_p161-165.tex
\setcounter{footnote}{174}

%=====================================================================
\printedpage{159}
% head: A. Grothendieck : Dérivateurs, Ch. XVIII, Édition des Universités de Montpellier II et Paris VII

\msp{page 419}
est $\pi$-accessible, quand $x_0$ est dans $\cE_{0\,\pi}$. Or par hypothèse sur $\pi$, il est $\pi$-accessible si $x_0$ est dans $\cE_{0\,\pi_0}$. Mais si $x$ est dans $\cE_{0\,\pi}$, on aura % [sic] x -> x_0, the formula that follows concerns x_0
\[
x_0 \;=\; \varinjlim_{i\in I} x_{0\,i} \qquad
\left\{
\begin{array}{l}
I\ \ \text{grande}\ \ \text{devant}\ \ \pi_0,\\
\text{les}\ \ x_{0\,i}\ \ \text{dans}\ \ \cE_{0\,\pi_0},\\
\card I \leqslant \pi
\end{array}
\right.
\]
(car $\pi_0$ est utile pour $\cE_0$ par hypothèse, et $\pi$ satisfait $\pi^{\pi_0} = \pi \geqslant \cardred \cE_{\pi_0}$). Comme $x_0 \longrightarrow H(x_0,x_1) : \cE_0 \longrightarrow \Ens$ satisfait à $\mathrm{L}_{\pi_0}$, pour tout $x_1$ dans $\cE_1$, on aura donc % [sic] x_0 \longrightarrow H(x_0,x_1) -> x_0 \longmapsto H(x_0,x_1), elsewhere such functors are written with \longmapsto
\[
H(x_0) \;\simeq\; \varprojlim_{i\in I} H(x_{0\,i})
\]
(où pour $x_0$ dans $\cE_0$, $H(x)$ désigne le foncteur $x_1 \longmapsto H(x_0,x_1) : \cE_1 \longrightarrow (\Ens)$). Comme % [sic] H(x) -> H(x_0), the functor depends on x_0
les foncteurs $H(x_{0\,i})$ sont $\pi$-accessibles, et $\card I \leqslant \pi$, il s'ensuit que leur $\varprojlim\limits_{I}$ $H(x_0)$ est $\pi$-accessible (SGA 4 I 9.6 (ii)). Cela achève de prouver donc 2°).

3°). Pour prouver que $\Ind_\pi(\cE(\pi)) \longrightarrow \cE$ est essentiellement surjectif, il suffit de voir que tout $x \in \Ob\cE_0$ et tout $x \in \Ob\cE_1$
\msp{page 420}
est dans l'image essentielle. Or cela résulte de l'hypothèse que $\pi$ est utile pour $\cE_0$ et $\cE_1$, et de 1°) qui implique que $\cE_0 \longrightarrow \cE$ et $\cE_1 \longrightarrow \cE$ satisfont $\mathrm{L}_\pi$.

Cela achève de prouver la proposition 4.14.24, et de même le résultat suivant, plus précis que l'implication c) $\Longrightarrow$ a).

\textbf{Corollaire 4.14.25.} \textit{Soient $\cE$ une catégorie, $\cE_0$ une sous-catégorie ouverte, $\cE_1$ la sous-catégorie fermée complémentaire,}
\[
H \;=\; ((x_0,x_1) \longmapsto \Hom_{\cE}(x_0,x_1)) \;:\; \cE_0^\circ \times \cE_1 \;\longrightarrow\; \Ens
\]
\textit{le foncteur recollement.}
\begin{enumerate}
\item[\textit{1°)}] \textit{Soit $\pi$ un cardinal infini. Pour ce que $\cE$ et $\cE_0$ satisfassent à $L_\pi$, et le foncteur d'inclusion $\cE_0 \hookrightarrow \cE$ itou, il faut et il suffit que $\cE_0$ satisfasse à $L_\pi$, et que pour tout $x_1$ dans $\cE_1$, le foncteur $x_0 \longmapsto H(x_0,x_1) : \cE_0 \longrightarrow \Ens^\circ$ satisfasse à $L_\pi$. Alors $\cE_1$ et $\cE_1 \hookrightarrow \cE$ satisfont aussi à $L_\pi$,}
\msp{page 421}
\textit{et on a}
\[
\cE_\pi \;\subset\; \cE(\pi) \;\defeq\;
\begin{array}{l}
\text{\itshape sous-catégorie\ \ strictement\ \ pleine\ \ de\ \ $\cE$\ \ telle\ \ que}\\
\Ob\cE(\pi) = \Ob\cE_{0\,\pi} \cup \Ob\cE_{1\,\pi}.
\end{array}
\]
\item[\textit{2°)}] \textit{Supposons que la condition c) de 4.14.24 soit satisfaite, et soit $\pi_0$ un cardinal infini tel que $\pi_0$ soit utile pour $\cE_0$, et que les foncteurs $x_0 \longmapsto H(x_0,x_1) : \cE_0 \longrightarrow \Ens^\circ$ définis par les $x_1 \in \Ob\cE_1$ soient $\pi_0$-accessibles. Soit $\pi$ un cardinal tel que $\pi^{\pi_0} = \pi \geqslant \sup(\pi_0, \cardred\cE_{0\,\pi_0})$, et tel que pour tout $x_0 \in \Ob\cE_{\pi_0}$, le foncteur $x_1 \longmapsto H(x_0,x_1) : \cE_1 \longrightarrow \Ens$ soit $\pi$-accessible. Alors on a}
\[
\cE_\pi \;=\; \cE(\pi)\ .
\]
% foot: Version 17 décembre 2025

\dgrblock{dgrb1106}{The missing closed-fibre hypothesis}
\begin{dcomment}{status}{p420}{The missing closed-fibre hypothesis}
The first part of Corollary 4.14.25 is stronger than the preceding proof and is false as printed. The proof assumed that both fibres had the relevant filtered colimits. The right-hand side of the corollary mentions only $\mathcal E_0$ and the functors $H(-,y)$, then attempts to deduce colimit existence in $\mathcal E_1$. Nothing in those conditions controls a diagram entirely within the closed fibre.

For a nonempty counterexample, fix the manuscript cardinal $\pi$ and put $\kappa=\pi^+$. Take $\mathcal E_0$ to be the terminal category, $\mathcal E_1$ the ordinal $\kappa$ regarded as an ordered category, and $H$ the constant singleton functor. The collage is the ordinal $\kappa$ with an additional least object. The open fibre has all filtered colimits, and every $H(-,y)$ takes its filtered colimits to limits. Nevertheless the identity diagram $\kappa\to\mathcal E_1$ is $\kappa$-filtered and has no upper bound in the collage, hence no colimit. This also shows why nonemptiness of the open fibre does not rescue the statement.

The correction is to add that $\mathcal E_1$ satisfies $\mathrm L_\pi$ on the right-hand side. Sufficiency is precisely the two-case argument above. For necessity, assume the left-hand side. A nonempty filtered diagram in $\mathcal E_1$ cannot have a colimit in $\mathcal E_0$, because there is no cocone from it to such an object. Its colimit in $\mathcal E$ therefore lies in $\mathcal E_1$ and is universal there. The preservation condition on $H$ follows by testing a colimit from $\mathcal E_0$ against an object of $\mathcal E_1$. The assertions about fibrewise small objects then follow by restricting the relevant representables.

This omission does not invalidate Proposition 4.14.24: there the cardinal was explicitly chosen useful for both fibres. It matters when the corollary is reused as a separate fixed-cardinal criterion, which is exactly what its formulation invites the reader to do.
\end{dcomment}


%=====================================================================
\printedpage{160}
% head: A. Grothendieck : Dérivateurs, Ch. XVIII, Édition des Universités de Montpellier II et Paris VII

\item[\textit{3°)}] \textit{Si $\pi_0$ et $\pi$ sont comme dans 2°), et si de plus $\pi$ est utile pour $\cE_1$, alors $\pi$ est utile pour $\cE$, donc (en vertu de 2°)}
\[
\cE \;\xleftarrow{\ \sim\ }\; \Ind_\pi(\cE(\pi))\ .
\]
\end{enumerate}
Nous pouvons maintenant énoncer la généralisation cherchée de 4.14.17 :

\msp{page 422}
\textbf{Théorème 4.14.26} (généralisation de 4.14.13). \textit{Soit $p : \cE \longrightarrow \cC$ un foncteur, avec $\cC$ petite. On suppose que $p$ satisfait à l'hypothèse H) de la page 409, qui équivaut (moyennant 4.14.24) à la suivante :}
\begin{enumerate}
\item[\textit{H$'$)}] \textit{$\forall\ s \in \Ob\cC$, la catégorie fibre $\cE_s$ est accessible, et pour tout $\lambda : s \longrightarrow t$ dans $\Fl\cC$, le foncteur}
\[
\left\{
\begin{array}{lcl}
H_\lambda : \cE_s^\circ \times \cE_t & \longrightarrow & \Ens\\[1ex]
\qquad (x_0,x_1) & \longmapsto & \Hom_\lambda(x_0,x_1)
\end{array}
\right.
\]
\textit{satisfait les deux conditions suivantes :}
\begin{enumerate}
\item[\textit{a)}] \textit{$\forall\ x_0 \in \Ob\cE_s^\circ$, $x_1 \longmapsto H_\lambda(x_0,x_1) : \cE_t \longrightarrow \Ens$ est accessible.}
\item[\textit{b)}] \textit{Il existe un cardinal infini $\pi_\lambda$, tel que pour tout $x_1 \in \Ob\cE_t$, le foncteur $x_0 \longmapsto H_\lambda(x_0,x_1) : \cE_s \longrightarrow (\Ens)^\circ$ soit $\pi_\lambda$-accessible.}
\end{enumerate}
\end{enumerate}
\textit{Ceci posé, la catégorie $\underline{\Gamma}(\cE/\cC)$ des sections de $\cE$ sur $\cC$ est accessible, et les foncteurs $\varepsilon_s : \underline{\Gamma}(\cE/\cC) \longrightarrow \cE_s$ sont accessibles. On peut préciser de façon similaire à 4.14.18 (énoncé en forme laissé au lecteur).}

\msp{page 423}
\textsc{Démonstration} de 4.14.26. Elle se fait exactement comme celle de 4.14.17 (p. 400 ff.), mais est conceptuellement plus simple. On suppose $\cC$ petite, et on introduit les catégories
\[
\Pi \;=\; \prod_{s\,\in\,\Ob\cC} \cE_s\ , \qquad \Pi' \;=\; \prod_{\lambda\,\in\,\Fl\cC} \cE_\lambda\ ,
\]
où
\[
\cE_\lambda \;=\; \cE \times_{\cC} (\boldsymbol{\Delta}_1, i_\lambda)\ ,
\]
et les deux foncteurs canoniques
\[
\Pi \;\underset{G}{\overset{F}{\rightrightarrows}}\; \Pi'
\]
de composantes
\[
\begin{aligned}
F_\lambda \;&=\; i_\lambda \circ \pr_{s(\lambda)}\\
G_\lambda \;&=\; j_\lambda \circ \pr_{b(\lambda)}\ ,
\end{aligned}
\]
où pour $\lambda : s \longrightarrow t$, on pose
\[
s(\lambda) \;=\; s\ , \qquad b(\lambda) \;=\; t\ ,
\]
et $i_\lambda$, $j_\lambda$ sont les inclusions canoniques
% In the print, the vertical arrow p_lambda and Delta_1 are enclosed together in a pair of large parentheses.
\[
\begin{tikzcd}
\cE_s \arrow[r,"i_\lambda"] \arrow[d,dash,"\wr"] & \cE_\lambda \arrow[dd,"p_\lambda"] & \cE_t\ . \arrow[l,"j_\lambda"'] \arrow[d,dash,"\wr"] \\
(\cE_\lambda)_0 & & (\cE_\lambda)_1 \\
& \boldsymbol{\Delta}_1 &
\end{tikzcd}
\]
% foot: Version 17 décembre 2025

%=====================================================================
\printedpage{161}
% head: A. Grothendieck : Dérivateurs, Ch. XVIII, Édition des Universités de Montpellier II et Paris VII

\msp{page 424}
On désigne par $\cM'$ la catégorie
\[
\cM' \;=\;
\begin{array}{l}
\text{catégorie des couples } (X,u)\text{, avec}\\
X \in \Ob\Pi,\ u : F(x) \longrightarrow G(x),
\end{array}
\] % [sic] F(x), G(x) -> F(X), G(X), the object is named X
qui est constructible par 4.14.22.1 est accessible (car $\Pi$ et $\pi'$ et les foncteurs $F$ et $G$ le sont, par l'hypothèse H) faite sur $p$). On peut l'interpréter comme la catégorie des sections % [sic] \pi' -> \Pi', the category introduced is \Pi'
de $\mathrm{Diag}(\cE)$ sur $\mathrm{Diag}(\cC)$. On a un foncteur strictement pleinement fidèle
\[
\cM \;\defeq\; \underline{\Gamma}(\cE/\cC) \;\hookrightarrow\; \cM'\ ,
\]
qui permet d'interpréter $\cM$ comme une sous-catégorie strictement pleine de $\cM'$. On voit alors comme dans le cas de 4.14.17 que cette sous-catégorie est accessible dans $\cM'$ -- c'est essentiellement sorital. Il s'ensuit que $\cM$ est accessible et que les foncteurs $\cM \longrightarrow \cE_s$ le sont aussi, comme composés de
\[
\cM \;\overset{\text{\upshape acc.}}{\hookrightarrow}\; \cM' \;\overset{\text{\upshape acc.}}{\hookrightarrow}\; \Pi \;\xrightarrow{\ \ \pr_s\ \ \text{\upshape acc.}\ \ }\; \cE_s\ ,
\]

\dgrblock{dgrb1107}{Sections without a fibration structure}
\begin{dcomment}{exposition}{p423}{Sections without a fibration structure}
Theorem 4.14.26 identifies how much of the section argument actually used cartesian or cocartesian transport. A single base arrow $\lambda:s\to t$ determines a collage $\mathcal E_\lambda$, whether or not a preferred lift exists. The hypotheses say that this collage and its endpoint inclusions are accessible. That is enough to choose the arrow data of a section by an inserter.

For a family $X=(X_s)$, an arrow
\[
 i_\lambda X_s\longrightarrow j_\lambda X_t
 \quad\hbox{in }\mathcal E_\lambda
\]
is exactly an arrow $X_s\to X_t$ of $\mathcal E$ lying over $\lambda$. Thus the new inserter is again the category of graph sections. The remaining equations are literal identities and compositions in $\mathcal E$. No comparison isomorphisms from a cleavage occur.

One accessibility check is compressed into the word ``sorital'' on page 424. Choose a common regular cardinal at which the fibres, the interval categories, their endpoint inclusions, and the graph-section construction admit and preserve filtered colimits. Given a filtered diagram of graph sections, form all its object values in the fibres. The structural arrow over $\lambda$ is the unique map induced by the corresponding colimit cocones in $\mathcal E_\lambda$. For composable $\lambda$ and $\mu$, composition of the induced arrows agrees with the arrow induced by the composites: precomposing with every object of the source colimit gives the same arrow, and the universal property gives uniqueness. Consequently the two functors recording $a_{\mu\lambda}$ and $a_\mu a_\lambda$ preserve these filtered colimits. The same argument handles identity arrows.

To impose their equality, regard the pair as a diagram of parallel arrows in $\mathcal E_{\mu\lambda}$, with fixed endpoint degrees. The category of equal pairs is accessible and accessibly embedded, just as on pages 405--406. There are only a set of objects and composable pairs in the small base. Pulling back these equation loci and intersecting them gives an accessible full subcategory of graph sections. Its evaluation functors are accessible by composition with the inserter projection and the product projections. The hooked arrow $\mathcal M'\hookrightarrow\Pi$ in the last display should therefore be an ordinary forgetful arrow; the argument only uses its accessibility.

This proof is local along arrows in a precise sense. It tests the categories of arrows over each edge, while functoriality supplies the relations among edges. It does not assert that the total category $\mathcal E$ is accessible for an arbitrary small base. The next section imposes an additional condition on idempotents in the base to address that separate conclusion.
\end{dcomment}


\hfill q.e.d.

\msp{page 425}

\subsection*{4.15\quad Catégorie accessible sur une autre}

\textbf{Proposition 4.15.1.} \textit{Soit $f : \cE \longrightarrow \cC$ un foncteur avec $\cC$ petite. On suppose de plus satisfaite la condition H (p. 409), qu'on explicite ici sous la forme équivalente H$'$ :}
\begin{enumerate}
\item[\textit{H 1)}] \textit{Pour tout $s \in \Ob\cC$, $\cE_s$ est accessible.}
\item[\textit{H 2)}] \textit{Pour toute $\lambda \in \Fl\cC$, $\lambda : s \longrightarrow t$, existe un cardinal infini $\pi_\lambda$ tel que pour tout $x_1 \in \Ob\cE_t$, le foncteur}
\[
x_0 \;\longmapsto\; \Hom_\lambda(x_0,x_1) \;:\; \cE_s \;\longrightarrow\; (\Ens)^\circ
\]
\textit{est $\pi_\lambda$-accessible ($\cE_s$ satisfaisant de plus à $L_{\pi_\lambda}$).}
\item[\textit{H 3)}] \textit{Pour toute $\lambda \in \Fl\cC$, $\lambda : s \longrightarrow t$, et pour tout $x_0 \in \Ob\cE_s$, le foncteur}
\[
x_1 \;\longmapsto\; \Hom_\lambda(x_0,x_1) \;:\; \cE_t \;\longrightarrow\; (\Ens)
\]
\textit{est accessible.}
\end{enumerate}
\textit{Alors les foncteurs d'inclusion $i_s : \cE_s \longrightarrow \cE$ sont accessibles. De plus, si $\cC$ ``est sans projecteurs'' (i.e. si tout endomorphisme $p$ d'un objet $s$ de $\cC$ tel que $p^2 = p$, est l'identité), alors $\cE$ est accessible, et $f : \cE \longrightarrow \cC$ est accessible.}

Cela résulte du résultat plus précis :
% foot: Version 17 décembre 2025

\dgrblock{dgrb1108}{A presentability test and a necessary rank repair}
\begin{dcomment}{development}{p425:Cela résulte du résultat plus précis}{A presentability test and a necessary rank repair}
Grothendieck's next aim is to reconstruct accessibility of the total category from its fibres. The explicit bounds in Corollary 4.15.2 are meant to make this reconstruction reusable. They omit one condition needed by the subsequent calculation: the initial presentation rank $\pi_0$ must already control the contravariant functors in H2. Bounding those ranks only by $\pi_1$ does not justify applying them to a $\pi_0$-filtered presentation. We make the issue concrete before giving a repair.

\textbf{Fibrewise presentability test over a finite poset.}\ours\ 

Let $p:\mathcal A\to P$ have finite poset base. Fix a regular infinite $\kappa$. Suppose every fibre $\mathcal A_s$ has $\kappa$-filtered colimits and, for $s\leq t$ and $y\in\mathcal A_t$, the functor
\[
 \Hom_{s\leq t}(-,y):\mathcal A_s^{\mathrm{op}}\to\mathrm{Set}
\]
sends these colimits to limits. Then the total category has $\kappa$-filtered colimits, computed in a cofinal fibre, and each fibre inclusion preserves them. An object $x\in\mathcal A_s$ is $\kappa$-presentable in $\mathcal A$ if and only if every functor $\Hom_{s\leq t}(x,-)$, for $t\geq s$, preserves $\kappa$-filtered colimits. The case $t=s$ includes presentability in the fibre.

\emph{Proof.} The image of a filtered diagram in the finite poset is a finite directed subset and therefore has a greatest element $t$. Choose an index mapping to $t$. The overcategory of that index is cofinal, and all its images equal $t$; the diagram has consequently reduced to one fibre. To show that a fibrewise colimit is universal in the total category, test maps to an object in every fibre. If there is no base arrow the relevant sets are empty; otherwise the assumed contravariant comparison is the required bijection. This proves colimit existence and preservation. If $x$ is presentable in the total category, restrict its representable to a diagram in $\mathcal A_t$, using the just-proved preservation. Conversely reduce any target diagram to its cofinal fibre $\mathcal A_t$; the asserted preservation of the corresponding cross-Hom gives precisely the presentability comparison for $x$. This proves the test.

If the fibres are $\kappa$-accessible and every fibrewise $\kappa$-presentable object passes this test, the union of their small generating subcategories consists of total $\kappa$-presentables and generates $\mathcal A$. Hence $\mathcal A$ is $\kappa$-accessible. This last statement follows directly from the definition and does not require additional colimits.

\textbf{Counterexample to the printed bound in 4.15.2(c).}\ours\ 

Use Grothendieck's convention: $\pi$-presentable sets have cardinal at most $\pi$, corresponding to modern $\pi^+$-presentability. Put
\[
 |K|=\aleph_1,\qquad
 \pi_0=\aleph_0,\qquad
 \pi_1=2^{\aleph_1},\qquad
 \pi=\beth_{\omega_1}.
\]
Take the opfibration over $[1]$ classified by
\[
 \varphi:\mathrm{Set}\to\mathrm{Set},\qquad X\longmapsto X^K.
\]
Its cross-Hom is $H(X,Y)=\mathrm{Set}(X^K,Y)$. The functor $\varphi$ preserves $\aleph_2$-filtered colimits: a $K$-indexed family and all equalities among its representatives are witnessed at one stage because $|K|<\aleph_2$. Thus H2 holds with $\pi_\lambda=\aleph_1$. Both fibres are useful at $\pi_0$, and their reduced $\pi_0$-presentable subcategories have size at most $2^{\aleph_0}$. For countable $X$, the set $X^K$ has size at most $2^{\aleph_1}=\pi_1$, so $H(X,-)$ is $\pi_1$-accessible. All three printed bounds of part (c) are satisfied.

The cardinal $\pi$ is strong limit of cofinality $\aleph_1$. Every countable sequence in $\pi$ is bounded, whence $\pi^{\aleph_0}=\pi$: for each bound $\nu<\pi$, one has $\nu^{\aleph_0}\leq2^{\nu\cdot\aleph_0}<\pi$, and there are $\pi$ possible bounds. In contrast, K\"onig's inequality gives $\pi^{\aleph_1}>\pi$. Choose a lower-fibre set $X$ of size $\pi$. It is $\pi$-presentable in its fibre, but $X^K$ has size greater than $\pi$. The $\pi^+$-filtered diagram of all subsets of $X^K$ of size at most $\pi$, placed in the upper fibre, has colimit $X^K$. The identity of $X^K$ does not factor through any of those subsets. Thus $H(X,-)$ does not preserve this colimit, and the fibrewise test proves that $X$ is not $\pi$-presentable in the total category. The asserted inclusion $\mathcal E(\pi)\subseteq\mathcal E_\pi$ fails at the printed bounds.

The qualitative proposition is repaired by choosing $\pi_0$ useful for all fibres \emph{and} at least every $\pi_\lambda$, then making the subsequent choices of $\pi_1$ and $\pi$ as before. Now a $\pi_0$-filtered presentation may legitimately be passed through each contravariant $H_\lambda(-,y)$; the small-limit calculation on page 419 applies. This is also the missing justification in the continuation of the proof on manuscript page 430. We found no discussion of this precise printed rank omission in the accessible-category references checked here. The test is a direct reorganization of the manuscript's proof, and the counterexample concerns its formulation; neither is asserted as a priority claim.
\end{dcomment}


%=====================================================================
\printedpage{162}
% head: A. Grothendieck : Dérivateurs, Ch. XVIII, Édition des Universités de Montpellier II et Paris VII

\msp{page 426}
\textbf{Corollaire 4.15.2} (\footnote{Pour a) et b) on n'utilise que H 1 et H 2, à l'exclusion de H 3.}). \textit{Les hypothèses sur $f : \cE \longrightarrow \cC$ sont celles de 4.15.1 ($\cC$ petite, et $f$ satisfait $H' =$ H 1) 2) 3)).}
\begin{enumerate}
\item[\textit{a)}] \textit{Soit $s \in \Ob\cC$ et soit $\pi_1$ un cardinal infini tel que pour toute $\lambda : s \longrightarrow t$ (flèche de $\cC$ de source $t$), on ait $\pi_1 \geqslant \pi_\lambda$ (où $\pi_\lambda$ est défini dans l'énoncé de H 2). Alors le foncteur $i_s : \cE_s \longrightarrow \cE$ est $\pi_1$-accessible.} % [sic] source t -> source s, lambda : s -> t has source s (cf. p. 163 "de source s")
\item[\textit{b)}] \textit{On suppose $\cC$ sans projecteurs. Soit $\pi_1$ un cardinal $\geqslant \sup(\pi_\lambda, \card\cC)$. Alors $\cE$ satisfait $L_{\pi_1}$, et $f : \cE \longrightarrow \cC$ est $\pi$-accessible.}
\item[\textit{c)}] \textit{Supposons encore que $\cC$ soit sans projecteurs (p.ex. une catégorie ``rigide'', p.ex. une catégorie préordonnée). Soient $\pi_0$, $\pi_1$ des cardinaux tels que :}
\begin{enumerate}
\item[\textcircled{\footnotesize $1^o$}] $\forall\ s \in \Ob\cC$, $\pi_0$ est utile pour $\cE_s$ \textit{(un tel $\pi_0$ existe grâce à H 1 et au fait que $\cC$, donc $\Ob\cC$ petits).}
\item[\textcircled{\footnotesize $2^o$}] $\pi_1 \geqslant \sup(\pi_0, \card\cC, \sup_{\lambda\in\Fl\cC}\pi_\lambda, \sup_{s\in\Ob\cC} \cardred \cE_{s\,\pi_0})$.
\item[\textcircled{\footnotesize $3^o$}] \textit{$\forall\lambda \in \Fl\cC$, et $x_0 \in \Ob\cE_{s,\pi_0}$, le foncteur $H_\lambda(x_0)$ de H 3, $x_1 \longmapsto \Hom_\lambda(x_0,x_1) : \cE_t \longrightarrow \Ens$ est $\pi_1$-accessible.}
\end{enumerate}
\textit{Soit $\pi$ un cardinal tel que $\pi^{\pi_0} = \pi \geqslant \pi_1$. Alors par b), $\cE$ satisfait $L_{\pi_1}$, donc $L_\pi$, et $f : \cE \longrightarrow \cC$ est $\pi_1$-accessible. De plus,}
\[
i_s(\cE_{s\,\pi}) \;\subset\; \cE_\pi\ ,
\]
\textit{ou encore}
\[
\cE(\pi) \;\subset\; \cE_\pi\ , \tag*{$(*)$}
\]
\textit{où $\cE(\pi)$ est la sous-catégorie pleine de $\cE$ telle que}
\[
\Ob\cE(\pi) \;=\; \bigcup_{s\in\Ob\cC} \Ob\,(\cE_{s\,\pi})\ .
\]
\textit{Comme $\cE_\pi$ satisfait $L_{\pi_1}$ par b), donc a fortiori $L_\pi$, on a donc par $(*)$ un foncteur} pleinement fidèle
\[
\Ind_\pi(\cE_{s\,\pi}) \;\xrightarrow{\ \sim\ }\; \cE\ ,
\]
\textit{et ce foncteur est une équivalence de catégories (donc $\cE$ est accessible).}
\item[\textit{d)}] \textit{En fait, sous les hypothèses de c) sur $\cC$, $\pi_0$, $\pi_1$, $\pi$, on a $\cE(\pi) = \cE_\pi$ (non seulement $\cE(\pi) \subset \cE_\pi$).}
\end{enumerate}

\msp{page 427}
\textsc{Démonstration de 4.15.1 et 4.15.2.} Il est clair que 4.15.1 est contenu dans 4.15.2 a) et c). Nous allons donc prouver 4.15.2.

a) Soit $(x_i)_{i\in I}$ un système inductif dans $\cE_s$, avec $I$ grand devant $\pi_0$, et $x$ sa limite inductive dans $\cE_s$ (qui existe, car $\cE_s$ satisfait à $\mathrm{L}_{\pi_\lambda}$, donc à $\mathrm{L}_{\pi_0}$). Il faut prouver que c'est une limite dans $\cE$, donc que pour tout $y \in \Ob\cE$, la flèche canonique
\[
\Hom_{\cE}(x,y) \;\longrightarrow\; \varprojlim_{I} \Hom_{\cE}(x_i,y) \tag*{$(*)$}
\]
% foot: Version 17 décembre 2025

%=====================================================================
\printedpage{163}
% head: A. Grothendieck : Dérivateurs, Ch. XVIII, Édition des Universités de Montpellier II et Paris VII
est un isomorphisme. Or $\Hom_{\cE}(x,y)$ et les $\Hom_{\cE}(x_i,y)$ sont des objets de $\Ens/\Hom_{\cC}(s,t)$, où $s = f(x)$, $t = f(y)$, et les flèches de transition sont des flèches dans $\Ens/\Hom_{\cC}(s,t)$. Cela montre que la bijectivité de $(*)$ revient à celle de
\[
\Hom_\lambda(x,y) \;\longrightarrow\; \varprojlim_{I} \Hom_\lambda(x_i,y) \tag*{$(*)$}
\]
pour tout $\lambda \in \Hom_{\cC}(x,y)$ \texttt{[plutôt} $\Hom_{\cC}(f(x),f(y))$\texttt{]}. Pour $\lambda$ fixé, et $y$ fixé dans $\Ob\cE_t$, la bijectivité dans $(*)$, pour tout système inductif grand devant $\pi_1$ dans $\cE_s$, de limite $x$, signifie que le foncteur
\[
x \;\longmapsto\; \Hom_\lambda(x,y) \;:\; \cE_s \;\longrightarrow\; (\Ens)^\circ
\]
est $\pi_0$-accessible. Or c'est bien le cas, par hypothèse sur $\pi_1$, quel que soit la flèche $\lambda : s \longrightarrow t$ de source $s$, et quel que soit $y \in \Ob\cE_t$.

\msp{page 428}
Soit $(x_i)_{i\in I}$ un système inductif grand devant $\pi_1$, et $(s_i)_{i\in I}$ l'image de ce système dans $\cC$. Nous prouverons plus bas le

% In the left margin, beside the statement of Lemme 4.15.3 and beside the following paragraph "Mais attention ...", the edition prints a wavy marginal sign (one for each).
\textbf{Lemme 4.15.3} (\footnote{Cf. démonstration page 452.},\footnote{sous réserve de la validité de 4.8.13 (Mais c'est à présent établi, sous une forme suffisante, dans 4.15.10).}). \textit{Soit $\cC$ une catégorie, $\pi$ un cardinal, tels que} $\card\cC \leqslant \pi$\textit{, et que $\cC$ soit sans projecteurs. Alors pour tout système inductif $(s_i)_{i\in I}$ grand devant $\pi$ dans $\cC$, existe une partie cofinale $J$ dans $I$ telle que le système induit $(s_i)_{i\in J}$ soit} constant, \textit{i.e. tel que les $s_i$ soient égaux à un $s \in \cC$, et que les morphismes de transition soient tous égaux à} $\id_s$.

Mais attention, on s'appuie pour prouver 4.15.3 sur le théorème 4.8.13 (page 212-216), que je n'ai toujours pas prouvé. Dans le cas où $\cC$ est rigide (i.e. $\End_{\cC}(s) = \{\id_s\}\ \forall\ s \in \Ob\cC$), la démonstration est cependant immédiate, cf. plus bas.

\msp{page 429}
Pour prouver que $x = \varinjlim\limits_{I} x_i$ existe dans $\cE$, et commute à $f$, on peut remplacer $I$ par une sous-catégorie cofinale. Cela nous ramène, grâce au lemme, au cas où $(s_i)_{i\in I}$ est constant, donc au cas où le système inductif provient d'un système inductif dans une catégorie fibre $\cE_s$. Mais par a), on sait que sa limite $x$ dans $\cE_s$ est une limite dans $\cE$, d'où l'existence de cette dernière, et la commutation à $f$, puisque $\varinjlim s_i = s$ ( $\varinjlim$ d'un système inductif filtrant constant).

\dgrblock{dgrb1109}{What an idempotent in the base can obstruct}
\begin{dcomment}{examples}{p428}{What an idempotent in the base can obstruct}
The restriction on the base is meant to force a sufficiently filtered diagram to become constant after a cofinal change of indices. If the base permits a nonidentity idempotent, such a diagram can instead settle on repeated applications of that idempotent. The simplest example shows why the local hypotheses alone do not guarantee accessibility of the total category.

Let $B$ have one object $*$ and endomorphism monoid $\{1,e\}$ with $e^2=e\ne1$. Take $p=\id_B$. Its strict fibre is terminal. For either base arrow $1$ or $e$, the pullback $B\times_B[1]$ is $[1]$, with its identity projection, so hypothesis H holds. Nevertheless $B$ is not accessible. For any infinite regular $\kappa$, take the diagram indexed by the ordinal $\kappa$, constant on $*$, with every strictly increasing transition equal to $e$. This is a well-defined $\kappa$-filtered diagram because $e^2=e$.

A cocone to $*$ is a family $(a_i)$ with $a_j e=a_i$ whenever $i<j$. Since both $1e$ and $ee$ equal $e$, and every index has a larger index, every $a_i$ must be $e$. There is exactly one cocone. Were it colimiting, the map from $\Hom_B(*,*)$ to this singleton cocone set would be bijective. It sends both $1$ and $e$ to the same cocone and is not injective. There is no other possible colimit object. Hence $B$ has no such colimit at any regular rank.

This is the elementary splitting obstruction: a colimit of the repeated-idempotent system would supply the missing image of $e$. The example explains the hypothesis without claiming it is necessary in every situation. A base with split nonidentity idempotents may be accessible, and additional total-category hypotheses can also supply the missing splittings.

For the general constant-subdiagram lemma the manuscript itself records a deferred dependency. We checked the continuation in the editors' complete Chapter XVIII: page 452 still uses 4.8.13, and Corollaries 4.15.10--11 on pages 462--463 finally justify the argument with the strict bound $\pi>\operatorname{card}\mathcal C$. Thus the qualitative conclusion of 4.15.1 is secured by increasing the chosen cardinal. The boundary equality $\pi=\operatorname{card}\mathcal C$ is not established by that continuation, and we do not report it as proved here. Separately, the two displayed slips on page 426 must be read as colimit existence in $\mathcal E$, and the equivalence $\operatorname{Ind}_\pi(\mathcal E(\pi))\simeq\mathcal E$ using all fibres. These pages end during the next presentability calculation.
\end{dcomment}


c) Il faut prouver que si $x \in \Ob\cE_s$ est $\pi$-accessible dans $\cE_s$, il l'est dans $\cE$, i.e. que pour tout système inductif $(y_i)_{i\in I}$ grande devant $\pi$, la flèche canonique
\[
\varinjlim_{i} \Hom(x,y_i) \;\longrightarrow\; \Hom(x,y) \tag*{$(**)$}
\]
est bijective. Or comme par 2°) $\pi \geqslant \pi_1 \geqslant \sup(\card S, \sup_\lambda \pi_\lambda)$, il résulte encore du lemme % [sic] card S -> card \mathcal{C}, no S is defined; the condition 2° bounds card \mathcal{C}
que, quitte à remplacer $I$ par une sous-catégorie cofinale, on peut supposer que le système inductif $(y_i)_{i\in I}$ provient d'un système dans une catégorie fibre $\cE_t$, et que $y$ est sa $\varinjlim$ dans $\cE_t$. Mais alors les ensembles $\Hom(x,y)$ et $\Hom(x,y_i)$, et les morphismes de transition entre celles-ci, sont encore dans $\Ens/\Hom(s,t)$, et on est ramené
% foot: Version 17 décembre 2025
\begin{otherlanguage}{english}
\dgroverview{XVIII-d}{Chapter XVIII \textperiodcentered\ \textup{XVIII-d} \textperiodcentered\ Overview}
\begin{otherlanguage*}{english}
\noindent\textbf{Accessible conditions, categories of structures, and gluing.}

These pages turn accessibility into a working calculus. Their central question is how a condition on the objects of a large category can be recognized, transported, and imposed without losing a small description of the resulting category. The first part completes a criterion for accessible full subcategories. The second applies it to functors preserving specified cones and cocones, to categories of sections, and finally to categories lying over a small base. The progression is deliberate: once membership in a subcategory can be read on a suitable presentation, many apparently new constructions become instances of the same argument.

One convention must be fixed before comparing the manuscript with current accounts. Grothendieck calls an indexing category \emph{grande devant $\pi$} when it admits cocones on diagrams having at most $\pi$ arrows. Thus the associated regular cardinal in present terminology is $\lambda=\pi^+$, and his $\pi$-accessible objects are $\lambda$-presentable objects. A useful $\pi$ provides an essentially small full subcategory $\mathcal M_\pi$ generating $\mathcal M$ under these filtered colimits. The stronger phrase accessible \emph{in} an ambient category keeps track of the inclusion and of ambient presentability, not just an abstract equivalence class of accessible categories.

The opening argument is a transfinite alternation between two presentations. A final functor $\varphi:I\to J$ and a full replete final subcategory $J'\subset J$ do not in general have a final strict inverse image: the even and odd natural numbers in their common ordered set already show the problem. Grothendieck adds an ordinal-colimit hypothesis. Starting from a chosen object of $I$, he constructs alternating maps
\[
 \varphi(i_\beta)\longrightarrow j_\beta
 \longrightarrow\varphi(i_{\beta+1}),
 \qquad j_\beta\in J',
\]
whose composites are transition maps. At the final limit, the two chains have isomorphic colimits. One colimit comes from $I$, and the other lies in $J'$; repleteness then gives an object of the inverse image beyond the prescribed starting object. Lemma 4.13.17 is therefore a genuine cofinality statement with an additional closure hypothesis. The commentary supplies its final comparison explicitly and shows that the successor-stage equalizations in the manuscript are redundant.

This lemma drives Theorem 4.13.22. Suppose $\underline S$ is an accessible full subcategory of $\mathcal M$ closed under the indicated colimit shapes. In an adapted presentation $x=\varinjlim_i x_i$, membership $x\in\underline S$ is equivalent to finality of the full subcategory of indices for which $x_i\in\underline S$. Adaptedness says that the indexing category contains the specified colimits and that the presenting diagram preserves them. It is what prevents a presentation from hiding the condition at every finite stage. The worked example of divisible abelian groups exhibits the criterion through countable divisible subgroups; the familiar presentation of $\mathbb Q$ by cyclic groups explains why an arbitrary filtered presentation is insufficient.

The subsequent calculus treats enlargement of the accessibility rank, small unions and intersections, and inverse images under accessible functors. It also exposes a boundary of the method. The essential full image of an accessible functor need not be accessibly embedded, even after closing under retracts. We work out the split-monomorphism example recorded in \citet{Rosicky2009Combinatorial,BrookeTaylorRosicky2017}: forgetting chosen retractions allows diagrams whose colimit no longer splits. This failure is distinct from the controlled envelope construction. Starting with a small family of objects and closing it under a fixed filteredness threshold and a small set of additional colimit shapes does yield an accessible, accessibly embedded subcategory. The commentary proves this by choosing a sufficiently closed cardinal, making bounded filtered envelopes of index categories, and identifying the result with a fully faithful Ind-image.

The historical framework is already visible in \citet[Expos\'e I, 8.7]{SGA4I1972}, where Ind-objects and their universal extension property are developed. The systematic accessible-category theory of \citet{MakkaiPare1989} was published in 1989, before these pages were written. The relation is therefore a comparison of formulations, not a claim that this manuscript first announced a theory later discovered by those authors. The modern presentation separates products, inserters, equifiers, and pseudopullbacks; Grothendieck often constructs the same operations through accessible loci in categories of diagrams. \citet{Positselski2024Limits} makes especially clear how additional chain-colimit hypotheses permit control of a specified rank.

Section 4.14 begins with marked cones and cocones in a small category. A functor is required to take the marked data to genuine limits or colimits. When the target has the relevant universal constructions, each condition is the inverse image of an isomorphism locus under an accessible comparison functor. Their small intersection is again accessible. This is the language of sketches and categories of models. Restriction to presentable objects then allows certain large-source functor categories to be treated by the same method.

Here the hypotheses require particular care. The proof of the fixed-threshold functor-category assertion invokes target colimits not supplied by accessibility alone. We distinguish the theorem established with the additional target-colimit hypothesis from the more general assertion the manuscript leaves to its later question about colimit cones. We also distinguish the two adjoint theorems: cocontinuous functors from a locally presentable category have right adjoints, whereas a limit-preserving covariant functor need not have a left adjoint unless a solution-set or accessibility condition is added. The stronger claim at pages 379--380 fails already for the category of groups, a warning that Grothendieck himself recalled near the beginning of this chapter. The counterexample and the appropriate repair are explained with the primary references.

The question about epimorphisms has a different status. Grothendieck observes that epimorphisms are stable under the relevant filtered colimits and asks whether their full subcategory in the arrow category is always accessible. Closure alone does not answer that question. The commentary proves that accessibility and accessible embedding of this epimorphism locus would force the underlying accessible category to be co-wellpowered. Published results show that the universal co-wellpoweredness assertion already has a large-cardinal lower bound, while stronger large-cardinal assumptions give positive results \citep{BagariaCasacubertaMathiasRosicky2015}. This identifies a real set-theoretic obstruction. It does not settle the exact strength of the epimorphism-locus assertion. The recent study of pure epimorphisms by \citet{Positselski2026Pure} concerns an additional purity condition; its publisher abstract was checked as a status lead, and no result about the unrestricted epimorphism class is inferred from it.

Categories of sections provide the main constructive application. For a fibration, a prospective section consists of a family of fibre objects and arrows $X_s\to\lambda^*X_t$. Those data form an inserter. Identity and composition impose a small family of equations, and cartesian or cocartesian requirements impose isomorphism conditions. The proof keeps track of a nonsplit cleavage through its canonical coherence isomorphisms. It is thereby a genuinely categorical construction, not a pointwise prescription with the functoriality omitted. The higher-category analogue is the accessible section theorem of \citet[Proposition 5.4.7.11]{Lurie2009HTT}.

The qualitative section theorem is robust; several printed rank assertions are not. For unrestricted sections of an opfibration, pointwise presentability need not imply presentability of the section unless transport has an appropriate presentability bound. Conversely, the category of sections can be very small while a forced component is large. The commentary gives explicit examples for both directions. Cocartesian sections admit a separate argument, because their structure isomorphisms identify the transported source with the target component. Keeping that distinction avoids discarding a correct corollary along with an incorrect proof reference.

The final pages replace fibrations by intervalwise accessible data. An arbitrary functor $p:\mathcal E\to\mathcal C$ determines, over each base arrow, a category obtained by gluing its two endpoint fibres. Its cross-Hom bifunctor
\[
 H:\mathcal E_0^{\mathrm{op}}\times\mathcal E_1\to\mathrm{Set}
\]
controls the gluing. Contravariant preservation makes fibrewise colimits universal in the total category; covariant preservation makes source objects presentable at a suitable rank. This proves the accessible-section theorem without requiring chosen lifts. An additional hypothesis on idempotents in the base addresses accessibility of the total category itself.

Two corrections matter for reusing this last argument. Corollary 4.14.25 omits the colimit hypothesis on the closed fibre from one side of its equivalence; an ordinal example disproves the printed version. Corollary 4.15.2 also fails to put the contravariant preservation ranks below the initial presentation rank. The power-functor collage $X\mapsto X^{\aleph_1}$ gives a counterexample satisfying the printed cardinal bounds. Both qualitative constructions survive a common-rank choice with the missing requirements included. These pages end during the proof; the complete chapter later justifies the constant-subdiagram argument with the strict inequality $\pi>|\mathcal C|$. That later repair is reported at its exact manuscript locations.

The audit records these distinctions assertion by assertion: what the source proves, sketches, states, or asks; what is merely a slip in the edition; and what needs a mathematical repair. A few exact formulations remain uncertified, including the unrestricted fixed-rank third variant of 4.13.21 and the broad weakening suggested in footnote 172. These are identified as limits of the present verification rather than declared counterexamples or new open problems. Proofs marked \emph{[ours]} are arguments supplied by the commentators, with the scope of the literature check stated and no claim of priority.

\end{otherlanguage*}
\end{otherlanguage}
\dgrchunkend
